\section{Recovery grades and deadlines}\label{sec:grades}

This section defines the four SG grades and proves the delivery and
compatibility facts used by the stable-root rule.
Figure~\ref{fig:recovery-round} in Section~\ref{sec:rounds} lists the
complete recovery round in one place; the sections through
Section~\ref{sec:rounds} define each of its steps.  Six recurring
concepts appear in this section before their full rules do; the
following interface fixes the meaning of each so every use in this
section is defined.

\begin{definition}[Stable root and action root, interface]
\label{def:root-interface}
Each recovery round fixes two per-round roots for every validator, both
derived by the rules of Definition~\ref{def:stable-root} in
Section~\ref{sec:recovery-proposals}.  The six interfaced concepts are:
\begin{enumerate}
 \item the \emph{stable root}: the root the validator fixes at its
   vote time --- a grade-backed root, or its own finality-backed
   Simplex selection --- and uses as the walk root of its ordinary
   Goldfish votes for the round, subject only to the re-derivation rule
   of Definition~\ref{def:rederivation};
 \item the \emph{action root}: the root the round's action fixes ---
   the round's vote-time stable root, possibly re-derived, when the
   action-state admission check accepts it, and no root otherwise.  The
   SG head the validator signs is derived from it: the deepest official
   confirmation when one exists, and otherwise the deepest admissible
   block between the action-state Simplex root and the action root, per
   the head-production rule of the confirmation section.  A validator
   whose vote-time stable root fails the action-state admission of
   those rules emits no SG head and no current-height pair; the guards
   themselves appear with Definition~\ref{def:stable-root};
 \item an \emph{active} grade: a grade counts only against the
   validator's processed finality evidence --- the filter is
   Definition~\ref{def:active-grade} in this section;
 \item the \emph{available confirmation} evaluated at each action:
   the time-shifted-quorum rule of
   Section~\ref{sec:available-confirmation}, which replaces both of
   Goldfish's confirmation rules --- \(\kappa\)-deep and optimistic
   fast --- throughout recovery;
 \item the round's distinguished proposal, carrying a \emph{proposed
   root} \(A_p^r\), a \emph{parent} \(Z_p^r\), and a
   \emph{proposal block} \(B_p^r\), with
   \(A_p^r\preceq Z_p^r\prec B_p^r\); their derivation, acceptance
   guards, and the unobstructed-round conditions appear with
   Definitions~\ref{def:stable-root}--\ref{def:usable-honest-round},
   and this section uses only the ancestry just stated; and
 \item the \emph{relay point}: the fixed schedule time \(a_0\) at
   which the analyzed recovery window opens, defined with
   Definition~\ref{def:relay-frontier}; this section uses only
   that it is a cutoff time preceding the first counted round.
\end{enumerate}
This section and the next use only the properties stated here;
the derivation rules and acceptance guards appear in
Section~\ref{sec:recovery-proposals} and
Section~\ref{sec:available-confirmation}.

In plain words, the stable root is where a validator's votes walk from
during the round, and the action root bounds what its SG head signs at
the round's end.
\end{definition}

Recovery changes the local root supplied to Goldfish, replaces the
confirmation rule with the available confirmation of
Section~\ref{sec:available-confirmation}, and, on proposer equivocation,
discards both proposal copies --- the recovery wrapper around proposal
selection in Definition~\ref{def:recovery-proposal}.  It does not change
Goldfish's proposal-view merge, raw votes, or the finality state
transition of Section~\ref{sec:state-machine}.

\begin{assumption}[Synchronous recovery]\label{ass:recovery-network}
Recovery begins after \(\GST\).  Every honest validator is awake and
responsive.  A valid object broadcast or relayed by an honest validator at
time \(t\) reaches every honest validator strictly before \(t+\Delta\).
On receipt, honest validators relay every new valid block, SG head, raw
Goldfish vote, finality-gadget attestation, equivocation proof, and
required ancestry.  A view at a cutoff time contains everything received
at or before that time; because delivery is strict, an object held by one
honest validator at a cutoff is in every honest view strictly before the
cutoff one delivery bound later, and no argument below depends on the
treatment of an object arriving exactly at a deadline.

In plain words, delivery takes less than one bound, honest validators
relay everything on receipt, and boundary arrivals never matter.
\end{assumption}

\begin{definition}[Recovery round and four views]\label{def:recovery-timing}
Recovery round \(r\) has a globally unique number and contains at least two
slots of length \(4\Delta\), the slot layout of Goldfish's optimistic
fast confirmations, as the phase list below states.

\emph{Proposer selection.}
Ordinary Goldfish proposer selection is receiver-local: at vote time each validator takes
the winning eligible proposal among the round proposals in its own view,
as the source's algorithm prescribes, so two honest validators holding
different proposal sets may take different winners.  The \emph{round
proposer at} validator \(u\) is the signer of \(u\)'s local winner.
Round \(r\) \emph{has an honest round proposer} when one honest
validator's first-slot proposal is the local winner at every honest
validator; every use of ``honest round proposer'' in this paper names
this common-recognition event --- the source's honest-leader event ---
never the mere eligibility of an honest signer, whose proposal may lose
some receiver's local selection to a competing eligible proposal
delivered only there.  Let \(d_r\) be the first slot's proposal time.

\emph{Selection state and boundary exception.}
At the slot's proposal time \(d_r\), validator \(u\) derives the
\emph{selection state} \(\Sigma_{u,\mathrm{sel}}^r\) from the main store:
the \emph{public selection cutoff} is \(d_r\) itself, one schedule time
for every validator, and the state is derived from everything received
strictly before it.  The proposal is a scheduled action with the
inclusive deadline input \(M_I\) of the store definition; the selection
state is a \emph{derived} quantity, computed from the sub-view of that
input received strictly before \(d_r\), so the immutable record stores
the full inclusive input while the derivation discards boundary
receipts.  This is the one deliberate exception to the
inclusive view convention of Assumption~\ref{ass:recovery-network},
whose boundary-irrelevance sentence is scoped to delivery deadlines, not
to this schedule cutoff: the exclusion keeps the round's own proposal,
broadcast at \(d_r\), out of the state that chose its root, and the
single-signer rule covers the proposal as it covers every scheduled
action.  The grade
snapshot \(X_u^0\) at the same instant stays inclusive; a block graded
there but received exactly at \(d_r\) is outside the selection tree and
not choosable this round, and the choice rule remains a total function
of the two views.

\emph{Vote state and support freeze.}
At \(d_r+\Delta\), after processing the
proposal decision and its ordinary Goldfish context, it derives the
\emph{vote state} \(\Sigma_{u,\mathrm{vote}}^r\) and fixes the stable root
used by its ordinary Goldfish vote for the rest of the round, subject
only to the re-derivation rule of
Definition~\ref{def:rederivation} below.  The
\emph{support freeze} for the first slot is the cutoff
\(f_r=d_r+2\Delta\).

\emph{Action time and per-round states.}
At the round's SG/FG action time
\[
 a_r=d_r+4\Delta+2\Delta,
\]
two network-delivery bounds into the round's second slot, validator \(u\)
first processes every delivery due by \(a_r\), then completes the first slot's
available-confirmation evaluation of
Section~\ref{sec:available-confirmation} from that view, and then
derives the \emph{action state} \(\Sigma_{u,\mathrm{act}}^r\).  All
three per-round states, and the current state consulted at any slot
boundary by the re-derivation rule of Definition~\ref{def:rederivation},
are derived under the
activation filter of Definition~\ref{def:activation-filter}; the filter
is a fixed function of the store at
\(a_{r-1}\), so the activated finality outcomes are constant
across all of one validator's round-\(r\) states, and only the store
maximum and block tree can change between them;
with that one selection guard they are the
complete fork-choice states of Definition~\ref{def:finality-root}, including
the store-global \(\hmax\), viable subtree, and Simplex root.  A validator
whose vote-time stable root fails the action-state admission of
Definition~\ref{def:action-root} emits no SG head or
current-height pair in that round.

Each slot runs the phases
\[
 d\quad\text{proposal},\qquad
 d+\Delta\quad\text{vote},\qquad
 d+2\Delta\quad\text{confirmation input},\qquad
 d+3\Delta\quad\text{slot-view freeze}.
\]
The available confirmation of one slot is evaluated \(2\Delta\) into
the following slot.
The SG grades use snapshots of the preceding round's SG-head batch at
\[
 X_u^-:d_r-\Delta,\qquad
 X_u^0:d_r,\qquad
 X_u^1:d_r+\Delta,\qquad
 X_u^2:d_r+2\Delta .
\]
At \(d_r+2\Delta\), validator \(u\) fixes grade~0.  At the action time
\(a_r\), it completes the first slot's available-confirmation evaluation,
evaluates the strong G0 rule against its action state, and performs the
round's SG and FG action.  Every slot of the round runs ordinary Goldfish
with the round's stable root --- fixed at the vote time, subject only to
the ungraded-fallback re-derivation above.
The action in round \(r+1\), which grades round \(r\)'s SG heads, is the
\emph{closing action} for round \(r\)'s proposal.
The public round spacing satisfies
\[
  a_r+\Delta\leq d_{r+1}-\Delta,
\]
so every honest SG head released in round \(r\) reaches every honest
validator before the first grade view of round \(r+1\).

In plain words, grade 2 uses the proposal-time view, grades 1 and 0 use
the \(+\Delta\) and \(+2\Delta\) views, and grade 3 is read on the
intersection of the \(-\Delta\) and \(+\Delta\) views.  Goldfish votes at \(+\Delta\); SG and FG
wait until the available confirmation used by their gate has been evaluated.
\end{definition}

This is the slot layout of Goldfish's optimistic fast
confirmations~\cite{goldfish}: \(4\Delta\) slots with a dedicated
confirmation phase at \(d+2\Delta\), one delivery bound before the
slot-view handling; basic Goldfish instead merges and confirms in one
phase at \(+2\Delta\) of its \(3\Delta\) slot.  The recovery profile
changes that dedicated phase in two ways: the source's fast-confirmation
test is replaced by the support freeze, with the evaluation moved two
delivery bounds into the following slot under a live-participation
denominator, per Section~\ref{sec:available-confirmation}; and the
source's buffer-into-tree merge at that phase is dropped --- the receipt
buffer is not merged before the slot-view freeze at \(d+3\Delta\), as
Definition~\ref{def:recovery-goldfish-vote} states.  The changed rule and the
profile's other departures are listed in Section~\ref{sec:refinement}.

\begin{definition}[Mid-round re-derivation of the stable root]
\label{def:rederivation}
The vote-time stable root is subject to one exception for ungraded
roots: if the fixed root is ungraded --- the vote
state's Simplex root, whether as fallback or as an accepted ungraded
proposed root, which equals that selection by the stable-root rule of the
proposal section; the graded
alternative is specified with the same rule --- and, at a later slot
boundary of the round, the validator's
current state selects a Simplex root different from the fixed root, it
re-derives its stable root at that
boundary and uses the re-derived root for the round's remaining slots.
The test is deterministic: it compares the fixed root with the current
state's Simplex root, both locally derived.  After re-derivation, every
later reference to the round's vote-time stable root --- later-slot
walks, the action-root rule, and the action head --- names the re-derived
root.  The re-derived root is
finality-backed by Definition~\ref{def:finality-root}; it follows the
state's own Simplex selection, which may sit below the abandoned
justification when the selection rule switches from a justification to
the finalized root.  A graded stable root is never abandoned mid-round
while its grade stays
active, with one exit: if the root's grade becomes inactive under
Definition~\ref{def:active-grade} --- the root conflicts with a
finalized root whose evidence the validator has since processed --- the
root counts as ungraded from the next slot boundary on and is subject to
the same exception.  Every ungraded root
--- fallback or accepted ungraded proposed root --- is subject to the
exception.

In plain words, an ungraded stable root follows the validator's own selection
when that selection moves mid-round, and a graded root is dropped only
when finality evidence kills its grade.
\end{definition}

\begin{definition}[Activation filter]
\label{def:activation-filter}
A finality outcome takes effect in a validator's round-\(r\) selection,
vote, and action states when its evidence was processed by that
validator's preceding-round action cutoff \(a_{r-1}\).

For the first
counted recovery round, the preceding action cutoff is the relay point
of Definition~\ref{def:relay-frontier}: every pre-window object
has crossed it, so the first round's activation filter and the aged
viability test below are total there, and the one-round staggering bound
has its base case at that cutoff.

This is
implemented in the derivation, not merely asserted: the three per-round
states are derived by running the ordinary fork-choice derivation with
the \emph{activation filter}: a justification or finalized pair is
excluded from the greatest-justification and greatest-finalized selection
of every round-\(r\) state
--- and treated as evidence --- unless some block that was received
\emph{and accepted} by \(a_{r-1}\) has a derived
state recording that pair.  Both conditions are fixed once \(a_{r-1}\)
has passed: receipt times are the store's \(r_S\) values, and a block
queued behind missing ancestry at the cutoff does not activate the pair
this round even if its ancestry arrives mid-round --- honest relay
carries required ancestry with every block
(Assumption~\ref{ass:recovery-network}), so for honestly relayed blocks
receipt and acceptance coincide.  The filter is therefore a fixed
function of the store at \(a_{r-1}\).  The filtered derivation is the algorithm of
Definition~\ref{def:finality-root} with this one added selection guard;
\(\hmax\), the viable subtree, and the Simplex root are all computed
from the filtered selection.  All other
processing is unchanged.

In plain words, newly learned finality steers selection only from the
next round; Remark~\ref{rem:activation-stagger} below records the two
consequences used throughout.
\end{definition}

\begin{remark}[Staggering and output discipline under the filter]
\label{rem:activation-stagger}
An outcome processed by one honest validator by
\(a_{r-1}\) reaches every honest validator strictly before \(d_r\), so
activation staggering spans at most one round: whenever any honest
validator activates an outcome at round \(r\), every honest validator
activates it by round \(r+1\).  A round with a staggered activation has
a finality advance in flight; before the first recovery confirmation such
advances are bounded by the finalized heights below \(H\), and after
healing each one is finality progress.

Within a round a
newly completed outcome is evidence only for root selection.  The filter
delays selection, never inertness or signed outputs: a grade for a block
conflicting with
any processed finalized root --- activated or evidence-only --- is
inactive by Definition~\ref{def:active-grade}; a fixed stable root
whose grade has become inactive in this way counts as ungraded from the
next slot boundary on, so the ungraded-root re-derivation applies to it;
and the official-confirmation gate and fallback head of
Definition~\ref{def:official-confirmation} test processed finality
directly, so no signed SG output ever names a block conflicting with
evidence its signer has processed, even in the round before that
evidence activates for selection.
\end{remark}

\begin{definition}[Recovery candidate roots]
\label{def:recovery-context}
For any fork-choice state \(\Sigma\) derived from the main store, define
\[
 \mathcal C(\Sigma)=
 \{B\in\viableTree(\Sigma):
       \simplexRoot(\Sigma)\preceq B\}.
\]
This is a prefix-closed tree rooted at \(\simplexRoot(\Sigma)\).  Its
viability test always uses the store-global value \(\Sigma.\hmax\) from
Definition~\ref{def:finality-root}.  Starting a Goldfish walk at a deeper
root does not recompute a smaller maximum over that root's descendants.

Write \(\mathcal C^{-}_u(\Sigma)\) for the \emph{aged} variant of the
candidate tree, used exactly where this definition says so.
For a \emph{receiver's grade-root selection only}, the viability test
uses aged witnesses.  A block \(B\in\mathcal C(\Sigma)\) is in
\(\mathcal C^{-}_u(\Sigma)\) exactly when some block \(W\) with
\(B\preceq W\) was received \emph{and accepted} by that validator's
preceding-round action cutoff \(a_{r-1}\) and has derived
state-height \(\sigma[W].h\geq\Sigma.\hmax-1\).  Here \(W=B\)
is allowed, whether \(W\) is a leaf of any tree is irrelevant, and
\(\Sigma.\hmax\) is the deriving state's own current maximum.  The
receipt-and-acceptance bound is the same condition as the activation
filter, so a witness queued behind missing ancestry at the cutoff does
not count even if its ancestry arrives later.  A witness received or
accepted later takes effect one round later --- in particular, a young
branch that just raised \(\Sigma.\hmax\) can leave older branches,
and even the walk root, without aged membership for one round; the
walk-standing rule of Definition~\ref{def:walk-standing} below keeps
every walk defined regardless.

In plain words, receivers age their witnesses by one round while the
proposer does not, and every recovery choice uses the main store's
ordinary \(\hmax-1\) filter.
\end{definition}

\begin{definition}[Walk-standing rule]
\label{def:walk-standing}
Every recovery walk --- the first-slot aged walk and the current-tree
walks of the second and later slots alike --- starts at the round's
fixed root whether or not that root is in the walk's candidate tree:
aging, or a mid-round viability change, constrains which branches may
attract a walk, never where the walk stands.  A block with no viable
(for the first slot, aged-viable) candidate child returns itself,
exactly as ordinary GHOST returns a childless block; totality is
Lemma~\ref{lem:aged-walk-total} below.  This rule is the recovery
profile's one exception to the viability requirement of
Definition~\ref{def:finality-root}: a root fixed while viable may
still be walked from after a mid-round eviction, and the action-state
admission separately decides whether it may still anchor SG/FG
outputs.  The acceptance and selection tests of
Definition~\ref{def:stable-root} use the witness-based aged
membership of Definition~\ref{def:recovery-context}, with no
walk-root stipulation.

In plain words, a walk always has a place to stand; membership tests
gate acceptance, never the walk itself.
\end{definition}

\begin{definition}[Counting rule and proposal-path exemption]
\label{def:counting-rule}
In every recovery walk, the candidate tree constrains the walk's
choices, not the counted weights: at each fork, a child's weight
counts every known vote for blocks in that child's subtree of the full
known block tree, whether or not those blocks are in the candidate
tree.  Fork weights are therefore independent of each validator's
viability view, so equal vote views give equal weights.  (Goldfish has
no viability filter, so the counting rule agrees with the source's on
Goldfish's own trees.)  Blocks on the accepted proposal's path --- the
proposed root \(A_p^r\) through the parent \(Z_p^r\) to the
proposal block \(B_p^r\), and their ancestors --- are exempt from
aging in the acceptance checks and the first-slot walk: the accepted
proposal itself carries them, and receiver lag is the wrong concern
for the proposer's own path.  The containment clause of
Definition~\ref{def:usable-honest-round} is accordingly stated modulo
this exempted path; every exempted block except the proposal block
itself lies on the proposer's own selected chain.

In plain words, trees restrict choices while weights see every vote,
and the proposer's own path is never aged away.
\end{definition}

\begin{remark}[Where aging applies]
\label{rem:aged-scope}
The round proposer's own grade-root selection uses its current tree:
the proposer selects grade roots in
\(\mathcal C(\Sigma_{p,\mathrm{sel}}^r)\).  A receiver chooses
its grade-3 lower root at the vote time, when \(X_u^1\) exists,
within the aged tree
\(\mathcal C^{-}_u(\Sigma_{u,\mathrm{sel}}^r)\), as
Definition~\ref{def:grade-root-choice} specifies, and before voting checks
its final stable root against the aged tree
\(\mathcal C^{-}_u(\Sigma_{u,\mathrm{vote}}^r)\).  Every aged use
in this paper is one of the two receiver checks just named, the
first-slot vote-time walk, and Definition~\ref{def:stable-root}'s
items using these trees; every other \(\mathcal C(\Sigma)\) is the
unaged tree.  The height filter and every other viability use are
unchanged: the first slot's vote-time walk uses a candidate tree with
the same aged witnesses as the receiver's grade-root selection, so
with equal maxima and active roots a receiver's walk tree never
exceeds an honest proposer's selection tree.  Aging is a first-slot
device tied to proposal acceptance; the walks of the round's second
and later slots use the validator's current candidate tree at each
walk time, with the round's fixed stable root, under the same
walk-standing rule.  At \(a_r\), the later action-state check
separately decides whether the fixed root may affect the SG/FG
action; a validator's finality lock constrains its own
finality-gadget signature and is not an additional fork-choice root.
Under Assumption~\ref{ass:recovery-network}, a witness counted by an
honest receiver was received and accepted, with its full ancestry, by
\(a_{r-1}\), and every block of that ancestry was relayed on receipt
with its required ancestry; the containment this yields is
Lemma~\ref{lem:aged-containment} below.  At vote time, a receiver
unions the proposal's ordinary blocks and raw Goldfish votes with its
own frozen slot view, never an SG head or grade view.
\end{remark}

\begin{lemma}[Aged receiver trees are contained in honest proposer trees]
\label{lem:aged-containment}
Suppose an honest receiver and an honest round proposer have equal store
maxima, equal active finality outcomes, and equal processed finalized
evidence at their round-\(r\) cutoffs.  Then the receiver's aged
candidate tree --- its grade-root selection tree and its first-slot
vote-time tree, excluding the exempted proposal path --- is a subtree of
the proposer's selection tree.

In plain words, whatever a receiver treats as live has long since
reached the proposer, so an adversarially timed release can enrich the
proposer's tree, never a receiver's alone.
\end{lemma}

\begin{proof}
A block of the receiver's aged tree is placed there by a witness block
received and accepted, with its full ancestry, by the receiver's cutoff
\(a_{r-1}\).  Each block of that ancestry was received by the receiver
by \(a_{r-1}\) and relayed on receipt together with its required
ancestry (Assumption~\ref{ass:recovery-network}), so the proposer
received and accepted every one of them strictly before
\(a_{r-1}+\Delta\leq d_r\), and the witness is accepted in the
proposer's selection state.  With equal maxima, active finality
outcomes, and processed evidence, the viability threshold and the
selection root agree, so the witness places the same blocks in the
proposer's selection tree.
\end{proof}

\begin{lemma}[Every recovery walk is total]
\label{lem:aged-walk-total}
Every recovery walk --- the first-slot vote-time walk over an aged
candidate tree and every later-slot walk over the current candidate
tree --- is defined and terminates: it starts at the round's fixed
root, each step moves to a viable (for the first slot, aged-viable)
candidate child, and a block with no such child returns itself.  In
particular, a fixed graded root evicted from the current tree by a
mid-round maximum rise still yields a vote, from its own subtree.

In plain words, aging or a mid-round rise can starve a walk of
children, never of a place to stand.
\end{lemma}

\begin{proof}
The walk-standing rule of Definition~\ref{def:walk-standing} does
not test the root, each step strictly deepens within the finite
accepted tree, and the childless case returns the current block, the
ordinary GHOST behavior at a childless block.  Eviction changes no vote
weight, because the counting rule counts every known vote over the full
known tree.  Definedness and termination follow.
\end{proof}

\begin{definition}[SG head and eligible batch]\label{def:sg-head}
At round \(r\), validator \(i\) derives
\[
 S_{i,\mathrm{act}}^r=\simplexRoot(\Sigma_{i,\mathrm{act}}^r).
\]
A round-\(r\) SG head is the \(\mathsf{head}\) field \(B\neq\bot\) of
validator \(i\)'s signed round-\(r\) attestation.  An honest validator puts a
head in that field only when \(S_{i,\mathrm{act}}^r\preceq B\) and \(B\) is
viable in \(\Sigma_{i,\mathrm{act}}^r\).  The two conditions are normative here; that the head-production rule
of Definition~\ref{def:official-confirmation} emits only heads
satisfying them is recorded with that rule
(Remark~\ref{rem:head-checklist}).  The state and ancestry are validation
context and are not extra signed fields.  Two signed attestations from one
identity in one round are SG equivocation exactly when their head fields
differ and at least one is nonempty.

Round \(r>0\) grades exactly the heads signed in round \(r-1\).  Older and
newer heads are ineligible.  The first recovery round has an empty eligible
batch.

In plain words, every recovery round grades the head fields signed in the
preceding round.
\end{definition}

\begin{definition}[Direct and favorable support]\label{def:grade-support}
Fix one eligible batch.  For a view \(X\) and block \(B\), let \(D_X(B)\)
collect the identities \(i\) whose single eligible head \(B_i\) in \(X\)
satisfies \(B\preceq B_i\), with no SG equivocation by \(i\) in \(X\); an
identity with SG equivocation supplies no direct support.  Let \(\Phi_X(B)\)
collect the identities with an eligible head \(B_i\) in \(X\) satisfying
\(B\preceq B_i\), together with every identity that has an SG equivocation
in \(X\), each identity counted once.  Let \(D_u^{-,1}(B)\) collect the
identities whose unique eligible head is the same block \(B_i\) in both
\(X_u^-\) and \(X_u^1\), satisfies \(B\preceq B_i\), with no SG
equivocation in \(X_u^1\).  Define
\[
 \directSupport_X(B)=\sum_{i\in D_X(B)}w(i),\qquad
 \favorableSupport_X(B)=\sum_{i\in \Phi_X(B)}w(i),
\]
\[
 \directSupport_u^{-,1}(B)=\sum_{i\in D_u^{-,1}(B)}w(i).
\]
In plain words, direct support removes equivocators, while favorable support
preserves an earlier live witness when a later live equivocation appears.
\end{definition}

\begin{definition}[Four SG grades]\label{def:grades}
For every valid block \(B\), validator \(u\) computes
\begin{align*}
 \Gthree_u(B)&\Longleftrightarrow
      \directSupport_u^{-,1}(B)\geq m,\\
 \Gtwo_u(B)&\Longleftrightarrow
      \directSupport_{X_u^0}(B)\geq m,\\
 \Gone_u(B)&\Longleftrightarrow
      \favorableSupport_{X_u^1}(B)\geq m,\\
 \Gzero_u(B)&\Longleftrightarrow
      \favorableSupport_{X_u^2}(B)\geq m.
\end{align*}
All four predicates use the same fixed SG electorate, weights, round, and
eligible batch.  They are ancestor-closed.  G3 is a retrospective
predicate: it reads the \(-\Delta\) snapshot but becomes available only
at the \(+\Delta\) view, where it certifies that a grade-2 direct
majority already held at every honest proposal-time view.  The predicates are facts about
the signed batch and do not depend on a validator's finalized root.

In plain words, grades 3 and 2 are direct strict majorities, while grades 1
and 0 carry those observations through later equivocation.
\end{definition}

\begin{definition}[Grades used by fork choice]\label{def:active-grade}
At a state \(\Sigma\), validator \(u\) uses grade \(j\) for \(B\) only when
\(\mathsf G_j{}_u(B)\) holds, \(B\) is in the accepted subtree rooted at
\(\Sigma.F\), and \(B\) conflicts with no finalized root whose evidence
\(u\) has processed --- whether or not the activation filter of
Definition~\ref{def:activation-filter} has activated that outcome for
root selection.  A grade for a conflicting block remains
signed evidence but has no effect on the stable root or G0.

In plain words, grade delivery is computed from the signed batch, while
finalized-conflicting grades do not affect a validator's choices.
\end{definition}

\begin{lemma}[Direct-majority roots are compatible]
\label{lem:direct-roots-chain}
Two blocks with direct support at least \(m\) in one view are compatible.

In plain words, one view cannot put a direct strict majority on two forks.
\end{lemma}

\begin{proof}
The direct supporter sets of conflicting blocks are disjoint.  Their
combined weight would be at least \(2m>W\), which is impossible.
\end{proof}

\begin{corollary}[Grade-3 blocks lie on one chain]
\label{cor:g3-chain}
Two blocks with two-view direct support at least \(m\) at one validator
--- in particular two grade-3 blocks --- are compatible, so the deepest
grade-3 block of a candidate tree is unique whenever one exists.

In plain words, the grade-3 selection in the stable-root rule is total
and deterministic.
\end{corollary}

\begin{proof}
An identity in both blocks' two-view direct-support sets has one
eligible head, the same block in \(X_u^-\) and \(X_u^1\), and support
means that head lies in each supported block's subtree; one head cannot
descend from two conflicting blocks, so the sets are disjoint and their
combined weight would be at least \(2m>W\).  Compatible blocks form a
chain, and a chain's deepest member in a prefix-closed tree is unique.
\end{proof}

\begin{lemma}[Grade 3 reaches grade 2]\label{lem:g3-g2}
If one honest validator has grade~3 for \(B\), every honest validator has
grade~2 for \(B\).

In plain words, an early direct majority that remains unique reaches every
proposal-time view.
\end{lemma}

\begin{proof}
Let \(M\) be the heads counted by
\(\directSupport_u^{-,1}(B)\).  Validator \(u\) relays them at
\(d_r-\Delta\), so every honest \(X^0\) contains them.  If another honest
validator knew a different head from a signer in \(M\) by \(d_r\) ---
conflicting or not, a second distinct head is SG equivocation and removes
direct support --- its
relay would reach \(u\) by \(d_r+\Delta\), removing that signer from the
intersection support.  Thus the same weight at least \(m\) is direct in
every honest \(X^0\).
\end{proof}

\begin{lemma}[Grade 2 reaches grade 1]\label{lem:g2-g1}
If one honest validator has grade~2 for \(B\), every honest validator has
grade~1 for \(B\).

In plain words, a proposal-time direct majority becomes favorable everywhere
one delay later.
\end{lemma}

\begin{proof}
The observer relays its direct witness at \(d_r\).  By \(d_r+\Delta\), each
represented identity either retains one head in \(B\)'s subtree, or a second head
makes it an equivocator.  Both cases count favorably, preserving weight at
least \(m\).
\end{proof}

\begin{lemma}[Grade 1 reaches grade 0]\label{lem:g1-g0}
If one honest validator has grade~1 for \(B\), every honest validator has
grade~0 for \(B\).

In plain words, a grade-1 upper check reaches every final veto view.
\end{lemma}

\begin{proof}
The observer relays its favorable witness at \(d_r+\Delta\).  It reaches
every honest \(X^2\) by \(d_r+2\Delta\), and favorable weight cannot decrease
when another head from a represented identity arrives.
\end{proof}

\begin{lemma}[A grade-1 root is accepted everywhere by the grade-0 freeze]
\label{lem:g1-acceptance}
Under Assumption~\ref{ass:sg-fault-bound}, if a block \(R\) holds
raw grade~1 at an honest validator in round
\(r\), then every honest validator receives and accepts \(R\), with
its full ancestry, strictly before the grade-0 freeze \(d_r+2\Delta\).

In plain words, a grade-1 root cannot be invisible to any grade-0 test.
\end{lemma}

\begin{proof}
\(R\)'s favorable support has weight at least \(m>b\), so some
supporter is honest, and an honest identity never SG-equivocates
(Definition~\ref{def:stored-data}'s single-signer rule), so that
supporter lies in the head-containing class of
Definition~\ref{def:grade-support}: its eligible head \(B_i\)
satisfies \(R\preceq B_i\), is in its holder's \(X^1\) view,
received by \(d_r+\Delta\) and relayed on receipt with required
ancestry.  Every honest validator therefore receives \(R\) with its
ancestry strictly before \(d_r+2\Delta\) and accepts it under
Definition~\ref{def:total-raw-replay}.
\end{proof}

\begin{lemma}[A compatible batch removes every conflicting grade]
\label{lem:fresh-support-clears}
Under Assumption~\ref{ass:sg-fault-bound}, if every eligible honest head is
compatible with \(B\), no block conflicting with \(B\) has grade
\(3,2,1,\) or \(0\) in an honest view.

In plain words, the next grade computation has no stale grade on the other
branch.
\end{lemma}

\begin{proof}
An honest head cannot support a block conflicting with \(B\), and an
honest identity does not equivocate, so favorable support on a conflicting
block is at most \(b<m\).  Direct support is no larger, so all four
predicates fail.
\end{proof}

\paragraph{Assumption usage.}
The clearance lemma uses only
Assumption~\ref{ass:sg-fault-bound}, through \(b<m\).  Grade delivery and
direct-root compatibility use no fault bound.  None of these lemmas uses
the finality threshold \(q\) or the stronger healing bound \(3b<W\).
The full-participation clause of
Lemma~\ref{lem:root-no-regression} below also uses
\(W-b\geq m\), which likewise follows from
Assumption~\ref{ass:sg-fault-bound} alone: \(2b<W\) gives
\(b\leq\lceil W/2\rceil-1\), so
\(W-b\geq\lfloor W/2\rfloor+1=m\).

\section{Available confirmation}\label{sec:available-confirmation}

This section defines the time-shifted-quorum available confirmation for one
slot and proves its two cross-view properties: two available confirmations
cannot conflict, and a time-shifted majority is a majority in every honest
view.  Section~\ref{sec:finality-votes} builds the official-confirmation
gate on these lemmas, and Section~\ref{sec:healing} drives the healing
induction with them.  Both proofs use only honest relay and the delivery
bound of Assumption~\ref{ass:recovery-network}; no fault bound appears.

\begin{definition}[Participants, first votes, and visible equivocation]
\label{def:tsq-views}
Fix one slot with proposal time \(d\) and the committee given for that
slot.  Raw Goldfish votes are counted by identity, one unit per committee
member.  For a view \(X\) and committee member \(i\), say \(i\)
\emph{visibly equivocates in \(X\)} when \(X\) contains two distinct votes
of this slot signed by \(i\).  Member \(i\) \emph{supports} block \(B\) in
\(X\), written \(\operatorname{supp}_X(i,B)\), when \(X\) contains exactly
one vote of this slot signed by \(i\) and that vote names a block in
\(B\)'s subtree.  No reception order enters the definition, so
simultaneous deliveries need no tie-break.  A \emph{participant} in \(X\)
is a committee member with at least one vote of this slot in \(X\), timely
or not.

Validator \(u\)'s \emph{support view} \(V_u^-\) is its view at the
support freeze \(f=d+2\Delta\).  Its \emph{evaluation view}
\(V_u^+\) is its view at the evaluation time \(e=d+4\Delta+2\Delta\), two
delivery bounds into the following slot.  Views grow monotonically, so
\(V_u^-\subseteq V_u^+\).  For the first slot of recovery round \(r\),
\(d=d_r\), \(f=f_r\), and \(e=a_r\).

In plain words, support is a single visible vote inside the subtree --- a
double vote supports nothing --- and participation counts every voter seen
by evaluation time.
\end{definition}

\begin{definition}[Time-shifted-quorum available confirmation]
\label{def:tsq-confirmation}
At the evaluation time \(e\) of Definition~\ref{def:tsq-views}, validator
\(u\) computes for each block \(B\) the
\emph{time-shifted support}
\[
 \operatorname{ts}_u(B)=\bigl|\{i:\operatorname{supp}_{V_u^-}(i,B)\ \text{and}\
 \operatorname{supp}_{V_u^+}(i,B)\}\bigr|
\]
and the \emph{live participation}
\(n_u=|\{\text{participants in }V_u^+\}|\).
Block \(B\) is \emph{available-confirmed} at \(u\) when \(2\operatorname{ts}_u(B)>n_u\).
Support counts a vote for every ancestor of its target, so available
confirmation is downward closed: every ancestor of an available-confirmed
block is available-confirmed.

In plain words, the support is frozen early and cleaned of every double
vote visible late, while the participation is counted late; a strict
majority of the late count confirms.  The support comes from the early
freeze so that every counted vote has reached every honest validator by
the slot-view freeze; the participation comes from the late view so that
no honest validator counts this slot's votes against a larger set.
\end{definition}

\begin{lemma}[Two available confirmations cannot conflict]
\label{lem:tsq-uniqueness}
Under the delivery and relay rule of
Assumption~\ref{ass:recovery-network}, let honest validators \(u\)
and \(v\) evaluate the same slot.  If \(B\) is
available-confirmed at \(u\) and \(B'\) is available-confirmed at \(v\),
then \(B\) and \(B'\) are compatible.  In particular, with \(u=v\), the
blocks available-confirmed in one evaluation form a chain.

In plain words, one slot cannot produce two conflicting confirmations, in
the same or in different honest views, and each view confirms along one
chain.
\end{lemma}

\begin{proof}
Suppose \(B\) and \(B'\) conflict.  Every vote counted by \(\operatorname{ts}_u(B)\) was
received by \(u\) by the freeze \(f=d+2\Delta\) and relayed
under Assumption~\ref{ass:recovery-network}, so it is in every honest view
by \(d+3\Delta\); in particular its signer is a participant in \(V_v^+\).
The same holds with \(u\) and \(v\) exchanged.  Now suppose one member
\(i\) is counted by both \(\operatorname{ts}_u(B)\) and \(\operatorname{ts}_v(B')\).  Then \(u\) holds a
vote by \(i\) in \(B\)'s subtree and \(v\) holds a vote by \(i\) in
\(B'\)'s subtree.  Conflicting blocks have disjoint subtrees,
so these are two distinct votes by \(i\); both were relayed by
\(d+3\Delta\), so both are in \(V_u^+\) and in \(V_v^+\) at
\(e=d+6\Delta\), and \(i\) visibly equivocates in both evaluation views.
Then \(\operatorname{supp}_{V_u^+}(i,B)\) fails, contradicting the count
of \(i\) in \(\operatorname{ts}_u(B)\).  The two counted member sets are therefore
disjoint, and each lies among the participants of \(V_u^+\) and among the
participants of \(V_v^+\).  Hence
\(\operatorname{ts}_u(B)+\operatorname{ts}_v(B')\leq n_u\) and \(\operatorname{ts}_u(B)+\operatorname{ts}_v(B')\leq n_v\).  The two
confirmations give \(2\operatorname{ts}_u(B)>n_u\) and \(2\operatorname{ts}_v(B')>n_v\); adding them gives
\(2\operatorname{ts}_u(B)+2\operatorname{ts}_v(B')>n_u+n_v\geq2\bigl(\operatorname{ts}_u(B)+\operatorname{ts}_v(B')\bigr)\), a
contradiction.
\end{proof}

\begin{lemma}[A time-shifted majority is a majority in every honest view]
\label{lem:tsq-adoption}
Suppose \(B\) is available-confirmed at honest validator \(u\).  Let \(w\)
be honest and let \(X_w\) be any view that contains \(w\)'s slot-view
freeze view of this slot and that \(w\) holds no later than \(e-\Delta\).
Then every member counted by \(\operatorname{ts}_u(B)\) satisfies
\(\operatorname{supp}_{X_w}(i,B)\), and \(B\)'s supporters in \(X_w\)
strictly outnumber half of the participants in \(X_w\).

In plain words, a confirmed block already had a strict majority of plain
support, not equivocation credit, in every honest view used for the
following slot's votes.
\end{lemma}

\begin{proof}
Let \(i\) be counted by \(\operatorname{ts}_u(B)\).  Its vote in \(B\)'s subtree reached
every honest validator by \(d+3\Delta\), so it is in \(w\)'s slot-view
freeze view and therefore in \(X_w\).  Suppose \(i\) visibly equivocated
in \(X_w\).  Then \(w\) received two distinct votes by \(i\) no later than
\(e-\Delta\) and relayed them under Assumption~\ref{ass:recovery-network},
so both are in \(V_u^+\) at \(e\), and \(i\) visibly equivocates there ---
contradicting \(\operatorname{supp}_{V_u^+}(i,B)\) in the count of
\(\operatorname{ts}_u(B)\).  Hence \(X_w\) holds exactly one vote by \(i\); that vote is
the one in \(B\)'s subtree, so \(\operatorname{supp}_{X_w}(i,B)\) holds.

Every participant in \(X_w\) has a vote that \(w\) received by
\(e-\Delta\); \(w\)'s relay puts it in \(V_u^+\) by \(e\), so every
participant in \(X_w\) is a participant in \(V_u^+\), and the participants
of \(X_w\) number at most \(n_u\).  Hence \(B\)'s supporters in \(X_w\)
number at least \(\operatorname{ts}_u(B)>n_u/2\), which is at least half the participants
in \(X_w\), strictly.
\end{proof}

\begin{lemma}[A single-vote majority wins every fork]
\label{lem:fork-margin}
Let a vote-time fork-choice walk count, for each child of a fork, the
identities with a vote in that child's subtree, one unit per identity per
child, over a view with \(P\) participants.  Suppose \(N\) identities
each have a single vote, all in one designated child's subtree, with
\(2N>P\).  Then the designated child strictly beats every other child at
that fork.

In plain words, a strict majority of single votes in one subtree decides
the fork under the ordinary counting rule.
\end{lemma}

\begin{proof}
The designated child counts at least \(N\).  Every identity counted for
another child lies outside the \(N\) single-vote identities, and each
identity is counted at most once per child, so another child counts at
most \(P-N<N\).
\end{proof}

\begin{corollary}[A confirmed block captures every honest vote-time walk]
\label{cor:tsq-walk}
In the setting of Lemma~\ref{lem:tsq-adoption}, let \(w\) run a vote-time
fork-choice walk on \(X_w\) --- or let the evaluator \(u\) run one on its
own evaluation view \(V_u^+\) --- from a root compatible with \(B\),
counting for each child the identities with a vote in that child's
subtree, one unit per identity per child, as ordinary Goldfish does, and
suppose every block on the path from the root to \(B\)
lies in the walk's candidate tree.  The walk's output descends from \(B\).

In plain words, after a confirmation, an honest vote from any compatible
root lands inside the confirmed subtree.
\end{corollary}

\begin{proof}
For \(V_u^+\), every member counted by \(\operatorname{ts}_u(B)\) supports \(B\) there by
definition and the participants number \(n_u\), so the counting below
applies with \(X_w=V_u^+\).  If the root descends from \(B\), every output
does.  Otherwise the root is a strict ancestor of \(B\).  Fix any fork on the path from the root to
\(B\), with \(c\) the child toward \(B\) and \(c'\) any other child.
A vote counts for every ancestor of its target, so each of the at least
\(\operatorname{ts}_u(B)\) members of Lemma~\ref{lem:tsq-adoption} has its single vote in
\(c\)'s subtree; none of them visibly equivocates in the view.  The
participants of \(X_w\) number at most \(n_u<2\operatorname{ts}_u(B)\), so
Lemma~\ref{lem:fork-margin} makes \(c\) beat \(c'\) strictly.
The premise keeps \(c\) among the walk's candidates at every step, so the
walk chooses \(c\) at every such fork; the path is finite because slots
strictly increase along a chain, so the walk reaches \(B\) and continues
inside \(B\)'s subtree.
\end{proof}

\begin{lemma}[An all-honest batch confirms]
\label{lem:tsq-liveness}
Under Assumption~\ref{ass:goldfish-committees}, if every honest committee
member of one slot votes for a block in \(B\)'s subtree, then \(B\) is
available-confirmed at every honest validator at that slot's evaluation.

In plain words, honest unanimity always clears the time-shifted majority.
\end{lemma}

\begin{proof}
Honest votes are cast at \(d+\Delta\) and delivered strictly before the
freeze \(f=d+2\Delta\), so each honest member is counted in every honest
\(V^-\); honest members do not equivocate, so their support persists in
every \(V^+\).  Hence \(\operatorname{ts}_u(B)\geq\eta\), the honest committee count, at
every honest \(u\).  Participants are committee members, so
\(n_u\leq\eta+f'\) with \(f'\) the faulty committee count, and
\(f'<\eta\) by Assumption~\ref{ass:goldfish-committees}.  Then
\(2\operatorname{ts}_u(B)\geq2\eta>\eta+f'\geq n_u\).
\end{proof}

\begin{corollary}[Every confirmation has an honest supporter]
\label{cor:tsq-honest-witness}
Under Assumption~\ref{ass:goldfish-committees}, the members counted by any
available confirmation at an honest validator include an honest one.

In plain words, faulty votes alone cannot confirm.
\end{corollary}

\begin{proof}
Suppose all counted members are faulty, so \(\operatorname{ts}_u(B)\leq a\) with \(a\) the
number of faulty members whose votes \(u\) counted.  Counted members are
participants, and every honest member is a participant in \(V_u^+\) by
Lemma~\ref{lem:tsq-liveness}'s delivery argument, so
\(n_u\geq\eta+a\).  Then \(2a\geq2\operatorname{ts}_u(B)>n_u\geq\eta+a\) gives \(a>\eta\),
contradicting \(a\leq f'<\eta\).
\end{proof}

\begin{lemma}[A confirmed block reaches every honest store before the action]
\label{lem:tsq-common-knowledge}
If \(B\) is available-confirmed at an honest validator at evaluation time
\(e\), then \(B\) and its ancestry are in every honest validator's store by
\(d+3\Delta\), three delivery bounds before \(e\).

In plain words, a confirmation never concerns a block that some honest
validator lacks at its action.
\end{lemma}

\begin{proof}
The confirmation gives \(\operatorname{ts}_u(B)\geq1\), so some counted vote names a block
in \(B\)'s subtree.  That vote was interpretable in \(V_u^-\), so the
store's dependency rule placed its target block and full ancestry ---
including \(B\) --- in \(u\)'s store by the freeze \(f=d+2\Delta\).
Honest relay under Assumption~\ref{ass:recovery-network} delivers the vote
with its block and ancestry to every honest validator by
\(f+\Delta=d+3\Delta\).
\end{proof}

\begin{lemma}[Later confirmations extend every earlier honest confirmation]
\label{lem:tsq-monotonicity}
Let \(B\) be available-confirmed at an honest validator at slot \(s\)'s
evaluation.  Suppose that, in every slot after \(s\) up to some slot
\(s'\), every honest committee member's vote-time stable root is compatible with
\(B\), the path from that root to \(B\) lies in its vote-time candidate
tree, and Assumption~\ref{ass:goldfish-committees} holds.  Then every
block available-confirmed at any honest validator at the evaluations of
slots \(s\) through \(s'\) is compatible with \(B\), and every honest
committee vote in those slots after \(s\) names a block in \(B\)'s
subtree.

In plain words, once one honest validator confirms, no later confirmation
leaves that subtree while honest roots remain compatible with it.
\end{lemma}

\begin{proof}
At slot \(s\)'s evaluation, Lemma~\ref{lem:tsq-uniqueness} makes every
confirmation compatible with \(B\).  For slot \(s+1\), each honest
member's vote-time view contains its slot-view freeze view and holds only
material received by the vote time, one delivery bound before the
evaluation, so Lemma~\ref{lem:tsq-adoption} applies to it; its root and
candidate tree satisfy the hypothesis, so Corollary~\ref{cor:tsq-walk}
puts its vote in \(B\)'s subtree.  Inductively, suppose every honest vote
of slot \(k>s\) names a block in \(B\)'s subtree.  First, a block \(B'\)
conflicting with \(B\) cannot be available-confirmed at slot \(k\)'s
evaluation: an honest member's first vote lies in \(B\)'s subtree, which
is disjoint from \(B'\)'s, so every member counted for \(B'\) is faulty,
contradicting Corollary~\ref{cor:tsq-honest-witness}.  Second, honest
votes of slot \(k\) give \(B\)'s subtree at least \(\eta\) supporters in
every honest vote-time view of slot \(k+1\) --- honest votes are delivered
one bound after casting.  Let \(n\) be the participant count of that
vote-time view for slot \(k\), in the sense of
Definition~\ref{def:tsq-views}: every participant is a member of slot
\(k\)'s committee, so
\(n\leq\eta+f'<2\eta\), and Lemma~\ref{lem:fork-margin} at every fork keeps
every honest slot-\((k+1)\) vote in \(B\)'s subtree.  The induction covers
every slot through \(s'\).
\end{proof}

\paragraph{Assumption usage.}
Uniqueness, adoption, the walk corollary, and common delivery use no
fault bound and no proposer assumption; they do use the delivery and
relay rule of Assumption~\ref{ass:recovery-network}.  The liveness lemma, the
honest-supporter corollary, and monotonicity additionally use the
per-committee honest majority of
Assumption~\ref{ass:goldfish-committees}; monotonicity also takes the
compatible-root hypothesis, which Section~\ref{sec:healing} discharges
with the root-safety and clearance lemmas.  Together these give the
division of labor used below: a confirmation is unique across honest
views, retroactively certifies a plain-support majority in every honest
vote-time view, is in every honest store one delivery bound after its
freeze,
and, while honest roots stay compatible, captures every later vote and
confirmation.  The confirmation numerator drops visible equivocators, while ordinary
vote-time head counting counts an identity once in each child subtree that
contains one of its votes; the walk corollary's margin holds under that
source counting because counted supporters never visibly equivocate.

The schedule's margins differ sharply between the two core lemmas.
Uniqueness retains three delivery bounds of slack.  The adoption lemma is
exactly tight at both ends: the support freeze \(f=d+2\Delta\), the
slot-view-freeze baseline \(d+3\Delta\), and the view bound \(e-\Delta\)
all sit at their boundary values, and shrinking any gap between them
admits a small counting counterexample.  The strict delivery bound of
Assumption~\ref{ass:recovery-network} absorbs every boundary case: an
object held by an honest validator at one cutoff lands strictly before
the next.  A deployment whose delivery bound is not strict can insert one
additional delivery bound before the evaluation; no lemma above is
affected.

\begin{lemma}[The exact batch preserves the honest proposal]
\label{lem:recovery-rooted-persistence}
Suppose every honest committee member in one slot votes for a descendant
of \(B\).  In every later slot, suppose every honest committee member's
fixed stable root is compatible with \(B\), the path from that root to \(B\)
lies in its vote-time candidate tree, and
Assumption~\ref{ass:goldfish-committees} holds.  Every honest raw-vote
target and fork-choice head in those slots descends from \(B\).

In plain words, the exact batch persists while the global height filter keeps
the proposal and later stable roots do not cross it.
\end{lemma}

\begin{proof}
The all-honest batch gives \(B\)'s subtree at least \(\eta\) supporters,
against at most \(\eta+f'<2\eta\) participants, in every honest vote-time view
of the following slot: honest votes are delivered one bound after casting,
honest members do not equivocate, and every participant of a vote-time
view is a committee member.  The compatible-root and candidate-tree
premises then let Lemma~\ref{lem:fork-margin}, applied at every fork on
the path with numerator \(\eta\) against at most \(\eta+f'<2\eta\)
participants, place every honest vote of that slot in \(B\)'s subtree.  Induction over the
stated slots, exactly as in Lemma~\ref{lem:tsq-monotonicity}, extends the
conclusion to every later slot.
\end{proof}

The exact vote batch alone does not prove its own viability.  A different
branch can raise the main store's \(\hmax\) and remove \(B\) from
\(\viableTree(\Sigma)\).  Every later use of
Lemma~\ref{lem:recovery-rooted-persistence} must therefore provide an
independent height-filter argument: whenever the global maximum rises to
\(x\), an accepted descendant of \(B\) must have state-height at least
\(x-1\).

\begin{lemma}[A persistent proposal is available-confirmed at the next evaluation]
\label{lem:recovery-rooted-confirmation}
If Lemma~\ref{lem:recovery-rooted-persistence} holds through a slot, then
\(B\) is available-confirmed at every honest validator at that slot's
evaluation, and every block available-confirmed there, at any honest
validator, is compatible with \(B\).

In plain words, persistence turns honest unanimity into a confirmation one
evaluation later, and into the only confirmation.
\end{lemma}

\begin{proof}
The persistent batch makes every honest committee member of that slot vote
for a descendant of \(B\), so Lemma~\ref{lem:tsq-liveness} available-confirms
\(B\) at every honest validator at the slot's evaluation, and
Lemma~\ref{lem:tsq-uniqueness} makes every block confirmed at that
evaluation compatible with \(B\).
\end{proof}


\section{Recovery proposal and ordinary Goldfish}\label{sec:recovery-proposals}

This section defines the stable root and the ordinary Goldfish proposal
interface.  The two bridge lemmas at the end of
Section~\ref{sec:available-confirmation} show how one exact honest vote
batch becomes
an available confirmation.


\begin{definition}[Grade-root choices]\label{def:grade-root-choice}
Write
\[
 S_{u,\mathrm{sel}}^r=
 \simplexRoot(\Sigma_{u,\mathrm{sel}}^r),
 \qquad
 S_{u,\mathrm{vote}}^r=
 \simplexRoot(\Sigma_{u,\mathrm{vote}}^r),
 \qquad
 S_{u,\mathrm{act}}^r=
 \simplexRoot(\Sigma_{u,\mathrm{act}}^r).
\]
\emph{Proposer choice.}
At \(d_r\), proposer \(p\) chooses the deepest block
\(A_p^r\in\mathcal C(\Sigma_{p,\mathrm{sel}}^r)\) for which it has grade~2
and which conflicts with no finalized root whose evidence the proposer
has processed --- the same processed-finality guard every grade use
obeys --- and supplies a \emph{grade-2 witness}: a set of signed preceding-round SG
heads of weight at least \(m\), each supporting \(A_p^r\).  Witness
validity is positive-only --- signatures, round, and eligibility --- and
never asserts the absence of other messages from the proposer's view;
the receiver's own raw \(\Gone_u\) predicate, whose favorable counting
tolerates equivocation copies the proposer may have omitted, supplies the
semantic acceptance guard.  If there is no graded block, the proposer
proposes the selection-state Simplex root \(S_{p,\mathrm{sel}}^r\) as an
ungraded base root.

\emph{Receiver lower root.}
Throughout,
call a graded block \emph{strict} at a validator when it strictly
descends from that validator's own vote-state Simplex root
\(S_{u,\mathrm{vote}}^r\).  At \(d_r+\Delta\), validator \(u\) chooses
the deepest \emph{strict} grade-3 block in the aged tree
\(\mathcal C^{-}_u(\Sigma_{u,\mathrm{sel}}^r)\) that remains in the aged
tree \(\mathcal C^{-}_u(\Sigma_{u,\mathrm{vote}}^r)\) --- unique by
Corollary~\ref{cor:g3-chain} --- or
\(S_{u,\mathrm{vote}}^r\) if no strict grade-3 block survives both
trees.  A grade-3 block equal to the selection itself is subsumed by the
fallback: the root is then the selection, classified as ungraded for
retention and re-derivation, so the two clauses name one rule.
Call this lower root \(L_u^r\).  Thus a grade made irrelevant by a newly
learned finalized root is not retained as a live lower bound.  These fallback
clauses make both selections total.

In plain words, the proposer offers its deepest grade-2 root or its own
selection, and each receiver fixes its own grade-3 lower bound.
\end{definition}

\begin{definition}[Stable root]\label{def:stable-root}
The vote-time stable root is derived from the proposal and the grade
roots of Definition~\ref{def:grade-root-choice} by the acceptance
guards and fallback below.

\emph{Acceptance guards.}
A proposal is \emph{timely} at receiver \(u\) when \(u\) receives it
strictly before its vote time \(d_r+\Delta\); timeliness and proposer
equivocation are decided in \(u\)'s vote-time view, locally and finally.
For a proposal from \(u\)'s round proposer that is timely at \(u\), with no
second distinct signed round proposal by the same proposer in \(u\)'s
vote-time view, let \(Z_p^r\)
be its parent and \(B_p^r\) its block.  The receiver accepts the proposal
exactly when:
\begin{enumerate}
 \item either \(L_u^r\preceq A_p^r\preceq Z_p^r\) and
 \(S_{u,\mathrm{vote}}^r\preceq Z_p^r\) --- and then
 \(S_{u,\mathrm{vote}}^r\preceq A_p^r\) holds automatically, by the
 first fact of Remark~\ref{rem:stable-root-coherence} below; or
 \(A_p^r\prec S_{u,\mathrm{vote}}^r\preceq Z_p^r\) with
 \(S_{u,\mathrm{vote}}^r=\Sigma_{u,\mathrm{vote}}^r.F\) --- the
 receiver's own newer finalized root already lies on the proposal path
 --- and \(L_u^r\preceq Z_p^r\);
 \item either the raw predicate \(\Gone_u(A_p^r)\) holds, \(A_p^r\)
 conflicts with no finalized root whose evidence the receiver has
 processed --- raw counting stays favorable to equivocation copies, but
 every grade use respects processed finality, as
 Definition~\ref{def:active-grade} requires --- and the proposal
 has a valid grade-2 witness, or \(A_p^r\) is ungraded and equals the
 receiver's own vote-time Simplex selection \(S_{u,\mathrm{vote}}^r\) ---
 in particular the genesis block, justified and finalized by stipulation,
 exactly when that is the selection.  The proposal may carry
 justification or finality certificates; the receiver processes them as
 ordinary deliveries, and under the activation filter of
 Definition~\ref{def:activation-filter} they enter root selection only at
 the next round boundary;
 \item the deepest block among
 \(\{A_p^r,L_u^r,S_{u,\mathrm{vote}}^r\}\), denoted
 \(R_{u,\mathrm{vote}}^r\), and the proposal block \(B_p^r\) are both in
 the aged tree \(\mathcal C^{-}_u(\Sigma_{u,\mathrm{vote}}^r)\), with
 the proposal-path exemption of Definition~\ref{def:counting-rule}.
\end{enumerate}
The deepest root is used for ordinary Goldfish for the rest of the
round; Remark~\ref{rem:stable-root-coherence} below records why that
deepest root is unique.

\emph{Failed-proposal fallback.}
If the proposal fails these checks --- or no proposal is recognized,
including the discard-both case --- the vote-time stable root is \(L_u^r\)
when \(L_u^r=S_{u,\mathrm{vote}}^r\), or when the raw predicate
\(\Gone_u(L_u^r)\) holds, \(L_u^r\) conflicts with no finalized root
whose evidence the receiver has processed, and \(L_u^r\) is in the aged
tree \(\mathcal C^{-}_u(\Sigma_{u,\mathrm{vote}}^r)\).  In every such case, write
\(R_{u,\mathrm{vote}}^r\) for the vote-time stable root so derived; the
action-root rule and the fallback SG-head rule consume this notation.  Otherwise it is
\(S_{u,\mathrm{vote}}^r\).

In plain words, G3 is the receiver's lower bound, G2 is the proposer's
witness, and G1 is the receiver's upper acceptance check.  A receiver never
moves behind its lower root.  If its own Simplex selection is a newer root on the
honest proposal's path, it starts farther forward on that same path rather
than rejecting the proposal.  The raw G1 predicate authenticates a proposed
root that may now be an ancestor of the receiver's finalized root; that old
root is not itself used for fork choice or G0.  An ungraded proposer root
never moves a receiver at all: it is accepted only when it already equals
the receiver's own selection, so every strict forward proposed root needs
the graded branch.  A root on
another branch makes the validator use its ordinary fallback.
\end{definition}

\begin{remark}[Coherence facts for the stable-root rule]
\label{rem:stable-root-coherence}
Four facts about
Definitions~\ref{def:grade-root-choice}, \ref{def:stable-root},
and~\ref{def:action-root} are used throughout.
First, \(L_u^r\succeq S_{u,\mathrm{vote}}^r\) always: both branches
of \(L_u^r\)'s definition select within
\(\mathcal C(\Sigma_{u,\mathrm{vote}}^r)\), which is rooted at the
selection.  Hence the first disjunct of item~1 forces
\(S_{u,\mathrm{vote}}^r\preceq A_p^r\), and the second disjunct is
the only way a root of the receiver may lie strictly after \(A_p^r\):
such roots are exactly a newer finalized selection and its grade-3
descendants.  Second, all three roots in item~1 lie on one chain, so the
deepest root \(R_{u,\mathrm{vote}}^r\) is unique.  Third, in the
action-root switch case, a grade-3 block classified as the equality
fallback does hold raw grade~1 --- its two-view direct supporters are in
the favorable-support set at \(X^1\) --- so retention there is a
real outcome.  Fourth, a proposal is never an authority for a root the
receiver's own state does not select: the receiver checks carried
certificates, never a claim about the proposer's store, and it rejects
an ungraded root different from its own selection, so the only strict
forward move \(S_{u,\mathrm{vote}}^r\prec A_p^r\) is grade-backed,
through the receiver's raw \(\Gone_u\) check and the grade-2
witness.
\end{remark}

\begin{lemma}[No forward move at the action]
\label{lem:no-forward-move}
Within one round, the Simplex selection of every per-round state and
every slot-boundary state is the round's activated justification \(J\)
or its activated finalized root \(F\preceq J\), and it moves at most
once, from \(J\) to \(F\).  Consequently
\(S_{u,\mathrm{act}}^r\preceq S_{u,\mathrm{vote}}^r\preceq
R_{u,\mathrm{vote}}^r\) for the root \(R_{u,\mathrm{vote}}^r\) that
Definition~\ref{def:stable-root} derives at vote time, and a
re-derived root of Definition~\ref{def:rederivation}, being the
re-deriving boundary's own selection, also has
\(S_{u,\mathrm{act}}^r\) at or below it.  A re-derived root may sit
strictly below \(S_{u,\mathrm{vote}}^r\) --- the abandoned
justification --- and no invariant claims otherwise.

In plain words, selections only move backward within a round, so the
action never needs a forward move, while a re-derivation may
legitimately move the walk root backward past the vote-time selection.
\end{lemma}

\begin{proof}
Within one round the activated pairs are constant and \(\hmax\) is
nondecreasing, so the selection of Definition~\ref{def:finality-root}
at any state of the round is the fixed \(J\) or the fixed \(F\) ---
an ancestor of \(J\), by the \(F\)-filter of that derivation ---
and once \(\hmax\) exceeds \(h_j+1\) it cannot return, so the
selection moves at most once, from \(J\) to \(F\).  The action
state derives last, so its selection sits at or below every earlier
selection of the round, including any re-deriving boundary's; this
covers the re-derived root, which is exactly that boundary's selection.
For the vote-time root, \(R_{u,\mathrm{vote}}^r\) is the deepest of
three blocks on one chain that include the vote-state selection, and
each fallback root is the lower root or the selection itself; every
case sits at or after \(S_{u,\mathrm{vote}}^r\) by the first fact of
Remark~\ref{rem:stable-root-coherence}.
\end{proof}

\begin{definition}[Action root and admission]\label{def:action-root}
At \(a_r\), the action root is the round's vote-time stable root ---
the re-derived root, after a mid-round re-derivation.  No forward move
is ever required: the action-state selection precedes the round's
vote-time root and any re-derived root, by
Lemma~\ref{lem:no-forward-move}.  The action evaluates the following
cases in order; the first that applies governs.
\begin{enumerate}
 \item \emph{Proposal-free.}  If no unique proposal was accepted ---
   a failed proposal, no timely proposal, or the discard-both case ---
   or an accepted proposal block has left the action candidate tree
   while the root ancestry holds --- a mid-round rise evicted the
   proposal --- the action reads no proposal block at all: the ancestry
   condition of case~2 is vacuous, the validator sets
   \(R_{u,\mathrm{act}}^r=R_{u,\mathrm{vote}}^r\), and only
   case~3's membership and admission tests apply.  A superseded
   proposal costs its proposer the round's context, never the validator
   its outputs.
 \item \emph{Off-path root.}  If the vote-time root or
   \(S_{u,\mathrm{act}}^r\) is not an ancestor of \(B_p^r\) --- in
   particular when a re-derivation moved the vote-time root off the
   proposal's chain --- no action root is defined and the validator
   emits no SG head or current-height pair.
 \item \emph{Admission.}  Otherwise the receiver uses the action root
   for SG and FG exactly when \(B_p^r\) --- in the proposal case ---
   and \(R_{u,\mathrm{act}}^r\) remain in
   \(\mathcal C(\Sigma_{u,\mathrm{act}}^r)\), and the action root
   either equals the action-state Simplex root or has raw grade~1
   while conflicting with no finalized root whose evidence the
   validator has processed by the action --- finality evidence
   processed after the vote thus withdraws a root from SG/FG use even
   when no slot boundary remains to re-derive it --- and the validator
   abstains otherwise.  When the vote-time root is an ungraded
   fallback equal to the round's justification selection and the
   action state has switched to the finalized root, this same
   admission decides: retention exactly when the root stays in the
   action candidate tree with raw grade~1 free of processed-finality
   conflict.  Either outcome is a deterministic function of the
   deadline view, and the switch requires a mid-round maximum move,
   one of the separately handled cases.
\end{enumerate}
On failure of any case, every output of the action is fixed --- empty
head, empty pair --- so the action is total with or without a proposal.
The vote-time stable root remains fixed for ordinary Goldfish, subject
only to the ungraded-fallback re-derivation of
Definition~\ref{def:rederivation}; the SG/FG action root is that same
root when the admission accepts it.

In plain words, the action first decides whether a proposal is in play,
then whether the roots are on its chain, then whether the root is
admissible; a validator abstains only on a real conflict, never merely
because a proposal was superseded.
\end{definition}


\begin{definition}[Recovery proposal]\label{def:recovery-proposal}
The first-slot proposal contains
\[
 (r,\mathcal O_p^r,\mathcal V_{p,\rm GF}^r,
   A_p^r,\mathsf{tag}_p^r,W_p^r,Z_p^r,B_p^r,\pi_p^r),
\]
where \(\mathsf{tag}_p^r\in\{\textsf{graded},\textsf{ungraded}\}\)
declares which branch of acceptance item~2 the proposer claims, and
\(W_p^r\) is the grade-2 witness when the tag is \textsf{graded} and
absent otherwise.  The tag binds the branch: a \textsf{graded} tag
without a valid witness, or an \textsf{ungraded} tag whose root differs
from the receiver's own selection, simply fails item~2 --- no case is
left open by a proposer that mismatches tag and witness.
The block context \(\mathcal O_p^r\) is the set of ordinary Goldfish
blocks the proposer attaches so that every receiver can connect and
validate the proposal: the ancestry of \(B_p^r\) down to a block the
receiver may already hold, together with every block referenced as a
target or parent by a vote in \(\mathcal V_{p,\rm GF}^r\) --- the
carried view is \emph{consistent} in the source's sense: every
reference resolves within the carried view or to a block every honest
validator already holds.  A proposal whose context does not connect its
block, or whose carried view is inconsistent, is invalid.  The available-vote view \(\mathcal V_{p,\rm GF}^r\) is
the proposer's set of raw Goldfish votes for the preceding slot; it is
the input of the proposal-view merge below.  Both fields carry ordinary
Goldfish data only --- no SG heads and no grades.  A graded proposed stable root has its separate grade-2 witness; an
ungraded base root is identified as such, and by item~2 of
Definition~\ref{def:stable-root} it is accepted only when it equals the
receiver's own vote-time Simplex selection.  Certificates carried with
the proposal are finality-state evidence, not an SG head view; the
receiver processes them as ordinary deliveries and never checks a claim
about the proposer's store.  The proposal block is valid, has parent \(Z_p^r\), and satisfies
\[
 S_{p,\mathrm{sel}}^r\preceq A_p^r\preceq Z_p^r\prec B_p^r.
\]
The credential \(\pi_p^r\) is the ordinary Goldfish proposer credential: the
witness by which ordinary proposer selection recognizes this proposal for that
slot.  How proposers are selected is outside this paper's scope, as with the
committees of Assumption~\ref{ass:goldfish-committees}.

A receiver validates the proposal and then applies Goldfish's ordinary
proposal-view merge.  It does not reject a Byzantine proposal merely because
its own GHOST calculation would choose a different parent.  An honest
proposer chooses \(Z_p^r\) by GHOST on its merged ordinary view,
choosing within its own selection-state candidate tree
\(\mathcal C(\Sigma_{p,\mathrm{sel}}^r)\) under the counting rule of
Definition~\ref{def:counting-rule} --- the same viability-filtered
walk every walk in this paper runs --- and extends
that parent with \(B_p^r\).  Honest receivers relay every timely round
proposal.  If a receiver's vote-time view contains two distinct signed
round proposals by its locally winning proposer, it uses neither as the
distinguished proposal and derives its vote-time stable root by the
failed-proposal fallback of Definition~\ref{def:stable-root}.  The
decision is a function of the receiver's vote-time view and is immutable
with the vote action; equivocation evidence arriving later does not
reopen it, though the proposals' blocks and votes remain ordinary
evidence.  Under a Byzantine proposer, two honest receivers may decide
differently --- one accepting a copy whose twin reached the other in
time --- and every acceptance is individually guarded by
items~1--3 of Definition~\ref{def:stable-root}; the results that need
common recognition state it as a premise
(Definition~\ref{def:usable-honest-round}).  An honest proposer's
unique proposal, broadcast at \(d_r\), is timely at every honest
receiver by Assumption~\ref{ass:recovery-network} and unpaired by the
single-signer rule of Definition~\ref{def:stored-data}: the proposal is
a scheduled action, and one honest identity releases at most one signed
result for it.  Timeliness and unpairedness do not alone make it the
local winner --- a competing eligible proposal delivered to a subset of
receivers can outrank it there --- so common recognition is exactly the
honest-round-proposer event of Definition~\ref{def:recovery-timing},
assumed to recur by Assumption~\ref{ass:frontier-opportunities} and
discharged by Lemma~\ref{lem:common-recognition}.

The discard-both rule is a recovery wrapper around proposal selection.  Once
the wrapper selects one proposal or none, block processing, the
proposal-carried ordinary view merge, raw Goldfish voting, and expiry
follow the ordinary Goldfish rules.  Confirmation does not: throughout
recovery, the confirmation rule is the available confirmation of
Section~\ref{sec:available-confirmation}, which replaces both of the
source's confirmation rules --- the \(\kappa\)-deep rule and the
optimistic fast confirmation whose slot layout this profile borrows ---
as stated in
Assumption~\ref{ass:recovery-goldfish}.

In plain words, an honest proposer supplies the ordinary merged view;
equivocation makes the recovery wrapper proceed without a distinguished
proposal.
\end{definition}

\begin{definition}[Unobstructed honest proposal]\label{def:usable-honest-round}
An honest round proposal with parent \(Z\) and block \(B\) is
\emph{unobstructed} when every honest validator recognizes it at vote time,
and the path checks in Definition~\ref{def:stable-root} give it a vote-time
root \(R_{u,\mathrm{vote}}^r\) on the chain segment from \(A_p^r\) through
\(Z\).  The required grade-2 witness is valid on the graded branch, and an
ungraded base root equals each receiver's own vote-time selection, and
both \(R_{u,\mathrm{vote}}^r\) and \(B\) are in the receiver's aged
vote-state candidate tree
\(\mathcal C^{-}_u(\Sigma_{u,\mathrm{vote}}^r)\) with the
proposal-path exemption --- item~3 of Definition~\ref{def:stable-root}.  In addition, every honest receiver's first-slot
candidate tree, excluding the proposal path that
Definition~\ref{def:counting-rule} exempts from aging, is contained
in the proposer's selection tree --- supplied by
Lemma~\ref{lem:aged-containment} whenever their store maxima, active
finality outcomes, and processed
finalized evidence agree; the exclusion is necessary
because \(B\) itself is created at \(d_r\) and lies in no selection
tree --- so the walks of
Lemma~\ref{lem:forward-root} run in subtrees of the proposer's up to
that path.  At action time, every newly selected Simplex root is an
ancestor of \(B\), so the action-state admission of
Definition~\ref{def:action-root} accepts the vote-time root.  A root that conflicts with \(B\),
or that strictly descends from it, supersedes the proposal.

This is not an additional recurring liveness assumption.
Lemma~\ref{lem:honest-proposal-confirms}, in the first-confirmation
development of Section~\ref{sec:first-confirmation}, proves that, before the first recovery
confirmation, the only conflicting root that can supersede such a proposal
is the unique old height-\(H-1\) justification.

In plain words, an honest proposal is unobstructed when every receiver can
start on its parent chain and every newer selection stays on the proposal
chain.
\end{definition}

\begin{lemma}[Common recognition at honest-proposer rounds]
\label{lem:common-recognition}
If round \(r\) has an honest round proposer --- the common-recognition
event of Definition~\ref{def:recovery-timing} --- then every honest
validator recognizes its proposal as the round's distinguished proposal
at vote time.

In plain words, at a common-recognition round the honest proposal is the
one proposal every honest voter acts on.
\end{lemma}

\begin{proof}
Three conditions make a proposal the distinguished proposal at a
receiver (Definition~\ref{def:recovery-proposal}): the receiver's
vote-time view contains it, it is the receiver's local selection winner,
and the discard-both wrapper does not fire.  The honest proposal,
broadcast at \(d_r\), reaches every honest validator strictly before
the vote time \(d_r+\Delta\) by
Assumption~\ref{ass:recovery-network}, so every vote-time view contains
it.  It is the local winner at every honest validator by the premise
itself; this clause is assumed, not derived, because timeliness and
honesty do not exclude a competing eligible proposal, released only to a
subset of receivers, that outranks the honest one under the
receiver-local selection of Definition~\ref{def:recovery-timing}.
Last, the proposal is a scheduled action, so the single-signer rule of
Definition~\ref{def:stored-data} yields at most one signed round
proposal from this honest proposer; no vote-time view holds two, and the
discard-both wrapper of Definition~\ref{def:recovery-proposal} does not
fire.
\end{proof}

\begin{lemma}[An unobstructed honest proposal puts every vote root on its parent chain]
\label{lem:honest-proposer-stable-root}
For an unobstructed honest proposal, every honest validator's vote-time
stable root \(R_{u,\mathrm{vote}}^r\) satisfies
\[
 A_p^r\preceq R_{u,\mathrm{vote}}^r\preceq Z_p^r.
\]

In plain words, a receiver may start farther forward because of its own
newer selection, but it starts on the path the honest proposer already
chose.
\end{lemma}

\begin{proof}
Definition~\ref{def:usable-honest-round} supplies every acceptance guard in
Definition~\ref{def:stable-root}.  That definition makes
\(R_{u,\mathrm{vote}}^r\) the deepest of three roots that are each ancestors
of \(Z_p^r\), including \(A_p^r\) itself.  The stated ancestry follows.
\end{proof}

\begin{assumption}[Goldfish interface during recovery]
\label{ass:recovery-goldfish}
Raw Goldfish votes are separate from weighted SG heads and FG attestations.
Goldfish votes are cast by the slot committees of
Assumption~\ref{ass:goldfish-committees} and counted one unit per member.
Recovery retains ordinary Goldfish proposer selection, proposal-view merge
after the recovery wrapper has selected a proposal, one-slot vote expiry,
vote equivocation handling, deterministic tiebreak, and the future-slot
in-limbo rule --- a block whose slot has not begun stays outside every
derived state, the slot-eligibility clause of
Definition~\ref{def:total-raw-replay}; the confirmation
rule of the recovery profile is Definition~\ref{def:tsq-confirmation}.

The proof imports one view-merge interface.  Suppose a recognized honest
round proposer holds, at proposal time, a view containing every vote and
block held by any honest committee member at the preceding slot-view
freeze; it starts from stable root \(A\), chooses GHOST parent \(Z\) on
its merged proposal view, and extends \(Z\) with block \(B\); every honest
committee member receives the proposal by the vote time, merges exactly
the proposal's carried view into its own frozen slot view, and no honest
committee member wakes or alters its frozen view between that freeze and
its vote.  Then every honest committee member's merged vote-time view
equals the proposer's published view together with \(B\), and its GHOST
walk from the proposer's root \(A\) votes for the exact block \(B\).

The source~\cite{goldfish} proves the corresponding
genesis-rooted property --- its view-merge lemma, in the analysis
section of the Goldfish paper --- under these hypotheses, including the
honesty of the proposer; a merely valid proposal carries no such
guarantee.  The rooted form stated here is this paper's interface, not
the source's; Remark~\ref{rem:view-merge-discharge} records why
recovery meets the hypotheses and how the rooted form follows.

In plain words, an honest proposer that includes everything and is
recognized by everyone hands every honest voter its exact view, so every
honest vote names its block.  Safety of confirmations does not rest on
this assumption; it comes from Section~\ref{sec:available-confirmation}.
\end{assumption}

\begin{remark}[Discharging the view-merge hypotheses during recovery]
\label{rem:view-merge-discharge}
The proposer-side containment holds under
Assumption~\ref{ass:recovery-network}: the preceding slot-view freeze is
one delivery bound before the proposal, so everything an honest member
held at that freeze has reached the proposer.  The no-wake clause is
immediate during recovery, where every honest validator is awake
throughout.  The rooted form follows because equal merged views make the
walk from the proposer's root replay the proposer's own choice sequence;
that lift is part of this paper's interface and is not attributed to the
source.
\end{remark}

\begin{lemma}[A forward root on the proposer's path returns the same vote]
\label{lem:forward-root}
In the setting of Assumption~\ref{ass:recovery-goldfish}, let a receiver
run its vote-time walk from a root \(R\) with \(A\preceq R\preceq Z\)
instead of \(A\), while \(R\) and \(B\) remain in its candidate tree and
the receiver's candidate-child set at every fork along the proposer's
walk is contained in the proposer's candidate-child set --- the
containment clause of Definition~\ref{def:usable-honest-round}.  The
walk returns \(B\).

In plain words, starting farther forward on the proposer's own path cannot
change the vote.
\end{lemma}

\begin{proof}
All merged vote-time views are equal by
Assumption~\ref{ass:recovery-goldfish}, so the vote weights at every fork
are equal at every honest validator, and the proposer's walk from \(A\)
is one fixed sequence of child choices whose visited path is the chain
segment from \(A\) through \(Z\); in particular it visits
\(R\).  The walk terminates at \(Z\), so \(Z\) has no candidate child
in the proposer's selection tree; \(B\), the fresh child of \(Z\), is
created at \(d_r\) and enters no selection tree.  A receiver's walk chooses within its candidate tree, but by the
counting rule of Definition~\ref{def:counting-rule} each child's
weight counts every known vote through the full known block tree, so
with equal merged vote views the weight of every child is the same
number at the proposer and at every receiver, whatever their candidate
trees.  The premise keeps the path to \(B\) inside the receiver's tree.
At each fork the proposer visited strictly before \(Z\), the chosen
child has the maximum weight among the proposer's candidate children and
wins the common deterministic tie-break; by the containment hypothesis
--- stated modulo the proposal path, whose only block off the segment
through \(Z\) is \(B\) itself --- the receiver's candidate-child set at
that fork is a subset of the
proposer's and contains the chosen child, and since
weights do not depend on the candidate tree, the same child wins there.
At \(Z\) the receiver's candidate-child set is contained in the
proposer's empty set united with the exempted path, hence is exactly
\(\{B\}\): \(B\) is present by the premise, and any other candidate
child of \(Z\) would lie in the proposer's selection tree, which has
none.  The receiver's walk therefore steps from \(Z\) to \(B\), and
\(B\), a first-slot block with no descendants at vote time, is the
output.  The walk is memoryless: its choices from any
visited block depend only on the view, the tree, and that block.  The
walk started at \(R\) therefore makes the same
choices from \(R\) onward and returns \(B\).
\end{proof}

\begin{definition}[Raw Goldfish vote]\label{def:recovery-goldfish-vote}
At \(d_r+\Delta\), an honest committee member merges the recognized
proposal's carried ordinary view into its own frozen slot view --- the
receipt buffer is not merged before the slot-view freeze --- runs
vote-time GHOST from its accepted stable root within the aged tree
\(\mathcal C^{-}_u(\Sigma_{u,\mathrm{vote}}^r)\), with the
proposal-path exemption of Definition~\ref{def:counting-rule}, and
votes for that output; the walk is total by
Lemma~\ref{lem:aged-walk-total}.  Without
a recognized proposal it runs the same rule on its frozen slot view and
vote-state candidate tree.  No SG head enters either merge, and no walk
recomputes \(\hmax\) below the walk's root.

In plain words, an accepted proposal supplies ordinary Goldfish data, not a
forced vote target.
\end{definition}

\begin{lemma}[An unobstructed honest proposal receives the exact honest vote batch]
\label{lem:honest-proposer-votes}
In an unobstructed honest-proposer round with proposal block \(B\), every
honest committee member votes for the exact block \(B\).

In plain words, an unobstructed honest proposal gets one exact vote batch.
\end{lemma}

\begin{proof}
The proposer chose \(Z_p^r\) by ordinary GHOST from \(A_p^r\) on the
proposal view and extended \(Z_p^r\) with \(B\).
Lemma~\ref{lem:honest-proposer-stable-root} puts every honest vote root on
the chain segment from \(A_p^r\) through \(Z_p^r\), and unobstructedness keeps
that root and \(B\) in the receiver's candidate tree.  Every honest voter
recognizes the proposal and applies the ordinary proposal-view merge.
Assumption~\ref{ass:recovery-goldfish} therefore gives every receiver the
proposer's exact merged view, and Lemma~\ref{lem:forward-root} makes the
receiver's vote-time GHOST return \(B\), even when its candidate tree
starts after \(A_p^r\).
\end{proof}


\section{SG outputs and ordinary finality-gadget signing}\label{sec:finality-votes}

This section defines the G0 gate, durable current-height choices, independent
latest-justification finality, and the one combined attestation used by the
state machine, and proves the two stable-root mechanism properties:
Lemma~\ref{lem:root-no-regression}, which drives the healing induction
of Section~\ref{sec:healing}, and Lemma~\ref{lem:root-backing}, the
accountability property consumed in Section~\ref{sec:refinement}.

\begin{definition}[Strong G0 gate and SG head]
\label{def:official-confirmation}
At action time \(a_r\), validator \(u\) completes the first slot's
available-confirmation evaluation of
Definition~\ref{def:tsq-confirmation}; let \(C_u^r\) be the deepest block
available-confirmed there, or \(\bot\) if none is.  Only the current
round's evaluation is read; no confirmation is carried from an earlier
round.  Validator \(u\) also computes the \emph{action head} \(G_u^r\): the
vote-time fork-choice walk of Definition~\ref{def:recovery-goldfish-vote},
run on the evaluation view \(V_u^+\), counting the evaluated first slot's
votes, from the round's vote-time stable root within the action state's
candidate tree.  This computation is guarded: if the round's fixed stable
root is absent from the action state's candidate tree --- for example
after a sibling maximum rise removed it from the viable band while its
grade stayed active --- the walk is skipped, \(G_u^r=\bot\), the
validator has no official confirmation at this action
(\(Q_u^r=\bot\) below), and it emits no SG head and no current-height
pair, the same abstention that a failed action-state check in
Definition~\ref{def:action-root} produces.  Every later output of this
definition is evaluated only under the guard.  The second slot's votes belong to the next evaluation and
are not counted here; this walk is a recovery-action computation, not the
ordinary next-slot Goldfish head.  Validator \(u\) also uses the grade~0
view fixed at \(d_r+2\Delta\).  A block \(Q\) is
\emph{veto-free} when no block conflicting with \(Q\) that descends from the
action state's finalized root has a grade~0 in \(X_u^2\) that is active
under Definition~\ref{def:active-grade} \emph{evaluated at the grade-0
freeze}: the activity filter counts exactly the finalized evidence the
validator had processed by \(d_r+2\Delta\), when \(X_u^2\) was fixed.
Evidence processed later does not silence a veto in the same round ---
its holder may not have seen it, so the veto stays conservative --- while
a grade for a block conflicting with finalized evidence processed by the
freeze, activated or not, casts no veto.  An
\emph{official confirmation} is a block \(Q\) that is
veto-free, conflicts with no finalized root whose evidence the validator
has processed by the action --- processed evidence blocks a confirmation
even in the round before it activates for root selection --- and satisfies
\[
 S_{u,\mathrm{act}}^r\preceq Q\preceq C_u^r,\qquad Q\preceq G_u^r.
\]
Available confirmation is downward closed, so such \(Q\) are themselves
available-confirmed.  Let \(Q_u^r\) be the deepest official confirmation,
or \(\bot\) if none exists.  When \(C_u^r\neq\bot\) and the walk behind
\(G_u^r\) runs from a root compatible with \(C_u^r\) with the path in its
candidate tree, Corollary~\ref{cor:tsq-walk} makes \(C_u^r\preceq G_u^r\),
so the head conjunct excludes nothing on the confirmed chain.

When \(Q_u^r\neq\bot\), the honest SG head is \(Q_u^r\), which is also the
confirmation gate for the current-height pair.  When \(Q_u^r=\bot\), let
\(R_u^r\) be the deepest veto-free block between
\(S_{u,\mathrm{act}}^r\) and the round's action root
\(R_{u,\mathrm{act}}^r\) that conflicts with no finalized root whose
evidence the validator has processed by the action.  The honest SG head is
\(R_u^r\), or empty if no such block exists; its current-height pair is
empty.
The finality pair for the already latest justification is computed
independently and is not suppressed by G0.

A validator whose fixed stable root fails the action-state check in
Definition~\ref{def:action-root}, or whose walk guard above fails, takes
the empty SG-head and current-height
branches directly; both cases fix every output of this action
(\(G_u^r=Q_u^r=\bot\), empty head, empty pair), so the completed action
is total.  The veto test applies the full filter of
Definition~\ref{def:active-grade}: a grade for a block conflicting with
the action state's finalized root, or with any processed finalized root,
is evidence only and casts no veto.

In plain words, G0 first chooses a veto-free official confirmation.  Without
one, the validator may still publish a veto-free prefix of the fixed stable
root, but it cannot sign a current-height choice.  Finalized-conflicting
grades do not veto either branch.
\end{definition}

\begin{remark}[The head-production rule meets the SG-head checklist]
\label{rem:head-checklist}
Every head this rule emits satisfies the two normative conditions of
Definition~\ref{def:sg-head}: a confirmation head satisfies
\(Q\preceq G_u^r\) with the action head in the action-state candidate
tree, and a fallback head lies between \(S_{i,\mathrm{act}}^r\) and
the admitted action root, both in that tree, so prefix closure and the
ancestor-closed viability filter give each honest head
\(S_{i,\mathrm{act}}^r\preceq B\) and viability.
\end{remark}

\begin{lemma}[Graded roots never conflict with a block supported by every honest head]
\label{lem:root-no-regression}
Suppose every honest nonempty SG head of the preceding round lies in
\(B\)'s subtree.  Then every graded root adopted in the current round ---
a strict grade-3 lower root, a proposer root accepted under the first
disjunct of item~1 of Definition~\ref{def:stable-root}, or an action
root retained from a graded vote-time root under its raw grade-1 check
--- is compatible with \(B\)\@.  Under the second disjunct of item~1
the adopted root is the receiver's grade-3 lower root, covered by the
first class when it is a strict graded block and exempt as an ungraded
selection otherwise.  If in
addition every honest validator emitted a nonempty head and \(B\)
remains in the adopting validator's \emph{aged} selection- and
vote-state candidate trees
\(\mathcal C^{-}_u(\Sigma_{u,\mathrm{sel}}^r)\) and
\(\mathcal C^{-}_u(\Sigma_{u,\mathrm{vote}}^r)\) --- the trees the
grade-3 selection itself uses --- every adopted
graded root is \(B\) or a descendant of \(B\).  Ungraded roots follow the finality
selection instead and are not constrained by this lemma; in particular,
an action root retained from an ungraded fallback root is exempt even
though its retention check uses raw grade~1 --- it may lawfully sit
behind \(B\) when the selection rule switches from a justification to
the finalized root.

In plain words, once the batch stabilizes on a block, grading cannot move
any adopted root onto a conflicting branch; with full honest
participation, and while the block stays in the adopting validator's
aged candidate trees, it cannot move an adopted root behind the block
either.
\end{lemma}

\begin{proof}
A block conflicting with \(B\) receives favorable support only from
faulty identities, of weight at most \(b<m\) by
Assumption~\ref{ass:sg-fault-bound}, because every honest
nonempty head lies in \(B\)'s subtree and honest identities do not
equivocate; so no conflicting block holds any grade, and every adopted
graded root is compatible with \(B\).  Under full honest participation,
honest weight \(W-b\geq m\) --- from \(2b<W\) of
Assumption~\ref{ass:sg-fault-bound}, since \(b\leq\lceil W/2\rceil-1\)
gives \(W-b\geq\lfloor W/2\rfloor+1=m\) ---
gives \(B\) direct support at least \(m\)
in every grade view, every ancestor-closed grade predicate holds for
\(B\), and the deepest graded block in a candidate tree containing
\(B\) is then \(B\) or a descendant of \(B\) by Lemma~\ref{lem:direct-roots-chain}.
This covers the first selected root class: a grade-3 lower root is the
deepest grade-3 block of the adopting validator's selection tree that
remains in its vote tree, and the hypothesis puts \(B\), with its direct
support \(m\), in both, so that deepest block is \(B\) or a descendant.
Under the first disjunct of item~1, every root of the triple is an
ancestor of \(A_p^r\), so the adopted root is \(A_p^r\) itself ---
grade-backed through the receiver's raw grade-1 check --- and
\(A_p^r\succeq L_u^r\succeq B\) by the first-disjunct guard and the
previous sentence.  Under the second disjunct the adopted root is the
receiver's grade-3 lower root: a strict graded block is covered by the
first class, and the ungraded fallback follows the finality selection
and is exempt by the statement.  The third class reduces to them by its own
definition: an action root retained from a graded vote-time root is that
vote-time root, which is a grade-3
lower root or an accepted grade-2 proposer root, already covered; the
retention of an ungraded fallback root is exempt by the statement.
\end{proof}

\begin{lemma}[Every adopted graded root has an honest gated supporter]
\label{lem:root-backing}
Under Assumption~\ref{ass:sg-fault-bound}, every graded root adopted by
an honest validator has favorable support of
weight at least \(m>b\) in the corresponding batch, and therefore an
honest supporter whose signed SG head contains it.  That head passed the
strong G0 gate of Definition~\ref{def:official-confirmation} when it was
signed.

In plain words, an adopted graded root is always backed by an honest head that itself
survived the veto; the head need not be an official confirmation.
\end{lemma}

\begin{proof}
Adoption requires the root's grade, whose threshold is \(m\); since
\(m>b\) under Assumption~\ref{ass:sg-fault-bound} alone ---
\(2b<W\) gives \(b\leq\lceil W/2\rceil-1<m\) --- some supporter
is honest.
Favorable support counts two classes --- identities whose eligible head
contains the root, and identities with an SG equivocation --- and an
honest identity never SG-equivocates: each recovery round schedules one
SG signing action per identity, and the store's single-signer rule
(Definition~\ref{def:stored-data}) has one active replica release at
most one signed result for it, so one honest identity emits at most one
SG head per round.  The honest supporter therefore lies in
the first class: its signed head contains the root.  An
honest SG head is emitted only through
Definition~\ref{def:official-confirmation}: it is the deepest official
confirmation when one exists and the veto-free fallback otherwise, so in
either case it passed the strong G0 gate.
\end{proof}

\begin{definition}[Current-height context]\label{def:ordinary-current-target}
The \emph{source proposal} for action \(r\) is the unique distinguished
proposal from round \(r-1\), whose SG batch this action grades.  If the
local round-\((r-1)\) record has no unique distinguished proposal --- the
proposer equivocated, or no proposal was recognized --- then action \(r\)
has no source proposal at that validator.  In that case the validator
skips every source-state option: it uses \(Y_i=Q_i^r\) and
\(\sigma_i=\bar\sigma_i=\sigma_a[Q_i^r]\) directly when its deepest
official confirmation \(Q_i^r\) is nonempty, with the same
\(k_i,\nu_i,T_i\) assignments as below, and emits no current-height
pair when \(Q_i^r\) is empty.  When local winners differ or a winner
equivocated, recognition may differ between honest validators, so this
branch is per-validator; the common-context results below therefore
always require universal recognition, stated as a premise by
Definition~\ref{def:usable-honest-round} and discharged at
honest-round-proposer rounds by Lemma~\ref{lem:common-recognition}.
Each proposal therefore supplies the context of exactly one action.
Suppose the source proposal is \(B\).  The action signs the
authenticated state of its deepest official confirmation \(Q_i^r\); no
separate symbol is needed for this block.  When
\(Q_i^r\) is empty, the validator emits no current-height pair by
Definition~\ref{def:height-vote-rule} below, no context block is signed, and
none of the quantities defined below is read.  Let
\(\bar\sigma_i=\sigma_a[Q_i^r]\) be the finality action state of
\(Q_i^r\) for this action's slot
(Definition~\ref{def:finality-action-state}): the authenticated state
with every pre-action slot closed.  The closure changes no height or
participation field (Lemma~\ref{lem:empty-slot-noop}); it can only
materialize a pending target name, which the fallback below would
otherwise supply, and reading it keeps one state name for every signed
field of the action.  If
\[
 B\preceq Q_i^r
 \quad\text{and}\quad
 \sigma_a[B].h=\bar\sigma_i.h,
\]
validator \(i\) uses \(Y_i=B\) and \(\sigma_i=\sigma_a[B]\), the
finality action state of \(B\) for the same slot.  Otherwise it
uses \(Y_i=Q_i^r\) and \(\sigma_i=\bar\sigma_i\).  Define
the action's current-height context by
\[
  k_i=\sigma_i.h,\qquad \nu_i=\sigma_i.\nj,
\]
and normally set \(T_i=\sigma_i.T_{k_i}\).  If that target field is empty
and \(Y_i.\slot\geq\sigma_i.s_{k_i}\), set \(T_i=Y_i\); otherwise set
\(T_i=\bot\).

Thus, when the source proposal is officially confirmed and its height is
still current, its
authenticated post-state fixes the complete triple
\[
  \operatorname{ctx}(B)=(k_B,T_B,\nu_B)
\]
for that action.  A proposal is carried for this purpose through exactly the
action that grades its SG batch, and no farther.

In plain words, validators that confirm the round proposal use that proposal's
height, target, and nonjustifiable flag even when their later tips differ.
\end{definition}

\begin{lemma}[The same-slot target materializes before inclusion]
\label{lem:same-slot-target}
If Definition~\ref{def:ordinary-current-target} returns the first block \(B\)
of height \(k_B\) while \(T_{k_B}=\bot\), every descendant that can include
the vote has already run the slot transition that records \(T_{k_B}=B\),
unless the chain
has advanced and the vote is stale.

In plain words, the same-slot rule never makes a block vote for its own
unknown root.
\end{lemma}

\begin{proof}
The vote cannot be included in \(B\).  The height-\(k_B\) transition
was consumed at a block (Lemma~\ref{lem:empty-slot-noop}), and the
following slot's \textsc{process\_slot} call records that transition
block as the target.  A prior height advance instead resets the
target and makes the height-\(k_B\) pair stale.
\end{proof}

\begin{lemma}[An official confirmation crosses no live stable root]
\label{lem:official-confirmation-root-safety}
In a recovery action, suppose one honest validator officially confirms \(Q\)
and every honest action-state Simplex root is compatible with \(Q\).  Then
every honest action root in that action is either compatible with \(Q\)
or conflicts with a finalized root whose evidence the confirming
validator has processed --- its activated root or a newer processed
outcome.  A root in the
second case is semantically inert in that validator's state and in every
honest state after the finalization evidence is relayed.

In plain words, G0 prevents a confirmation from crossing another honest
validator's live SG root; a root already excluded by finality needs no veto.
\end{lemma}

\begin{proof}
Write \(u\) for the confirming validator throughout, and fix another
honest validator \(v\) with action root \(R\).  The action-admission
guard in Definition~\ref{def:stable-root} gives two cases.
If \(R=S_{v,\mathrm{act}}^r\), it is compatible with \(Q\) by premise.
Otherwise that guard requires grade~1 for \(R\), regardless of whether
\(R\) originated as a proposed root, a grade-3 lower root, or an earlier
base root.  Lemma~\ref{lem:g1-g0} then gives the confirming validator
grade~0 for \(R\).  Let \(F_Q\) be that validator's action-state finalized
root.  \(F_Q\) is activated in the action state, so some block received
and accepted by \(a_{r-1}\) records it
(Definition~\ref{def:activation-filter}); the public spacing gives
\(a_{r-1}\leq d_r-2\Delta\), strictly before the grade-0 freeze
\(d_r+2\Delta\), so \(F_Q\) is among the finalized roots the confirming
validator had processed by the freeze.  If \(R\) conflicts with no
finalized root whose evidence the
confirming validator had processed by its grade-0 freeze, then in
particular \(R\) is compatible with \(F_Q\).  \(R\) is in the
confirming validator's accepted subtree, with its full ancestry,
strictly before the freeze, by Lemma~\ref{lem:g1-acceptance}.  That grade~0 is
therefore active in the veto test of
Definition~\ref{def:official-confirmation}; if in addition \(R\) descends from
\(F_Q\) and conflicts with \(Q\), the active grade
violates the veto-free guard in
Definition~\ref{def:official-confirmation} --- contradiction --- and a
strict ancestor of \(F_Q\) is compatible with \(Q\), because
Definition~\ref{def:finality-root} keeps \(F_Q\preceq S_{u,\mathrm{act}}^r\preceq Q\),
so in this case \(R\) is compatible with \(Q\).  Otherwise \(R\)
conflicts with some finalized root whose evidence the confirming
validator has processed --- the second case of the statement --- and
Definition~\ref{def:active-grade} makes it inert in the confirming state.
The confirming validator relays the finalization evidence under
Assumption~\ref{ass:recovery-network}; once another honest validator
processes it, Definition~\ref{def:active-grade} removes \(R\) from
grade-based use immediately --- the activation filter delays only root
selection --- and once that validator's activation filter admits the
outcome at its next round boundary, Lemma~\ref{lem:finalized-lockin} keeps
every later live tree block and fork-choice output in that finalized
root's subtree.  These
cases prove the claim.
\end{proof}

\begin{definition}[Current-height signing rule]\label{def:height-vote-rule}
Let \(C_i=Q_i^r\) be the action's confirmation gate --- its deepest
official confirmation --- and let
\((k,T,\nu)\) be Definition~\ref{def:ordinary-current-target}'s output for the
action, and let \(\mathcal H_i\) be durable history.  If
\(C_i=\bot\), the current-height pair is empty.  Otherwise choose in this
order:
\begin{enumerate}
 \item If \(\mathcal H_i.\tau(k)\) is true, repeat \((k,\bot)\) when the
   closed state satisfies \(\sigma_a[C_i].h\geq k\)
   (Definition~\ref{def:finality-action-state}), and otherwise emit empty.
 \item If \(\mathcal H_i.\mathrm{lock}(k)=T_L\neq\bot\), repeat
   \((k,T_L)\) when \(T_L\preceq C_i\), and otherwise emit empty.
 \item If \(\mathcal H_i.T(k)=T_0\neq\bot\), repeat \((k,T_0)\) when
   \(T_0\preceq C_i\).  If not, record and sign \((k,\bot)\) when the
   closed state satisfies \(\sigma_a[C_i].h\geq k\); otherwise emit empty.
 \item With no history and \(\nu=\mathsf{true}\), record and sign
   \((k,\bot)\) when \(\sigma_a[C_i].h\geq k\).
 \item With no history and \(\nu=\mathsf{false}\), record and sign \((k,T)\) when
   \(T\neq\bot\) and \(T\preceq C_i\).  Otherwise record and sign
   \((k,\bot)\) when \(\sigma_a[C_i].h\geq k\); if neither test
   passes, emit empty.
\end{enumerate}
Every durable write completes before signature release.  Timeout history has
precedence over target history.  A finality lock never switches to timeout.

In plain words, confirmed on-chain saved targets repeat, confirmed off-chain
unlocked targets become timeouts, and an unconfirmed height stays empty.
\end{definition}

\begin{definition}[Independent finality signing rule]
\label{def:finality-vote-rule}
The finality pair is read from one explicit state: the fork-choice
action state \(\Sigma_{i,\mathrm{act}}^r\) of
Definition~\ref{def:recovery-timing}.  Its fields \((J,h_j)\) and
\((F,h_F)\) are that state's activated greatest-justification and
greatest-finalized selections --- not the chain state of any context
block, whose recorded justification may lag the store.  The rule is
total: the fork-choice action state exists at every action, including
one whose SG head and current-height pair are empty.  Validator \(i\)
signs finality pair \((h_j,J)\)
exactly when
\[
 h_j>h_F,\qquad F\preceq J,
\]
it knows the justification certificate \(\mathsf{JC}(h_j,J)\) of
Definition~\ref{def:certificates},
\(\mathcal H_i.T(h_j)=J\),
\(\mathcal H_i.\tau(h_j)=\mathsf{false}\), and
\(\mathcal H_i.\mathrm{lock}(h_j)\in\{\bot,J\}\).
It records \(\mathcal H_i.\mathrm{lock}(h_j)=J\) before first release.  Otherwise
the finality pair is empty.  A newer justification is the only eligible
finality pair because the chain state keeps only the latest-\(J\) tally.

This rule is independent of the current-height confirmation gate.  A valid
finality pair may therefore accompany an empty current-height pair and an
empty SG head.

In plain words, finality piggybacks the latest saved justification even when
strong G0 vetoes new SG or height work.
\end{definition}

\begin{definition}[Ordinary combined attestation]\label{def:fg-rule}
At an ordinary SG/FG action, an honest validator first derives the SG head
and current-height confirmation gate, then evaluates the two pair rules
\emph{in order}: first the finality pair of
Definition~\ref{def:finality-vote-rule}, whose durable lock write
completes before the next step, then the current-height pair of
Definition~\ref{def:height-vote-rule}, whose history read sees that
write.  The ordering is load-bearing: a lock recorded in this action
already governs the current-height rule's lock case, so the two fields
of one attestation can never form E1 evidence.  It
signs the one attestation of Definition~\ref{def:fg-message}.  The attestation
contains no state-block root or confirmed-block root.

In plain words, one ordinary message carries three separately checked
fields, with G0 applying only to the first two.
\end{definition}

\begin{lemma}[Honest finality-gadget attestations are not slashable]
\label{lem:signer-safety}
An honest validator following
Definitions~\ref{def:height-vote-rule} and~\ref{def:finality-vote-rule}
never creates E1 or E2 evidence.

In plain words, durable history prevents every honest target, timeout, and
finality conflict.
\end{lemma}

\begin{proof}
Timeout history is checked first and prevents any later finality pair at that
height.  A first target is written once, and every later target repeats it,
so E2 is impossible.  An unlocked saved target may switch to timeout only
before a finality entry exists --- including one recorded earlier in the
same action: the ordering of Definition~\ref{def:fg-rule} evaluates the
finality rule first and makes its lock write visible to the
current-height rule's history read, so an action that signs a finality
pair at a height can no longer take the timeout branch at that height
and instead repeats the locked target or emits empty.  The timeout, once
recorded, permanently blocks finality at its height.
After a finality entry exists, the height rule repeats that exact target or
emits empty, so E1 is impossible.  Empty pairs trigger neither rule.
\end{proof}

\begin{definition}[Retransmission and pending references]
\label{def:earlier-votes}
While a height can affect honest fork choice, honest validators retransmit
every saved target or timeout still usable on a live branch, the blocks and
individual attestations needed to replay latest justified and finalized
states, every justification certificate retained by a finality lock, and
every live SG head and raw Goldfish vote needed by an uncompleted action.
Unknown-reference objects remain queued until their blocks and ancestry
arrive.  A signature changes consensus state only through ordinary inclusion
in a valid block.

An entry is discarded only after all honest validators have finalized a
common descendant and no competing branch at or below that height can affect
protocol choices.  Local height advance alone is not enough.

In plain words, old signatures remain available until neither liveness nor
accountability can need them.
\end{definition}

\section{Recovery action and liveness assumptions}\label{sec:rounds}

This section puts the preceding rules on one timeline and states the network,
proposer, inclusion, and SG/FG composition premises used by the healing
proof.  Figure~\ref{alg:recovery-action} summarizes one recovery round.

\begin{figure}[H]
\caption{One recovery round}\label{fig:recovery-round}
\label{alg:recovery-action}
\begin{center}
\begin{boxedminipage}{0.98\textwidth}
\scriptsize
\begin{algorithmic}[1]
\Function{recovery\_round}{$r$}
  \State \textbf{require} one history-compatible result for every preceding
    action
  \State at \(X^-\), snapshot eligible preceding SG heads and relay evidence
  \State at \(X^0\), compute G2 and relay every witness
  \If{\(i\)'s ordinary proposer credential wins the first slot}
    \State choose the stable root; run ordinary proposer-view merge and GHOST
    \State build, sign, and broadcast the valid proposal extending that parent
  \EndIf
  \State at \(X^1\), compute G3 and G1
  \State validate the unique timely proposal; use none if the vote-time view holds two signed round proposals
  \State merge only the accepted proposal's ordinary Goldfish view
  \State run ordinary vote-time GHOST and broadcast the raw Goldfish vote
  \State at \(X^2\), compute G0 and freeze the first slot's support view
  \State at \(d_r+4\Delta\), process the second slot's ordinary proposal
  \State at \(d_r+5\Delta\), run ordinary vote-time GHOST with the fixed root
    and broadcast the raw Goldfish vote
  \State at \(a_r\), process every delivery due, then evaluate the first slot's
    available confirmation
  \State derive the action state and read the available confirmation and head
  \State derive the SG head and current-height confirmation gate
  \State select the source proposal and current-height context by
    Definition~\ref{def:ordinary-current-target}
  \State derive the latest-\(J\) finality pair first and complete its
    durable lock write, per the ordering of Definition~\ref{def:fg-rule}
  \State then derive the current-height pair under the G0 gate, its
    history read seeing that write
  \State sign and broadcast the individual combined attestation of
    Definition~\ref{def:fg-rule}
  \State after the action, the rest of the second slot and any later slots
    continue ordinary Goldfish until the next recovery round
\EndFunction
\end{algorithmic}
\end{boxedminipage}
\end{center}
\end{figure}

\begin{definition}[Recovery window]\label{def:recovery-window}
A recovery window is a finite post-\(\GST\) interval
\([t_0,t_1]\) such that Assumptions~\ref{ass:recovery-network},
\ref{ass:recovery-goldfish}, and~\ref{ass:frontier-opportunities}
hold throughout \([t_0,t_1)\); the right endpoint \(t_1\) is the
end of the analyzed period or the first event that violates one of
them --- the \emph{terminating event}, such as the admission of a
released pre-window branch failing the entry condition.  The
terminating event belongs to the window: it is exactly what the
disruption outcome of the healing theorem names, while every earlier
point satisfies the assumptions.
The public recovery-round schedule continues throughout the interval,
independently of local fork choice.  The protocol does not switch to a
different Goldfish state machine at the end; this paper claims persistence
only for later actions to which the same assumptions apply.

In plain words, recovery is the interval analyzed by the proof, not a local
mode bit or a change to block validity.
\end{definition}

\begin{definition}[Relay point and normalized frontier]
\label{def:relay-frontier}
The \emph{relay point} is the preceding action cutoff \(a_0\) of the
first counted recovery round.  This sentence defines the term; every
later use, including Definition~\ref{def:activation-filter}'s, refers
here.  A recovery round or action is \emph{counted} when it lies at or
after the relay point, inside the analyzed window.  An opportunity or
round is \emph{charged} when it is one of the budget items enumerated
in the statement of Proposition~\ref{prop:opportunity-budget}, a
separate notion from countedness.  At the relay point, the
\emph{admitted frontier} is the greatest reachable state-height of the
admitted store, and \(h_F^0\) is the greatest finalized height
recorded on any admitted branch there --- branch states record
finalized pairs of the one safe finalized chain, so this maximum is
well defined, and every admitted branch's recorded finalized height is
at most \(h_F^0\), so a debt test against \(h_F^0\) implies the
same test against each branch's own recorded height.  A valid branch
is \emph{relevant} when some honest validator admits it to its store
during the window; a branch never admitted appears in no honest state
and affects no honest selection.  The guarded height \(H\) is the
\emph{least} height at or above the admitted frontier with
\(K\mid H\) and \(H-h_F^0>D_{\rm debt}\); by the maximality of
\(h_F^0\), every admitted branch entering \(H\) satisfies the same
entry condition with its own recorded finalized height, so \(H\) is
nonjustifiable on every such branch.  The height below \(H\) is
ordinary on every branch: \(K\mid H\) and \(K\geq2\) give
\(K\nmid H-1\), and the nonjustifiable latch requires
\(K\mid h\), so \(\nj=\mathsf{false}\) at \(H-1\) regardless
of any branch's recorded debt.

In plain words, the relay point opens the analyzed window, and \(H\)
is the first periodic height past the frontier with enough room below
it.
\end{definition}

\begin{assumption}[Recovery opportunities]\label{ass:frontier-opportunities}
This assumption has five parts.
\begin{enumerate}
\item \emph{Action histories.}
For every honest actor, each counted recovery action has one
history-compatible result for every required predecessor action.  A merge of
replicas with two divergent completed histories remains safe under
Definition~\ref{def:store-join}, but it is outside this liveness assumption
because its successor actions abstain.

\item \emph{Normalized frontier.}
By the relay point of Definition~\ref{def:relay-frontier}, after
\(\GST\), every piece of pre-recovery evidence held by an honest
validator --- a finite set --- has been relayed and delivered under
the retransmission rule of Definition~\ref{def:earlier-votes}.  No
honest validator signed a nonempty height pair above \(H\) before
that point.  Honest validators may already have saved targets,
timeouts, or locks at \(H\), and a certificate may already have
advanced one branch from \(H-1\) to \(H\).  The height filter
places every honest selected branch at height \(H-1\) or \(H\).
A certificate on one branch does not replay its transition on a
sibling.  Finality advances completing during the window are covered
by the debt dichotomy of Lemma~\ref{lem:debt-dichotomy} below.  The
formal results of this paper are stated for this normalized frontier:
the admitted maximum is \(H\) and every honest selected branch is at
\(H-1\) or \(H\), as required above.  A recovery from a frontier
further behind repeats the same mechanisms height by height, but its
lemmas and opportunity count are not restated here.  An object
admitted by one honest validator at a cutoff is relayed and is
available to every honest validator by the next cutoff.  The proof
does not infer that two stores at the same cutoff contain the same
adversarially released objects.

\item \emph{Proposer recurrence.}
Within the finite security horizon considered here, with probability at
least \(1-\operatorname{negl}(\lambda)\), where \(\lambda\) is the
cryptographic security parameter of Assumption~\ref{ass:crypto}, every
interval of \(\rho\)
consecutive recovery rounds after that relay point contains an honest round
proposer in the common-recognition sense of
Definition~\ref{def:recovery-timing}: one honest validator's proposal
is the local selection winner at every honest validator.  This is the
source's honest-leader recurrence event restated for the abstracted
proposer selection of Definition~\ref{def:recovery-proposal}; because
proposer selection is outside this paper's scope, the event is assumed
at this strength rather than derived from a lottery analysis, and a
round whose honest eligible proposal loses one receiver's local
selection to a competing released proposal does not witness it.  (The
symbol \(\kappa\) is reserved for Goldfish's confirmation depth and
names only the source's rule.)  The window length \(\rho\) is a protocol constant.  This event is the proposer premise
used for SG and available-chain healing; the later finality restart has its
own clean target premise.

\item \emph{Inclusion and head persistence.}  This part has five
   sub-clauses.
   (a)~\emph{Delivery and inclusion behavior:} by
   Assumption~\ref{ass:recovery-network}, every honest SG/FG
   attestation reaches every honest proposer strictly before one
   delivery bound has elapsed, and an honest round proposer includes
   every pending valid honest attestation whose signed height and
   ancestry checks pass on its proposal branch.
   (b)~\emph{Assumed includer:} if the closing action for a proposal
   \(B\) releases honest current-height pairs of weight at least
   \(q\), then within \(\rho\) rounds an honest round proposer
   proposes a descendant of \(B\) whose prestate is still at that
   height and includes those pairs, unless the branch has already
   advanced or the permitted heights ---
   Definition~\ref{def:permitted-heights} below --- have moved, in
   which case no such inclusion is needed.  The \emph{occurrence} of
   this including proposal at the pre-confirmation heights is assumed,
   not derived: before the first recovery confirmation no induction
   keeps honest roots inside \(B\)'s subtree across the round between
   the pair release and the includer; after the first confirmation the
   healing induction derives it, and this clause is not used.
   (c)~\emph{Head persistence:} from a charged unobstructed proposal's
   closing action until every honest validator has processed the
   includer's advance --- \emph{common} always means exactly this,
   processed by every honest validator, which relay on receipt delivers
   within one bound of the first honest processing --- honest SG heads
   remain in the confirmed block's subtree, and once the advance is
   common the first later charged honest proposal descends from the
   advanced block, so no third honest-proposer round at the same height
   is a fresh opportunity outside the budget list.  This
   pre-confirmation head persistence is disclosed assumption residue;
   the corresponding post-confirmation property is derived by the
   healing induction.
   (d)~\emph{Finality-pair inclusion:} the same rule includes pending
   finality pairs only on a branch whose prestate has processed, and
   still names, their justification as its latest unfinalized one.
   (e)~\emph{Includer delivery:} the including block, its direct
   height check, and its post-state reach every honest validator before
   the next counted SG/FG action.

\item \emph{Round shape.}
A recovery round contains at least two slots, so its action follows the
first slot's evaluation, by Definition~\ref{def:recovery-timing}.
\end{enumerate}

In plain words, the standard every-\(\rho\)-round event supplies honest
round proposers; each honest actor has one actual action history, and
synchrony and the grade rules supply the remaining liveness argument.
\end{assumption}




\section{Reaching the first recovery confirmation}
\label{sec:first-confirmation}

This section defines the first-confirmation event and proves that the
honest-proposer recurrence of
Assumption~\ref{ass:frontier-opportunities} forces it: it bounds
pre-confirmation height movement, classifies every honest-proposer
round, and derives the closing-action confirmation.

\paragraph{Proof map.}
Lemma~\ref{lem:selectable-roots} and
Lemma~\ref{lem:pre-confirmation-cap} bound pre-confirmation roots and
maxima; Lemma~\ref{lem:common-recognition} and
Lemma~\ref{lem:honest-proposer-coverage} classify every
honest-proposer round into an opportunity or one of three bounded
obstructions; Lemma~\ref{lem:closing-action-existence} and
Lemma~\ref{lem:honest-proposal-confirms} --- through the shared
Lemma~\ref{lem:gate-closure} and
Lemma~\ref{lem:pre-confirmation-finality-compatible} --- turn an
unobstructed proposal into universal official confirmations.
Section~\ref{sec:healing} then proves catch-up and the height
advances, roots every later action at the first confirmation through
the drain lemmas and the persistence induction, and assembles
Proposition~\ref{prop:opportunity-budget} and
Theorem~\ref{thm:eventual-healing} from these pieces.

\begin{definition}[First recovery confirmation]
\label{def:first-recovery-confirmation}
Let \(r_\star\) be the first recovery action in which an honest validator has
an official confirmation \(B_\star\) whose state-height is at least \(H+1\).
If several honest confirmations first occur in that action, choose any one
as \(B_\star\); the definition does not assume that the other confirmations
are compatible with it.  Write \(C_\star\) for the confirming validator's
deepest available-confirmed block at \(r_\star\)'s evaluation;
\(B_\star\preceq C_\star\) by Definition~\ref{def:official-confirmation}.
The names \(r_\star\), \(B_\star\), and \(C_\star\) are used only
in the proof.  Validators do not test whether an action or confirmation is
first.  The \emph{drain action}, denoted \(r_{\rm drain}\), is the closing
action that grades the SG heads signed in \(r_\star\).

In plain words, the safety proof starts from the first action whose evaluation
gives some honest validator an official confirmation that has reached the
recovery height.
\end{definition}

\begin{lemma}[Every honest height vote has a confirmed witness]
\label{lem:height-vote-confirmation}
If an honest validator signs a nonempty current-height pair at a height
\(x\geq H+1\) --- necessarily at a counted recovery action, by
Assumption~\ref{ass:frontier-opportunities} ---
its official confirmation \(Q\) is nonempty and
\(\sigma[Q].h=\sigma_a[Q].h\geq x\) --- the closed state of the signing
action and \(Q\)'s own post-state agree on the height by
Lemma~\ref{lem:empty-slot-noop}, so the confirmed block itself has
state-height at least \(x\).

In plain words, an honest target or timeout at height \(x\) proves that its
signer first confirmed a chain that had reached \(x\) by its action.
\end{lemma}

\begin{proof}
Definition~\ref{def:official-confirmation} makes the current-height pair empty
when the official confirmation is empty.  For a timeout, every applicable
case of Definition~\ref{def:height-vote-rule} checks
\(\sigma_a[Q].h\geq x\) in the closed state.  For a target, the rule checks that the height-\(x\)
target lies on the chain ending at \(Q\).  A stored or repeated target is the
unique first block of that height by
Lemma~\ref{lem:chain-target-uniqueness}; a same-slot target has the same
height by Definition~\ref{def:ordinary-current-target} and
Lemma~\ref{lem:same-slot-target}.  Thus the target's authenticated post-state
has height \(x\), so the descendant \(Q\) has block post-state height at
least \(x\), and \(\sigma_a[Q].h=\sigma[Q].h\) by
Lemma~\ref{lem:empty-slot-noop}.
These cases exhaust every nonempty current-height pair.
\end{proof}

\begin{lemma}[Every height increase has an honest confirmed witness]
\label{lem:height-increase-confirmation}
Whenever a valid chain advances from a height \(x\geq H+1\) to
\(x+1\), some honest signer in the advancing quorum has an official
confirmation whose chain has reached height \(x\).  The height
restriction makes every honest quorum pair post-relay:
Assumption~\ref{ass:frontier-opportunities} says no honest validator
signed a nonempty pair above \(H\) before the relay point, so an honest
pair at height \(x\geq H+1\) was signed at a counted recovery action.

In plain words, neither a justification nor a timeout can move the height
ahead of every honest confirmed chain.
\end{lemma}

\begin{proof}
Lemma~\ref{lem:height-progression} gives a target or progress quorum of weight
at least \(q\).  Since \(q>b\) by
Lemma~\ref{lem:integer-thresholds}, the quorum contains an honest signer.
Lemma~\ref{lem:height-vote-confirmation} supplies that signer's confirmed
witness at height \(x\).
\end{proof}

\begin{lemma}[The store maximum is capped until a first recovery confirmation]
\label{lem:pre-confirmation-cap}
As long as no honest validator has an official confirmation whose chain has
state-height at least \(H+1\), no honest store admits a branch of
state-height above \(H+1\), and every store-global maximum is at most
\(H+1\).

In plain words, the cap does not presuppose that a first confirmation ever
happens; it holds exactly until one does.
\end{lemma}

\begin{proof}
A branch of state-height \(H+2\) contains a valid height increase from
\(H+1\) to \(H+2\).  Lemma~\ref{lem:height-increase-confirmation} gives
an honest signer in its advancing quorum with an official confirmation
whose chain has reached height \(H+1\), contradicting the premise.
\end{proof}

\begin{lemma}[A finality advance preserves nonjustifiability or restarts finality]
\label{lem:debt-dichotomy}
Let a finality advance activate finalized height \(h'>h_F\) at any time
during the window.  Exactly one
of the following holds.
\begin{enumerate}
 \item \(H-h'>D_{\rm debt}\).  Then every branch entering height \(H\)
   while every finalized height yet activated satisfies the same debt
   condition --- that is, before any later case-2 advance --- latches
   \(\nj=\mathsf{true}\), and every
   result stated under ``height \(H\) is nonjustifiable'' continues to
   hold over that span.
 \item \(H-h'\leq D_{\rm debt}\).  Two subcases, by where the
   finality transition itself completed.  If it completed inside the
   window, the advance is a fresh post-\(\GST\) finalization within
   \(D_{\rm debt}\) of the recovery frontier: the finality-restart
   outcome itself, which the healing theorem reports directly.  If
   instead the transition completed before the window on a branch whose
   evidence was withheld and first released during the window ---
   activation of history, not new finality --- the guarded height's debt
   test fails against \(h'\), the premises of this window's analysis at
   \(H\) are withdrawn, and the window ends as a finite disruption.
   Recovery is then repeated: the guarded height re-targets the least
   \(H''\geq h'\) with \(K\mid H''\) and \(H''-h'>D_{\rm debt}\), where
   \(H''-h'\leq D_{\rm debt}+K\), and a subsequent recovery window under
   re-instantiated assumptions runs at the then-admitted frontier, as
   the repetition corollary after the healing theorem formalizes.  Any
   official confirmation already made remains safe by
   Lemma~\ref{lem:tsq-uniqueness}, which is per-slot and unaffected.
   Each advance triggers at most one such re-targeting.
\end{enumerate}
Moreover, the first case-2 advance activating before the first
recovery confirmation finalizes at height at most
\(H-1\); the claim says nothing about later advances.

In plain words, an advance far below the frontier never revives the
guarded height; an advance near the frontier either is already the
healing goal or moves the guarded height one period up.
\end{lemma}

\begin{proof}
Definition~\ref{def:nonjustifiable} latches
\(\nj=(K\mid H)\wedge(H-h_F>D_{\rm debt})\) at entry, from the entering
transition's post-finality state, and the latch is immutable until the
next height transition.  \(K\mid H\) holds by the choice of \(H\).
Activated finalized heights are monotone along one chain by
Lemma~\ref{lem:finalized-chain}, so in case~1 every finalized height a
branch can record at entry --- the pre-window value or any advanced value
up to \(h'\) --- satisfies the debt inequality, every entrant latches
\(\nj=\mathsf{true}\), and branches that entered earlier keep their
latch.  For case~2, the first subcase restates the advance itself; the
second subcase makes no claim beyond the arithmetic of the re-target ---
\(H\) fails the debt condition against \(h'\), the least qualifying
\(H''\) satisfies \(H''\leq h'+D_{\rm debt}+K\), and re-targetings are
keyed to the monotone activated heights, at most one per advance ---
and the disruption framing, which the repetition corollary supplies.
For the final claim, a
finalization at height \(H\) requires a height-\(H\) justification;
before any case-2 advance, case~1 held for every earlier advance, so
every branch at \(H\) has \(\nj=\mathsf{true}\) and
Lemma~\ref{lem:nonjustifiable-fg} excludes that justification.  A
finalization at height \(h'\geq H+1\) activating before the first
recovery confirmation is impossible: its justification's target quorum
needs honest weight \(q-b\) in
height-\(h'\) pairs, and
Lemma~\ref{lem:height-vote-confirmation} would give such a signer an
official confirmation of state-height at least \(H+1\), a
contradiction.  Hence the
first pre-confirmation case-2 advance finalizes below \(H\); the claim
says nothing about advances after the confirmation, which the healing
induction and the disruption subcase govern.
\end{proof}

Until the healing theorem, the analysis proceeds under case~1 of
Lemma~\ref{lem:debt-dichotomy} throughout the window: every finality
advance except the confirmed chain's own above-\(H\) finalizations
keeps \(H-h'>D_{\rm debt}\), so height
\(H\) stays nonjustifiable on every relevant branch, before and after
the first recovery confirmation.  The theorem
states the two subcases of case~2 as alternative outcomes: a restart, or
a rerun at the recomputed height.

\begin{lemma}[No honest state ever holds a height-\(H\) justification]
\label{lem:no-height-h-justification}
Under case~1 of Lemma~\ref{lem:debt-dichotomy} throughout the window, no
justification at height \(H\) appears in any honest state during the
window, before or after the first recovery confirmation.

In plain words, the guarded height stays unjustifiable for every branch an
honest validator ever uses, so no later argument needs the pre-window
hypothesis again.
\end{lemma}

\begin{proof}
A justification enters an honest state only on a branch the validator has
admitted, and an admitted branch is relevant by definition.  A branch
fixes its nonjustifiable latch at its height-\(H\) entry from its own
on-chain post-finality state (Definition~\ref{def:nonjustifiable}), and
every descendant still at height \(H\) inherits the latch.  A relevant
branch entering \(H\) satisfies the debt condition at entry: for a branch
entering before the window this is
Assumption~\ref{ass:frontier-opportunities}; a branch entering during
the window with on-chain finalized height \(h_F\geq H-D_{\rm debt}\)
carries an included finality certificate at that height.  That height
is below \(H\): a chain's recorded finalized height satisfies
\(h_F\leq h_j<h\) (Lemma~\ref{lem:finalized-before-justified}), so a
branch entering \(H\) records \(h_F\leq H-1\), and the carved-out
finalizations above \(H\) cannot be the recorded one.  The honest
validator that admits the branch processes that certificate, activating
a finalized height \(h'\) with
\(H-D_{\rm debt}\leq h'\leq H-1\) --- an advance the hypothesis
excludes.  So every relevant branch at height \(H\) has
\(\nj=\mathsf{true}\), Definition~\ref{def:height-outcome} disables its
target branch, and Lemma~\ref{lem:nonjustifiable-fg} applies: no
justification at \(H\) forms on any relevant branch.  A justification
formed on a branch no honest validator ever admits appears in no honest
state.
\end{proof}

\begin{lemma}[Selectable roots before the first confirmation]
\label{lem:selectable-roots}
Before the first recovery confirmation, every justification in any
honest store has height at most \(H-1\), and every finalized root has
height below \(H\).  Consequently every processed finalized root whose
height is below a valid block's state-height lies on that block's chain,
by Lemma~\ref{lem:past-finalized}, and every finality-derived Simplex
root at an honest validator is a justification of height at most
\(H-1\) or a finalized root of height below \(H\).

In plain words, until the first confirmation no root at or above the
disabled height exists, so every selectable root sits below the recovery
heights.
\end{lemma}

\begin{proof}
Height \(H\) can be neither justified nor finalized, by
Lemma~\ref{lem:no-height-h-justification}.  A justification at height
at least \(H+1\) has a target quorum of weight \(q>b\), hence an
honest signer; that signer's official confirmation has state-height at
least \(H+1\) by Lemma~\ref{lem:height-vote-confirmation} --- an
honest confirmed witness at the recovery height, so the first recovery
confirmation has occurred, a contradiction.  A finalized root at height
at least \(H\) copies a justification at that height, excluded by the
same two cases.  The Simplex root of Definition~\ref{def:finality-root}
is the selected justification or the finalized root.
\end{proof}


\begin{definition}[Permitted heights]\label{def:permitted-heights}
The \emph{permitted heights} of a store are \(\hmax-1\) and \(\hmax\):
the state-heights of the leaves the viability filter of
Definition~\ref{def:finality-root} admits.  A store-maximum rise moves
both.

In plain words, the fork choice always keeps the two highest heights,
and raising the maximum raises the pair.
\end{definition}

\begin{lemma}[Every honest round proposer is a recovery opportunity]
\label{lem:honest-proposer-coverage}
After the relay point in
Assumption~\ref{ass:frontier-opportunities}, consider an honest round
proposer --- the common-recognition event of
Definition~\ref{def:recovery-timing} --- before the first recovery
confirmation.  One of the following happens:
\begin{enumerate}
 \item its proposal is unobstructed in the sense of
   Definition~\ref{def:usable-honest-round};
 \item the store-global maximum rises enough to move the next proposal to
   higher permitted heights (Definition~\ref{def:permitted-heights}); or
 \item the unique height-\(H-1\) justification is first admitted and
   redirects the Simplex root.  This case can delay at most one honest
   proposal; or
 \item a finalized height is first processed at the round's cutoffs ---
   by a transition completing inside the window or by the released
   evidence of a pre-window transition --- and either its
   activation is staggered across the round boundary, or the newly
   processed evidence deactivates a graded candidate root mid-round by
   Definition~\ref{def:active-grade} before any activation.  Such rounds
   number at most two per distinct finalized height above \(h_F^0\)
   first processed before the first recovery confirmation --- one for
   the mid-round deactivation, one for the staggered activation ---
   each height below \(H\), so at most twice the finalized-height
   budget; a first-processed height at or below \(h_F^0\) obstructs
   no round, its root lying below an already-recorded one on the safe
   finalized chain.
\end{enumerate}
In plain words, proposer fairness supplies actual protocol opportunities:
there is no separate unquantified assumption that a chosen honest round
is unobstructed beyond the cases listed.
\end{lemma}

\begin{proof}
We check the selection, vote, and action boundaries in order.

At selection time, all pre-recovery evidence known to an honest validator has
crossed the relay point by
Assumption~\ref{ass:frontier-opportunities}.  Accepted finalized roots form
one chain by Lemma~\ref{lem:finalized-chain}.  Before the first recovery
confirmation, Lemma~\ref{lem:no-height-h-justification} excludes a height-\(H\)
justification.  The only justification that can newly become the selected
Simplex root on a conflicting branch while the maximum is \(H\) is therefore
the unique height-\(H-1\) justification.  If it is first admitted during the
round, item~3 holds.  Once its certificate and carrying state have been
relayed, Lemma~\ref{lem:target-uniqueness} prevents another such redirection.

Suppose no such redirection occurs.  If the proposer has a grade-2 root
\(A_p^r\), every receiver's strict grade-3 lower root \(L_u^r\) ---
strict at that receiver --- reaches grade~2 in the
proposer's view by Lemma~\ref{lem:g3-g2}.  Direct-majority compatibility
(Lemma~\ref{lem:direct-roots-chain}) and the proposer's deepest-choice rule
then give \(L_u^r\preceq A_p^r\) whenever both roots remain live in the same
candidate tree.  Lemma~\ref{lem:g2-g1} gives every honest receiver grade~1
for \(A_p^r\).  If a newly learned finalized root removes a former lower root
from that tree, Definition~\ref{def:stable-root} drops the inert grade and
uses the newer Simplex root instead.

If instead the proposer's and a receiver's candidate trees disagree on
which of those roots is live, their store-global maxima differ at the
cutoff.  Every witness counted by an honest receiver's aged test of
Definition~\ref{def:recovery-context} reached the proposer strictly
before its selection cutoff, so a receiver-live lower root is live in the
proposer's tree whenever their maxima, active finalized roots, and
processed finalized-evidence sets agree --- the last because
Definition~\ref{def:active-grade} filters grades by processed evidence,
so equal evidence sets veto the same roots.
A guard failure therefore forces an \(\hmax\) difference, a staggered
finality activation, or a divergence in processed finality evidence.  In
either finality case a finalized height is first processed at the
round's cutoffs, whether its transition completed inside the window or
on a released pre-window branch:
the evidence is relayed on receipt, so it is processed by every honest
validator within one delivery bound and active at every honest validator
by the following round boundary --- the one-round staggering bound of
Remark~\ref{rem:activation-stagger}.  Such rounds are case~4 of the
statement; each charges the finalized height above \(h_F^0\) first
processed at it, at
most two rounds per height as the statement counts, bounded before the
first recovery confirmation by twice the
finalized-height budget below \(H\).  A maximum divergence
means a maximum rise is in flight; its witness, held by one honest
validator at its cutoff, is in every honest store strictly before the same
cutoff of the following round, whose selection therefore uses the higher
common maximum.  Such a round is case~2 of the statement, and
Lemma~\ref{lem:filtered-hmax-monotone} and
Lemma~\ref{lem:pre-confirmation-cap} bound the number of maximum rises
before any first recovery confirmation, so this charges case~2 at most
once per rise and the four cases remain exhaustive.

We now place each newer Simplex root.  Let \(R\) be a finalized root whose
height is below the proposal block's state-height \(k\).
Lemma~\ref{lem:past-finalized} gives \(R\preceq B_p^r\).  Since \(B_p^r\) is
the fresh child of \(Z_p^r\), the strict height inequality gives
\(R\preceq Z_p^r\).  A finalized root at height at least \(k\) either exposes
the unique height-\(H-1\) justification in item~3, or has height at
least \(H\), excluded before the first confirmation by
Lemma~\ref{lem:selectable-roots}.  The
Simplex-root rule likewise leaves only the unique height-\(H-1\)
justification as a possible off-path justification root.

No Simplex root strictly descends from \(B_p^r\) before the first recovery
confirmation.  Such a root is a justification or finalized root whose
target strictly descends from \(B_p^r\).  By chain-target uniqueness
(Lemma~\ref{lem:chain-target-uniqueness}), the height-\(k\) target of
\(B_p^r\)'s branch is \(B_p^r\) or an ancestor of it, so that justification's
height exceeds \(k\geq H-1\).  A justification at height \(H\) or
above is excluded before the first confirmation by
Lemma~\ref{lem:selectable-roots}.
The action-admission dichotomy is therefore exhaustive: an admitted honest
root lies on the chain ending at \(B_p^r\), or off that chain and
supersedes the proposal through item~2 or item~3; an off-path root
arriving with a finalized height first processed at the round's cutoffs
is item~4.

If no grade-2 root exists, the proposer supplies its Simplex root as
\(A_p^r\), an ungraded base root, and relays the justification or
finality certificate that backs it as ordinary evidence.  By
Lemma~\ref{lem:g3-g2}, a receiver whose grade-3 lower root is strict at
that receiver would give the proposer grade~2 for it, so no receiver
holds one, and every
receiver's lower root falls back to its vote-state Simplex root.  In the
present branch the receivers' store maxima, active finality outcomes, and
processed finalized evidence agree with the proposer's, so the selection
rule of Definition~\ref{def:finality-root}, a function of exactly these
inputs, gives every receiver
\(S_{u,\mathrm{vote}}^r=A_p^r\), and the equality disjunct of item~2 of
Definition~\ref{def:stable-root} accepts.  In both the graded and base
cases the receiver's vote-time stable root lies
on the path from \(A_p^r\) through \(Z_p^r\).  Item~3 holds as well:
blocks on the proposal path are exempt from witness aging by
Definition~\ref{def:counting-rule}, and the receiver's remaining
roots are aged-backed or exempt --- its grade-3 lower root is chosen
within its own aged trees, and its vote-state Simplex root lies on the
exempted proposal path in every case: at or below \(A_p^r\), or
between \(A_p^r\) and \(Z_p^r\) under the newer-finalized disjunct,
and Definition~\ref{def:counting-rule} exempts that path's blocks
together with their ancestors.  All vote-time
acceptance checks therefore hold.

The remaining checks can fail only if the proposal or its root leaves the
store-global viable subtree, or if the action Simplex root moves across it.
The first requires the global maximum to rise by enough to place the
proposal below the \(h_{\max}-1\) band, which is item~2.  An action Simplex
root on the proposal path keeps the action-state admission of
Definition~\ref{def:action-root} satisfied for the vote-time root; only
an off-path root fails that action
check.  Such a root is the unique justification redirection already
classified, unless a
height-\(H+1\) honest confirmed witness exists.  That witness is already the
first recovery confirmation.  Recognition, the first conjunct of
Definition~\ref{def:usable-honest-round}, holds by
Lemma~\ref{lem:common-recognition}: that lemma's premise is this
lemma's premise --- the common-recognition event, not the mere existence
of an honest eligible signer.  In every remaining execution the proposal
is unobstructed, which is item~1.  The four cases are exhaustive.

\end{proof}



\begin{lemma}[Every honest action has reached the recovery height at the first confirmation]
\label{lem:first-confirmation-simplex-roots}
In action \(r_\star\), every honest action state has store-global maximum
\(H+1\).  Its Simplex root is its finalized root and is an ancestor of
\(B_\star\).

In plain words, the first honest confirmation cannot outrun delivery of the
recovery-height chain, and the nonjustifiable height puts every finality
fork-choice root below that confirmation.
\end{lemma}

\begin{proof}
Let \(C_\star\) be the block from which the confirming validator derives
\(B_\star\).  By Definition~\ref{def:official-confirmation},
\(B_\star\preceq C_\star\) and \(C_\star\) is available-confirmed at the
round's evaluation; only the current round's evaluation is read, so no
case for a carried earlier confirmation arises.
Lemma~\ref{lem:tsq-common-knowledge} puts \(C_\star\) and its full
ancestry in every honest validator's store by \(d_r+3\Delta\), three
delivery bounds before the action.  Since \(B_\star\) has state-height at
least \(H+1\) by Definition~\ref{def:first-recovery-confirmation} and
\(B_\star\preceq C_\star\), every honest action store contains a
state-height-\(H+1\) chain.

Equivalently, if \(t\) is the first time an honest validator processes a
chain at height \(H+1\), no honest confirmation at that height occurs before
\(t+\Delta\).  The counted support contains an honest member by
Corollary~\ref{cor:tsq-honest-witness}, and that member's vote followed
its own processing of an \(H+1\) chain.
Under the strict delivery bound, material processed at \(t\) is common
strictly before \(t+\Delta\), so the earliest confirmation follows the
first honest \(H+1\) chain by a full bound.  Faulty votes sent in
advance do not change this argument because faulty weight alone cannot
supply the required honest vote witness.

Fix one of these honest action stores and write \(F\) for its greatest
finalized root.
No finalized root has height \(H\), because finalizing at \(H\) would require
a height-\(H\) justification, which
Lemma~\ref{lem:no-height-h-justification} excludes.  Nor can \(F\) have height at
least \(H+1\).  Its finality certificate contains a justification certificate
at that height.  The target quorum of that justification contains an honest
signer because \(q>b\), and
Lemma~\ref{lem:height-vote-confirmation} gives that signer an official
confirmation at height at least \(H+1\) before the certificate entered this
action state.  That contradicts the choice of \(r_\star\).
Thus \(h_F<H\).  Since \(B_\star\) has state-height at least \(H+1\),
Lemma~\ref{lem:past-finalized} gives \(F\preceq B_\star\).  The
height-\(H+1\) chain established above remains in this store's
finalized-root subtree and makes its filtered maximum at least \(H+1\).

No admitted branch before the action in \(r_\star\) has state-height above
\(H+1\).
Otherwise the first transition from \(H+1\) to \(H+2\) on that branch would,
by Lemma~\ref{lem:height-increase-confirmation}, have an honest confirmed
witness at height at least \(H+1\) before \(r_\star\).  This contradicts
Definition~\ref{def:first-recovery-confirmation}.  Hence every honest action
store has maximum exactly \(H+1\).

Lemma~\ref{lem:no-height-h-justification} now makes each action-state Simplex root
equal to that store's finalized root \(F\).  The preceding paragraph already
proved \(F\preceq B_\star\), which is the claim.
\end{proof}

\begin{lemma}[An unobstructed honest round emits SG heads compatible with its proposal]
\label{lem:honest-proposal-sg-batch}
Let \(B\) be the proposal in an unobstructed honest-proposer round.  Every honest
SG head signed in that round is compatible with \(B\), and an honest
validator whose G0 checks admit no such head abstains.

In plain words, the exact proposal batch produces one compatible SG batch
even if stale grade~0 evidence vetoes \(B\) itself.
\end{lemma}

\begin{proof}
Lemma~\ref{lem:honest-proposer-votes} gives the exact \(B\)-vote batch.
The action checks in Definition~\ref{def:usable-honest-round} keep \(B\)
viable and every vote-time and action root on its ancestry chain.
Lemma~\ref{lem:tsq-liveness} therefore available-confirms \(B\) at every
honest validator's evaluation, and Corollary~\ref{cor:tsq-walk} --- its
root and candidate-tree premises are exactly those action checks --- makes
the head \(G_u^r\) at \(a_r\) descend from \(B\).  A nonempty SG head
\(Q_u^r\preceq G_u^r\) is therefore compatible with \(B\): both \(Q_u^r\)
and \(B\) lie on the ancestry chain of \(G_u^r\).
If no official confirmation exists, the fallback is the deepest veto-free
block between the action Simplex root and action root.  The action root is an
ancestor of \(B\) by Definition~\ref{def:stable-root}, so the fallback is compatible with
\(B\).  If no such veto-free block exists, the validator abstains.
All signed cases are therefore compatible with \(B\).
\end{proof}

\begin{lemma}[The following batch has no conflicting grade]
\label{lem:next-batch-clears}
In the round after Lemma~\ref{lem:honest-proposal-sg-batch}, no block
conflicting with \(B\) has grade \(3,2,1,\) or \(0\) in an honest view.

In plain words, the exact honest-proposer round removes stale grades in one
eligible batch.
\end{lemma}

\begin{proof}
Every eligible honest head is compatible with \(B\), so
Lemma~\ref{lem:fresh-support-clears} applies.
\end{proof}

\begin{lemma}[Mid-round re-derivation stays compatible with the proposal]
\label{lem:rederivation-ancestry}
In a round in which every action-state Simplex root is compatible with
a block \(B\), no store-maximum rise supersedes the proposal, and every
justification selectable by any of the validator's round states, and
every activated finalized root, is compatible with \(B\), every
root produced by the mid-round re-derivation of
Definition~\ref{def:rederivation} is compatible with \(B\).

In plain words, re-deriving at a slot boundary cannot move a stable root
off the proposal's chain unless the chain itself has moved on.
\end{lemma}

\begin{proof}
The re-derived root is the validator's current Simplex selection at a
slot boundary.  Its activated finality outcomes are the round's own by
the filter constancy of Definition~\ref{def:activation-filter}, so the
selection can differ from the action state's only through the store
maximum, which decides between the two branches of
Definition~\ref{def:finality-root}.  On the justification branch the
selected root is a selectable justification, compatible with \(B\)
directly by the third premise.  On the finalized branch the selected
root is the
round's activated finalized root, compatible with \(B\) by the same
premise.  A
superseding rise is excluded by the second premise.
\end{proof}

\begin{lemma}[Processed finality conflicts with no proposal before the first confirmation]
\label{lem:pre-confirmation-finality-compatible}
Let \(B\) be the proposal of an unobstructed honest round with
post-state height \(k\in\{H-1,H,H+1\}\), before the first recovery
confirmation, and suppose every finalized root whose evidence an honest
validator has processed either has height strictly below \(k\) or
lies on \(B\)'s chain.  Then no such root conflicts with \(B\).
The extra premise is automatic for \(k\geq H\), because
Lemma~\ref{lem:selectable-roots} bounds every processed root strictly
below \(H\leq k\); for \(k=H-1\) it excludes exactly one
situation --- a processed height-\((H-1)\) root conflicting with
\(B\) --- which is the redirection case the callers classify
separately.

In plain words, until the first confirmation, processed finality stays
on every unobstructed honest proposal's chain, except possibly the one
old root at the proposal's own height.
\end{lemma}

\begin{proof}
A processed root of height strictly below \(k\) lies on \(B\)'s
chain by Lemma~\ref{lem:past-finalized}, applied to the chain ending
at \(B\); a root on \(B\)'s chain does not conflict by definition.
Roots at height \(H\) or above do not exist before the first
confirmation, by Lemma~\ref{lem:selectable-roots}.
\end{proof}

\begin{lemma}[Gate closure at a closing action]
\label{lem:gate-closure}
At a closing action, suppose that at every honest validator: the deepest
available-confirmed block contains \(B\); the action-state Simplex
root is an ancestor of \(B\); no active grade~0 conflicts with
\(B\); no processed finalized root conflicts with \(B\); the
action-admission accepts, so no honest validator abstains; and every
vote-time stable root of the closing round is compatible with \(B\)
with the path to \(B\) in its candidate tree.  Then every honest
validator officially confirms in that action, and its deepest official
confirmation contains \(B\).

In plain words, once the gate's conjuncts and admission hold with
\(B\) on every path, the gate produces a confirmation containing
\(B\) everywhere.
\end{lemma}

\begin{proof}
The head conjunct holds by Corollary~\ref{cor:tsq-walk} applied to
\(B\): every vote-time stable root of the closing round is compatible
with \(B\) with the path in its candidate tree, so every action head
descends from \(B\).  With the available confirmation, the ancestral
Simplex root, the veto-free conjunct, and the processed-finality
conjunct supplied by the premises, all five gate conditions of
Definition~\ref{def:official-confirmation} hold for \(Q=B\), and
admission accepts by premise.  An official confirmation therefore exists
at every honest validator.  Every official confirmation is an ancestor
of the deepest available-confirmed block, which contains \(B\), and
\(Q=B\) itself passes, so the deepest official confirmation is at
least as deep as \(B\) on that chain and contains \(B\).
\end{proof}

\begin{lemma}[The closing action of a recovery-height honest round confirms officially]
\label{lem:closing-action-existence}
Let \(B\) be the proposal of an unobstructed honest round with post-state
height \(H+1\), and suppose no first recovery confirmation occurs before
the round's closing action.  Then at the closing action every honest
validator has an official confirmation, and its deepest official
confirmation contains \(B\).

In plain words, the gate does not merely constrain confirmations at the
closing action; it produces one at every honest validator.
\end{lemma}

\begin{proof}
Lemma~\ref{lem:honest-proposer-votes} gives the exact honest first-slot
vote batch for \(B\).  We first check the premises of
Lemma~\ref{lem:recovery-rooted-persistence} for every later slot through
the following round's first slot.  Within the proposal round, the
vote-time stable roots are fixed on \(B\)'s ancestry by the action checks
of Definition~\ref{def:usable-honest-round}.
Lemma~\ref{lem:rederivation-ancestry} covers every mid-round
re-derivation --- its justification premise holds vacuously at any
re-deriving validator: that validator accepted \(B\), so its maximum
is at least \(H+1\), a selectable justification would need height
\(\hmax-1\geq H\), and Lemma~\ref{lem:selectable-roots} leaves none
at or above \(H\) --- so the roots used by the round's
later slots stay compatible with \(B\).  For the following round, three root sources need checking.  Its G3
lower roots come from the proposal round's SG batch, which
Lemma~\ref{lem:honest-proposal-sg-batch} keeps compatible with \(B\).
Its Simplex roots are ancestors of \(B\).  By
Lemma~\ref{lem:selectable-roots}, every selectable justification has
height at most \(H-1\) and every finalized root height below \(H\).
The
unique height-\((H{-}1)\) justification cannot be selected once
\(B\) is accepted, because every honest store maximum is then at least
\(H+1\geq(H-1)+2\), so the justification guard \(\hmax=h_j+1\) of
Definition~\ref{def:finality-root} fails for it.  An accepted ungraded
proposed root adds no further case, because item~2 of
Definition~\ref{def:stable-root} accepts it only when it equals the
receiver's own Simplex selection.  Last, a finality outcome active at
any honest validator lies on the finalized chain by
Lemma~\ref{lem:finalized-chain}, and every finalized root below
\(B\)'s state-height lies on \(B\)'s chain by
Lemma~\ref{lem:past-finalized}, so even a staggered activation leaves
every vote-time stable root compatible with \(B\).  For the candidate trees,
Lemma~\ref{lem:pre-confirmation-cap} keeps every store maximum at most
\(H+1\), and \(B\) has state-height \(H+1\geq\hmax-1\), so \(B\) and
its ancestry stay in every honest candidate tree.  With these premises,
Lemma~\ref{lem:recovery-rooted-persistence} keeps every honest head and
raw vote in \(B\)'s subtree through the round's remaining slots and the
following round's first slot.  Every
honest first-slot vote of the following round therefore descends from
\(B\), so Lemma~\ref{lem:tsq-liveness} available-confirms \(B\) at every
honest validator at the closing action's evaluation, and the deepest
available-confirmed block \(C_u\) contains \(B\).

The action-state Simplex root is an ancestor of \(B\):
Lemma~\ref{lem:selectable-roots} excludes a justification or finalized
root at height \(H\) or above, and Lemma~\ref{lem:past-finalized}
places every lower finalized root on \(B\)'s chain.  Finally, Lemma~\ref{lem:next-batch-clears} removes every
grade on a block conflicting with \(B\), so \(B\) is veto-free.  The
head conjunct holds by Corollary~\ref{cor:tsq-walk} applied to \(B\):
every vote-time stable root of the closing round is compatible with \(B\) with
the path in its candidate tree, so every action head descends from
\(B\).  The processed-finality conjunct is
Lemma~\ref{lem:pre-confirmation-finality-compatible}, whose
strict-height premise is automatic at \(k=H+1\).  No honest validator takes the abstention branch of
Definition~\ref{def:official-confirmation} at the closing action: each
honest vote-time root of the closing round is a grade-3 lower root,
which carries raw grade~1 because
\(D_u^{-,1}(X)\subseteq\Phi_{X_u^1}(X)\); an accepted grade-2 proposer
root, which carries raw grade~1 by item~2 of
Definition~\ref{def:stable-root}; or the vote-state Simplex selection,
which equals \(S_{u,\mathrm{act}}\) because activation is constant
within the round and Lemma~\ref{lem:pre-confirmation-cap} pins the
store maximum, so the admission accepts in every case; the ancestry
conditions hold for any closing proposal by the low common Simplex
roots, the proposal-free branch reads no proposal block, and no
re-derivation fires absent the separately handled moves.  The premises
of Lemma~\ref{lem:gate-closure} are now all in place, and it gives
every honest validator an official confirmation whose deepest contains
\(B\).
\end{proof}

\begin{lemma}[An unobstructed proposal is confirmed or one old root redirects it]
\label{lem:honest-proposal-confirms}
Let \(B\) be the proposal in an unobstructed honest round with post-state
height \(k\in\{H-1,H,H+1\}\), and suppose the first recovery confirmation has not
yet occurred.  Let \(r_B\) be the closing action of the proposal round: the following
round's action, which grades the proposal round's SG batch and has
completed its own round's first-slot available-confirmation evaluation.  One of the
following occurs before or in \(r_B\):
\begin{enumerate}
 \item every honest validator officially confirms in \(r_B\), and every
   honest official confirmation in \(r_B\) contains \(B\);
 \item if \(k=H-1\), the store-global maximum rises to at least \(H+1\), so
   the next selected proposal has a higher permitted height;
 \item if \(k\in\{H-1,H\}\), the unique height-\(H-1\) justification becomes the
   Simplex root on a fork of \(B\); after that state is delivered, the next
   unobstructed proposal at a permitted height extends that root and this case cannot
   recur;
   or
 \item the first recovery confirmation has already occurred.
\end{enumerate}

In plain words, an old height-\(H-1\) justification can redirect one honest
proposal.  A rise on another branch can raise the permitted heights, but is not
called progress on \(B\)'s branch.
\end{lemma}

\begin{proof}
Lemma~\ref{lem:honest-proposer-votes} gives the exact \(B\)-vote batch, and
Lemma~\ref{lem:honest-proposal-sg-batch} puts every honest SG head from that
round compatible with \(B\) or makes it empty.

We first classify the two ways in which that exact batch could fail to
persist.  The first is removal by the height filter while \(B\) remains in
the finalized-root subtree.  Because \(B\) has state-height \(k\), it remains
viable there while the store-global maximum is at most \(k+1\).  If it ceases
to be viable for this reason, the maximum has reached at least \(k+2\).  For
\(k=H-1\), this means that the store-global maximum is at least \(H+1\),
which is item~2.  This changes the next permitted heights; it does not replay the
transition on \(B\)'s branch.  For \(k\geq H\), the
transition from \(H+1\) to \(H+2\) has already occurred.  By
Lemma~\ref{lem:height-increase-confirmation}, that transition has an honest
confirmed witness at height at least \(H+1\), which is item~4.

The second possible failure is a new action-state Simplex root \(R\) that
conflicts with \(B\).  If \(R\) is finalized at height \(y<k\),
Lemma~\ref{lem:past-finalized} gives \(R\preceq B\), so a conflicting
finalized root must have \(y\geq k\).  Height \(H\) cannot be justified or
finalized by Lemma~\ref{lem:no-height-h-justification}.  Thus a conflicting finalized
root at height at least \(H\) has height at least \(H+1\), and its
justification quorum contains an honest confirmed witness at height at least
\(H+1\) by Lemma~\ref{lem:height-vote-confirmation}.  This is item~4.  The
only remaining finalized-root case is \(k=H-1\) and \(y=H-1\).  Processing
that finality certificate reveals its height-\(H-1\) justification and
carrying state, which is item~3.

Suppose instead that \(R\) is a justification root.  Before the first
recovery confirmation, the admitted maximum is at most \(H+1\), by the same
height-increase argument as in the preceding paragraph.  The Simplex-root
guard selects a justification only one height below that maximum.
The height-\(H\) justification is disabled, and none exists above
\(H\) (Lemma~\ref{lem:selectable-roots}).  The recovery store already has
maximum at least \(H\), and
Lemma~\ref{lem:filtered-hmax-monotone} preserves that lower bound as finalized
evidence arrives.  Thus the only remaining selected justification that can
conflict with \(B\) is at height \(H-1\) while the maximum is \(H\).  If
\(k=H-1\), the carrying state raises the selected band but does not advance
\(B\)'s branch; this is item~3.  If \(k=H\), it is also item~3.  If
\(k=H+1\), no such justification is selectable at all: the unobstructed
proposal's recognition places its height-\((H{+}1)\) chain in every
honest store, so every honest maximum is at least \(H+1\) and the
selection guard \(\hmax=h_j+1\) fails for \(h_j=H-1\); items~2 and~3
are rightly inapplicable there.
Lemma~\ref{lem:target-uniqueness} prevents a second conflicting
height-\(H-1\) justification.  A later finalized root below \(H\) is either
an ancestor of the retry proposal by Lemma~\ref{lem:past-finalized}, or is
the same height-\(H-1\) target by
Lemma~\ref{lem:finalized-chain}.  Hence item~3 cannot recur after the
justification and its state are delivered.

It remains to prove item~1 when neither failure occurs.
Lemma~\ref{lem:next-batch-clears} removes every grade conflicting with \(B\).
Every action-state Simplex root through \(r_B\) is compatible with \(B\) by
the failure classification.  A selected justification lies on \(B\)'s chain,
and every compatible finalized root lies on that chain as well.  Every stable root after the grade computation
is therefore compatible with \(B\): it is either that Simplex root or a
grade-supported root --- an accepted ungraded proposed root equals the
receiver's own Simplex selection by item~2 of
Definition~\ref{def:stable-root}, so no third case exists --- and a
conflicting root strict at its holder would require a
conflicting grade.  Lemma~\ref{lem:rederivation-ancestry} covers every mid-round
re-derivation --- its justification premise holds by the Simplex-root
placement above: no selectable justification conflicts with \(B\) or
strictly descends from it in the present branch, so every one is an
ancestor --- so the roots used by the
round's later slots stay compatible with \(B\).  The subtree of \(B\) remains viable.  Starting from the
exact \(B\)-vote batch, Lemma~\ref{lem:recovery-rooted-persistence} now makes
every honest Goldfish head and raw vote descend from \(B\).

No honest validator takes the abstention branch of
Definition~\ref{def:official-confirmation} at this action.  The closing
round's action-state Simplex roots are, by the enumeration above, the
finalized root or the unique low justification, each of height below
\(B\)'s and an ancestor of every branch a closing proposal can extend
--- acceptance requires the receiver's selection below the parent ---
so the ancestry conditions hold whether or not the closing proposal
extends \(B\); the vote-time root is an ancestor of the accepted
parent, or the same-checked fallback root otherwise; the proposal-free
branch reads no proposal block; and absent the
separately handled mid-round moves, no re-derivation fires.  The guard
therefore passes at every honest validator, and since each validator's
walk root is compatible with \(B\) with the path in its candidate
tree, Corollary~\ref{cor:tsq-walk} gives \(B\preceq G_u\) even when
the closing proposal lies elsewhere, so the gate conditions decide for
\(Q=B\).  At that round's evaluation,
Lemma~\ref{lem:recovery-rooted-confirmation} available-confirms \(B\) at
every honest validator, so every honest confirmed chain reaches \(B\).
Before the first recovery confirmation, no Simplex root strictly descends
from \(B\): such a root is a justification or finalized root whose
target strictly descends from \(B\), which requires a justification at a
height above \(B\)'s, excluded exactly as in the strict-descendant
paragraph of Lemma~\ref{lem:honest-proposer-coverage}'s proof.  The
action-state Simplex root is therefore an ancestor of \(B\).  The
premises of Lemma~\ref{lem:gate-closure} are in place: the deepest
available-confirmed block contains \(B\); the Simplex root is an
ancestor; no grade~0 conflicts with \(B\); every vote-time stable root
is compatible with \(B\) with the path in its candidate tree;
admission accepts by the abstention paragraph above; and the
processed-finality conjunct is
Lemma~\ref{lem:pre-confirmation-finality-compatible}, whose
strict-height premise holds in the present branch: the failure
classification has excluded a conflicting processed root at \(B\)'s
own height, and every other processed root is strictly lower or on
\(B\)'s chain.
Lemma~\ref{lem:gate-closure} gives every honest validator an official
confirmation at \(r_B\) whose deepest contains \(B\).  This is
item~1.
\end{proof}

This is the liveness route to
Definition~\ref{def:first-recovery-confirmation}.  It is not the definition
of the safety event: if a faulty proposer produces the first honest official
confirmation at height \(H+1\), the proof below starts from that earlier
confirmation.

\section{Healing from the first recovery confirmation}\label{sec:healing}

This section first proves that the first honest official confirmation at
height \(H+1\) remains a viable prefix.  It then proves that an honest round
proposer forces such a confirmation if none occurred earlier, and gives the
height and finality steps that follow.

\begin{lemma}[No pre-recovery quorum exists above \(H\)]
\label{lem:no-hidden-h-advance}
Before the relay point of Definition~\ref{def:relay-frontier}, no
valid target or progress quorum exists at height \(H+1\) or above.

In plain words, old votes may already move height \(H\), but they cannot also
move the next height.
\end{lemma}

\begin{proof}
Every target or progress quorum has weight at least \(q\), so it contains an
honest signer because \(b<q\).  Assumption~\ref{ass:frontier-opportunities}
states that no honest validator signed a nonempty height pair above \(H\)
before the relay point.  Faulty signatures alone cannot form the quorum.
\end{proof}

\begin{lemma}[The first recovery confirmation classifies every SG head]
\label{lem:first-confirmation-heads}
Let \(F_\star\) be the finalized root in the action state of the honest
validator that officially confirms \(B_\star\).  Every honest action root
admitted by the action-state check, every available confirmation at that
action's evaluation, and every nonempty SG head in action
\(r_\star\) is either compatible with \(B_\star\) or conflicts with a
finalized root whose evidence that validator has processed ---
\(F_\star\), or a newer processed outcome on the same finalized chain.
Evidence in the second case is inert in the confirming state
and in every honest state after the finalization evidence is relayed,
which happens strictly before the following round's grade snapshots.

In plain words, every root that can affect the SG/FG action is either on the
confirmed chain or is excluded by the confirmer's finalized root; a
validator still holding an unusable older vote root abstains.
\end{lemma}

\begin{proof}
Lemma~\ref{lem:first-confirmation-simplex-roots} gives store-global maximum
\(H+1\) in every action state in \(r_\star\), keeps the height-\(H+1\)
chain viable there, and makes every action-state Simplex root an ancestor of
\(B_\star\).  Lemma~\ref{lem:official-confirmation-root-safety} therefore
makes every action root admitted by the action-state check compatible with
\(B_\star\) or conflicting with a finalized root whose evidence the
confirming validator has processed; processed finalized pairs lie on one
chain through \(F_\star\) by Lemma~\ref{lem:finalized-chain}, the
confirming validator relays the evidence at once, and one delivery bound
lands strictly before the drain round's earliest grade snapshot.  A validator whose fixed vote
root fails that check emits no SG head or current-height pair, by
Definition~\ref{def:stable-root}.

Lemma~\ref{lem:tsq-uniqueness} makes every available confirmation at the
action's evaluation compatible with \(C_\star\), and hence with
\(B_\star\): a confirmation at \(C_\star\) or a descendant of it
contains \(B_\star\), and an ancestor of \(C_\star\) lies on its
chain.  No
persistence fact is used at this point: the conclusion is compatibility,
not descent from \(B_\star\).

There are two SG-head cases in
Definition~\ref{def:official-confirmation}.
\begin{enumerate}
 \item If the validator has an official confirmation, its SG head is that
   confirmation and satisfies the stated alternative.
 \item Otherwise its nonempty SG head lies between its Simplex root and
   action root.  If the action root is compatible with \(B_\star\), every
   block on that ancestry path is compatible as well.  If it conflicts
   with a processed finalized root, each block on that path is either
   still on \(B_\star\)'s ancestry or lies after the fork and conflicts
   with the same processed finalized root.
\end{enumerate}
These cases exhaust every nonempty honest SG head.
\end{proof}

\begin{lemma}[The batch after the first confirmation has no live conflicting grade]
\label{lem:first-confirmation-clears}
In the first action that grades the SG heads from \(r_\star\), no block
conflicting with \(B_\star\) has an active grade \(3,2,1,\) or \(0\) in an
honest view.

In plain words, one fresh SG batch removes every live grade on the other
fork; finalized-conflicting evidence may remain stored.
\end{lemma}

\begin{proof}
By Assumption~\ref{ass:recovery-network}, the evidence for \(F_\star\) ---
and for every finalized root any honest head's holder had processed,
since a holder relays what it processes --- has
reached every honest validator before this action.
Lemma~\ref{lem:first-confirmation-heads} splits the eligible honest
heads: a head compatible with \(B_\star\) supports no conflicting
block; a head excluded by a processed finalized root \(R\) ---
\(F_\star\) itself or a newer processed outcome on the same chain, with
\(R\preceq B_\star\) by Lemma~\ref{lem:past-finalized} since \(R\)'s
height is below \(B_\star\)'s state-height --- can support a
conflicting block \(X\) only if \(X\) lies in the head's subtree.  If
\(X\preceq R\), then \(X\preceq R\preceq B_\star\) and \(X\) does not
conflict with \(B_\star\) at all; otherwise \(X\) conflicts with
\(R\), whose evidence every honest grader has processed by this
action's grade freeze, so Definition~\ref{def:active-grade} makes every
grade for \(X\) inactive at every honest grader.  An active block
conflicting with \(B_\star\) therefore receives favorable support only
from faulty
identities, of weight \(b<m\).  Definition~\ref{def:active-grade} and
Definition~\ref{def:grades} exclude all four active grades.
\end{proof}

\begin{lemma}[Height increases before the drain retain the first confirmation]
\label{lem:predrain-height-retention}
Let \(F_\star\) be the finalized root named in
Lemma~\ref{lem:first-confirmation-heads}.  After every height increase
accepted between \(r_\star\) and the drain action, \(B_\star\)
remains viable in every honest store, unless that store's finalized root
already descends from \(B_\star\), in which case every accepted block
of its restricted tree descends from \(B_\star\).  The store-global
maximum rises at most once before the drain, from \(H+1\) to
\(H+2\).

In plain words, no height advance between the confirmation and the drain
pushes the confirmed block out of the viable band.
\end{lemma}

\begin{proof}
We check every height increase accepted before the drain action.  A branch-local increase that does not raise the
store-global maximum leaves the viability threshold unchanged, so the accepted
\(B_\star\)-subtree keeps its existing witness.  Suppose instead that such an
increase raises the maximum
from \(x\) to \(x+1\).  If \(x\leq H\), the block \(B_\star\), whose
state-height is \(H+1\), already
witnesses the height-filter threshold \(x\) in its own subtree, and no
witness lemma is needed.  Otherwise \(x\geq H+1\), and before the drain
exactly \(x=H+1\), since a larger \(x\) would require an honest height
pair above \(H+1\); Lemma~\ref{lem:height-increase-confirmation},
stated for increases from \(x\geq H+1\), now applies: its quorum
contains an honest
signer with an official confirmation \(Q\) at height at least \(x\).
The honest signature cannot predate \(r_\star\), by
Definition~\ref{def:first-recovery-confirmation}.  It was released in
\(r_\star\), because no other SG/FG action occurs before
\(r_{\rm drain}\), and
Lemma~\ref{lem:first-confirmation-heads} makes \(Q\) compatible with
\(B_\star\) unless \(Q\) conflicts with a processed finalized root.  The latter branch is
outside the finalized-root subtree once the relayed finalization is
processed, by Lemma~\ref{lem:finalized-lockin}, and cannot raise that
subtree's maximum.  In the remaining case,
either \(B_\star\preceq Q\), so its subtree has
the height-\(x\) witness, or \(Q\preceq B_\star\), so state-height
monotonicity gives the witness at \(B_\star\).  No larger \(x\) can occur
before the drain, because its quorum would require an honest height pair
above \(H+1\), and no intervening action can produce one.
Assumption~\ref{ass:frontier-opportunities} delivers every witness and its
ancestry before the new maximum is used.  Thus every possible increase before
the drain retains \(B_\star\), unless the finalized root has already passed
it, in which case every accepted block descends from it.  A timeout cannot
stand in for a second increase because
Definition~\ref{def:height-outcome} resets the participation arrays.
\end{proof}

\begin{lemma}[Simplex-root updates through the drain stay ancestral]
\label{lem:predrain-simplex-roots}
Every finality-derived Simplex root selected by an honest validator from
\(r_\star\) through the drain action is an ancestor of \(B_\star\).

In plain words, no root selected between the confirmation and the drain
leaves the confirmed chain.
\end{lemma}

\begin{proof}
Before the drain the store-global maximum is \(H+1\) or \(H+2\):
it cannot exceed \(H+1\) at \(r_\star\) by
Lemma~\ref{lem:first-confirmation-simplex-roots}, and
Lemma~\ref{lem:predrain-height-retention} permits one rise.  We account
for every update.
If the store-global maximum is \(H+1\),
Lemma~\ref{lem:no-height-h-justification} selects the finalized root.  Its finalized
height is below \(H\), and Lemma~\ref{lem:past-finalized} puts it on the
ancestry of \(B_\star\).  If the maximum is \(H+2\), the only justification
that can be selected has height \(H+1\).  Its quorum contains an honest
target signer because \(q>b\).  That signature cannot predate \(r_\star\);
Lemma~\ref{lem:first-confirmation-heads} says that its official confirmation
is compatible with \(B_\star\) or conflicts with a processed finalized
root, in which case relay makes it inert everywhere by
Lemma~\ref{lem:first-confirmation-heads}.
The confirmation has state-height at least \(H+1\) by
Lemma~\ref{lem:height-vote-confirmation}.  Every finalized root the
signer had processed by its \(r_\star\) action has height below
\(H\) --- that processing happened before the first confirmation, so
Lemma~\ref{lem:selectable-roots} applies --- and each such root
therefore lies on the confirmation's chain by
Lemma~\ref{lem:past-finalized}, applied to that root; the confirmation
conflicts with none of them, which excludes the second alternative.  Hence the
confirmation is compatible with \(B_\star\), and
Definition~\ref{def:height-vote-rule} makes the target an ancestor of that
confirmation.  Lemma~\ref{lem:chain-target-uniqueness} therefore puts the
target on the ancestry of \(B_\star\).  A newly finalized root is either
below \(H\), in which case Lemma~\ref{lem:past-finalized} applies, or copies
this height-\(H+1\) target.  No height-\(H\) root exists, by
Lemma~\ref{lem:no-height-h-justification}, and no root above \(H+1\) can be selected
before the drain, by the height-pair argument of
Lemma~\ref{lem:predrain-height-retention}.  These cases exhaust Definition~\ref{def:finality-root} and
make every selected root through the drain an ancestor of \(B_\star\).
\end{proof}

\begin{lemma}[Drain stable roots are compatible with the first confirmation]
\label{lem:drain-stable-roots}
At the drain action, every honest validator's action-state finalized
root contains \(F_\star\), and every honest drain stable root is
compatible with \(B_\star\).

In plain words, by the drain action the confirmed finality outcome is
active everywhere and every root respects it.
\end{lemma}

\begin{proof}
Lemma~\ref{lem:first-confirmation-clears} removes every grade conflicting
with \(B_\star\) that is active in \(r_{\rm drain}\).  Activation has
caught up by then: \(F_\star\) was activated at the confirming validator
in \(r_\star\), so its carrying accepted block was received by
\(a_{r_\star-1}\) and relayed on receipt, reaching every honest
validator strictly before \(a_{r_\star}\) --- the drain round's
activation cutoff --- so every honest validator activates \(F_\star\) at
the drain, and activated finalized pairs form one chain by
Lemma~\ref{lem:finalized-chain}, so every honest action-state finalized
root at the drain contains \(F_\star\).  The action-state Simplex roots are
compatible by Lemma~\ref{lem:predrain-simplex-roots}, so the three cases in
Definition~\ref{def:stable-root} make every drain stable root compatible as
well.
\end{proof}

\begin{lemma}[Later-slot roots and votes stay on the confirmed side]
\label{lem:confirmation-round-slots}
In \(r_\star\)'s round, every root used by the second and later
slots at any honest validator is compatible with \(B_\star\), and
either lies with \(B_\star\) in that validator's candidate tree or
is a strict descendant of \(B_\star\) evicted by a mid-round rise,
whose walk the standing rule keeps inside \(B_\star\)'s subtree;
every honest raw vote of those slots and of the drain round's first
slot lies in \(B_\star\)'s subtree.

In plain words, the slots between the confirmation and the drain vote
inside the confirmed subtree even though their roots were fixed early.
\end{lemma}

\begin{proof}
The round's later slots need a separate
argument, because their roots were fixed at \(d_{r_\star}+\Delta\),
possibly before the height-\((H{+}1)\) chain reached every honest
validator.  Fix an honest validator \(v\) and split
on how its round root was selected.

\emph{Case: a graded root.}  If the root \(R\) is a graded root,
suppose toward a contradiction that \(R\) conflicts with \(B_\star\).
The root has grade~1 at \(v\) by admission, and
Lemma~\ref{lem:g1-g0} gives the confirming validator \(u\) raw grade~0
for \(R\).  Split on the finalized evidence conflicting with \(R\), if
any, that \(u\) had processed by its grade-0 freeze
\(d_{r_\star}+2\Delta\).  If there is none, then \(R\) conflicts
with no finalized root processed by the freeze; \(F_\star\) is among
those roots --- it is activated in \(u\)'s action state, so its carrying
block was received and accepted by \(a_{r_\star-1}\leq
d_{r_\star}-2\Delta\)
(Definition~\ref{def:activation-filter}), strictly before the freeze --- so \(R\) is
compatible with \(F_\star\), and \(R\) descends
from \(F_\star\), since compatibility
and the assumed conflict with \(B_\star\succeq F_\star\) exclude the
ancestor side.  \(R\) is in the confirming validator's accepted
subtree, with its full ancestry, strictly before the freeze, by
Lemma~\ref{lem:g1-acceptance}.  Hence \(u\)'s grade~0 for \(R\) is active at the freeze
and vetoes \(B_\star\)'s official confirmation by
Definition~\ref{def:official-confirmation}, contradicting the choice of
\(r_\star\).  If instead some finalized evidence conflicting with
\(R\) was processed by \(u\) by \(d_{r_\star}+2\Delta\), it was relayed
on receipt and processed by every honest validator strictly before
\(d_{r_\star}+3\Delta\); in particular \(v\) processes it before the
second slot boundary \(d_{r_\star}+4\Delta\), the grade-inactivity exit
of Definition~\ref{def:rederivation} makes \(R\) count as ungraded
from that boundary, and the ungraded case below re-derives \(v\)'s root
before the second-slot vote at \(d_{r_\star}+5\Delta\).  If instead
\(R\)'s grade is inactive against
the processed \(F_\star\) evidence, the root is not a live graded root
at all: the confirming validator activated \(F_\star\) for its round, so
the carrying block was received by the preceding action cutoff and
relayed on receipt, every honest validator processes the evidence
strictly before the round's first proposal time, and
Definition~\ref{def:active-grade} removes the conflicting root from
grade-based retention from vote time on.  The validator holds an
ungraded root instead, and the ungraded case below covers it.  In every
case, no honest second-slot walk starts from a live graded root
conflicting with \(B_\star\).

\emph{Case: an ungraded root.}  If instead the root is ungraded, it is the validator's own vote-state
Simplex selection --- as a fallback, or as an accepted ungraded proposed
root, which equals that selection by item~2 of
Definition~\ref{def:stable-root} --- possibly the unique
height-\((H{-}1)\) justification admitted by the validator's own
selection guard.  It may lie on a branch
that conflicts with \(B_\star\), and the validator may have lacked the
height-\((H{+}1)\) chain at vote time.  Lemma~\ref{lem:tsq-common-knowledge} puts
\(C_\star\), its ancestry, and the counted votes in every honest store by
\(d_{r_\star}+3\Delta\), strictly before the second slot boundary, so
every such validator's maximum rises to \(H+1\) and the
ungraded-root re-derivation
of Definition~\ref{def:rederivation} replaces its root at that
boundary.  The re-derivation state is derived under the activation
filter (Definition~\ref{def:activation-filter}), so it activates the same
finality outcomes as the round's action state; with maximum \(H+1\),
Definition~\ref{def:finality-root} selects the finalized root there,
the same root as in the action state, and
Lemma~\ref{lem:first-confirmation-simplex-roots} makes it an ancestor of
\(B_\star\).  Every honest root used
in the round's second and later slots is therefore compatible with
\(B_\star\), and \(B_\star\) lies in each candidate tree: the walks of
the second and later slots use the validator's current, unaged candidate
tree by Remark~\ref{rem:aged-scope}, the ancestry and leaf
witness of \(C_\star\) are in every honest store strictly before the
second slot boundary as above, and viability holds by
Lemma~\ref{lem:predrain-height-retention} --- every accepted height
increase between \(r_\star\) and the drain retains \(B_\star\)'s
viability --- together with \(B_\star\)'s own state-height.  Each such
voter's view contains its slot-view freeze view and is held by the
following evaluation's view bound, so Lemma~\ref{lem:tsq-adoption} and
Corollary~\ref{cor:tsq-walk} put every honest second-slot vote in
\(B_\star\)'s subtree; for the round's later slots and the drain round's
first slot, every preceding slot's honest votes lie in \(B_\star\)'s
subtree, so the induction of Lemma~\ref{lem:recovery-rooted-persistence}
keeps every later honest vote there as well.  Last, if a mid-round
sibling rise evicts a compatible root from a current tree while
\(B_\star\) stays viable, that root is necessarily a strict
descendant of \(B_\star\) --- an ancestor would inherit
\(B_\star\)'s own viability witness --- so the walk-standing rule of
Definition~\ref{def:walk-standing} has its walk stand inside
\(B_\star\)'s subtree, and the vote confinement above is unaffected.
\end{proof}

\begin{lemma}[The drain action confirms the first confirmation everywhere]
\label{lem:drain-confirms}
At the drain action, \(B_\star\) is available-confirmed at every
honest validator's evaluation, \(B_\star\) passes every guard of
Definition~\ref{def:official-confirmation}, and every honest deepest
official confirmation and nonempty SG head at the drain descends from
\(B_\star\).

In plain words, the drain action turns the one confirmation into a
universal official confirmation.
\end{lemma}

\begin{proof}
Lemmas~\ref{lem:predrain-height-retention}--\ref{lem:confirmation-round-slots}
establish the viability-or-finalized-past condition, the compatible-root
condition for every slot from the round's second slot onward, and the
confinement of every honest vote in those slots to \(B_\star\)'s
subtree.  Since every honest committee member of the drain
round's first slot votes for a block in \(B_\star\)'s subtree,
Lemma~\ref{lem:tsq-liveness} available-confirms \(B_\star\) at every
honest validator at the drain action's evaluation, and
Lemma~\ref{lem:tsq-monotonicity}, applied from that evaluation with the
established root and candidate-tree conditions, keeps every later honest
vote in \(B_\star\)'s subtree and every later available confirmation
compatible with \(B_\star\).  Lemma~\ref{lem:drain-stable-roots} gives
\(S_{u,\mathrm{act}}^{r_{\rm drain}}\preceq B_\star\), and the ordinary head
also descends from \(B_\star\).  For the veto-free guard: a vetoing
block would conflict with \(B_\star\), descend from the validator's
action-state finalized root --- which contains \(F_\star\) by
Lemma~\ref{lem:drain-stable-roots}, so the block descends from \(F_\star\) ---
and hold a grade~0 active at the drain's grade-0 freeze, exactly what
Lemma~\ref{lem:first-confirmation-clears} excludes: the finalization evidence, relayed as in
Lemma~\ref{lem:drain-stable-roots}, is processed at every honest
validator strictly before the drain round begins, hence before that freeze, so every conflicting grade is inactive
there, and the guard of
Definition~\ref{def:official-confirmation} tests freeze-active grades
only.
No honest validator takes the abstention branch at the drain: every
honest drain-round root is compatible with \(B_\star\) and lies in its action-state candidate tree by
Lemmas~\ref{lem:drain-stable-roots}
and~\ref{lem:confirmation-round-slots}; an accepted drain
proposal has a compatible proposed root --- graded roots have no
conflicting grade, and an ungraded root equals the receiver's own
selection --- and the root ancestry
\(S_{u,\mathrm{act}}\preceq S_{u,\mathrm{vote}}\preceq Z\prec
B_p\) holds by the acceptance guards together with
Lemma~\ref{lem:no-forward-move}.  If the proposal block \(B_p\) is
still in the action candidate tree, the admission accepts.  If the one
pre-drain rise permitted by Lemma~\ref{lem:predrain-height-retention}
--- from \(H+1\) to \(H+2\) --- has evicted \(B_p\), the eviction
case of Definition~\ref{def:action-root} treats the round as
proposal-free, and likewise a failed or absent proposal takes the
failed-proposal branch: in every such case the checked root is the
round's vote-time root, compatible with \(B_\star\) by
Lemma~\ref{lem:drain-stable-roots} and in the action candidate tree
because it is the action selection or a graded \(B_\star\)-side root
at the confirmed heights, which the same maximum bound keeps viable.
So the admission accepts in every case.  An ungraded re-derivation, if
any fires, lands on the action selection, ancestral by
Lemma~\ref{lem:predrain-simplex-roots}.
Thus \(B_\star\) passes every guard in
Definition~\ref{def:official-confirmation}, which makes \(B_\star\) itself an
official confirmation at every honest validator, so every honest deepest
official confirmation and SG head in that action descends from
\(B_\star\).
\end{proof}

\begin{lemma}[The first recovery confirmation remains canonical and viable]
\label{lem:first-confirmation-persists}
Let \(F_\star\) be the finalized root named in
Lemma~\ref{lem:first-confirmation-heads}.  From \(r_\star\) to
\(r_{\rm drain}\), every
honest action root admitted by the action-state check, every available
confirmation at an action's evaluation, and every nonempty SG head is either
compatible with \(B_\star\) or conflicts with a finalized root whose
evidence the confirming validator processed, from
Lemma~\ref{lem:first-confirmation-heads}.  From
\(r_{\rm drain}\) onward, every honest ordinary Goldfish head, raw vote,
deepest available-confirmed block, and nonempty SG head descends from
\(B_\star\); every available or official confirmation is compatible with
\(B_\star\); and every honest Simplex root and stable root is compatible
with it.  At every action, either \(B_\star\) is in the actor's store-global
viable subtree or the actor's finalized root descends from \(B_\star\).

In plain words, the first confirmation is already safe; the following SG
batch removes the stale vetoes and turns that safe prefix into the permanent
rooted path.
\end{lemma}

\begin{proof}
We separate the proof into three time intervals.

\emph{1. Before \(r_\star\).}
As shown in Lemma~\ref{lem:first-confirmation-simplex-roots}, the store-global maximum
cannot exceed \(H+1\): an earlier increase beyond that value would have an
honest confirmed witness at height at least \(H+1\), contradicting the
choice of \(r_\star\).  Thus, once \(B_\star\) is relayed, it is viable in
every honest store unless that store has finalized a descendant of
\(B_\star\).  Lemma~\ref{lem:first-confirmation-simplex-roots} also proves
directly that every action-state Simplex root at \(r_\star\) is an ancestor
of \(B_\star\).

\emph{2. From \(r_\star\) to \(r_{\rm drain}\).}
Lemma~\ref{lem:first-confirmation-heads} gives stable roots, available
confirmations, and SG heads that are compatible with \(B_\star\) or already
excluded by \(F_\star\).  Honest relay and
Lemma~\ref{lem:finalized-lockin} make the second class inert before the drain
action.
Lemma~\ref{lem:predrain-height-retention} keeps \(B_\star\) viable
through every accepted height increase,
Lemma~\ref{lem:predrain-simplex-roots} keeps every selected Simplex
root an ancestor of \(B_\star\),
Lemma~\ref{lem:drain-stable-roots} makes every drain stable root
compatible, Lemma~\ref{lem:confirmation-round-slots} confines the
round's later-slot roots and votes, and Lemma~\ref{lem:drain-confirms}
makes the drain action official-confirm \(B_\star\) at every honest
validator.  These five lemmas prove the interval's claims.

\emph{3. From \(r_{\rm drain}\) onward.}
We prove the stated claims together by induction over recovery actions.  At
the start of an action, assume that every honest ordinary head, raw vote,
deepest available-confirmed block, deepest official confirmation, and nonempty
SG head produced since
\(r_{\rm drain}\) descends from \(B_\star\), that every available or official
confirmation produced since \(r_{\rm drain}\) is compatible with \(B_\star\),
that every current Simplex and
stable root is compatible with \(B_\star\), and that each honest store either
keeps \(B_\star\) viable or has finalized a descendant of it.  Lemma~\ref{lem:drain-confirms} is the base case.
Lemma~\ref{lem:filtered-hmax-monotone} and the accepted height-\(H+1\) block
\(B_\star\) keep the store maximum at least \(H+1\) throughout this
induction.

First process the evidence admitted at the action cutoff.  For a new
justification, choose an honest target signer from its quorum.  A
justification at height at most \(H-1\) cannot be selected as the Simplex
root once the store maximum is at least \(H+1\).  A height-\(H\)
justification does not exist by Lemma~\ref{lem:no-height-h-justification}.  For a
justification at height at least \(H+1\), the honest signer's signature
cannot predate \(r_\star\).  If it was signed in \(r_\star\), compatibility
starts from the two alternatives in
Lemma~\ref{lem:first-confirmation-heads}.  Its official confirmation has
state-height at least \(H+1\), and every finalized root the signer
had processed by \(r_\star\) has height below \(H\)
(Lemma~\ref{lem:selectable-roots}, applied before the first
confirmation); Lemma~\ref{lem:past-finalized}, applied to each such
root, puts them all on the confirmation's chain and excludes the
finalized-conflicting alternative.  If it was signed later,
its deepest official confirmation, which is the signing gate, descends from
\(B_\star\) by the induction hypothesis.  Definition~\ref{def:height-vote-rule} makes the justified target
an ancestor of that confirmation, so the target is compatible with
\(B_\star\).  A new finalized root below \(H\) is an ancestor of
\(B_\star\) by Lemma~\ref{lem:past-finalized}; no finalized root exists at
\(H\); and a higher finalized root copies one of the justifications just
considered.  Thus every action-state Simplex root is compatible with
\(B_\star\).

Next process every branch-local height increase admitted at the same cutoff.
There are two cases.
\begin{enumerate}
 \item If the increase does not raise the store-global maximum, the
 viability threshold is unchanged.  The accepted \(B_\star\)-subtree keeps
 its previous witness and cannot be removed by
 Definition~\ref{def:finality-root}.
 \item If the increase raises the maximum from \(x\) to \(x+1\), then the
 old maximum is \(x\), so \(x\geq H+1\).  Its honest confirmed witness is
 supplied by Lemma~\ref{lem:height-increase-confirmation}.  The signature
 cannot predate \(r_\star\).  The preceding action cases make its official
 confirmation compatible with, or descending from, \(B_\star\).  The two
 ancestry cases of Lemma~\ref{lem:predrain-height-retention} ---
 \(B_\star\preceq Q\), or \(Q\preceq B_\star\) with state-height
 monotonicity --- give a height-\(x\) witness in the
 \(B_\star\)-subtree when the maximum rises to \(x+1\).
\end{enumerate}
Delivery precedes use by
Assumption~\ref{ass:frontier-opportunities}.  Consequently
\(B_\star\) remains viable unless the finalized root already descends from
it.  This establishes the height-filter premise before fork choice is run.

The eligible SG batch in the action descends from \(B_\star\) by the
induction hypothesis.  Lemma~\ref{lem:fresh-support-clears} excludes every
conflicting grade.  Together with the compatible action-state Simplex root,
the three stable-root cases therefore give a stable root compatible with
\(B_\star\).  Mid-round re-derivations preserve this:
Lemma~\ref{lem:rederivation-ancestry} applies with \(B:=B_\star\) ---
its justification premise holds by this induction's own Simplex-root
accounting, which keeps every selectable justification and every
activated finalized root compatible with
\(B_\star\) --- descendants of the healed block included --- and a
superseding rise is the separately handled
viability case --- so the roots of every later slot of each round stay
compatible with \(B_\star\).  If the finalized root already descends from \(B_\star\), every
accepted Goldfish block descends from it directly.  Otherwise the viability,
compatible-root and candidate-tree premises of
Lemma~\ref{lem:tsq-monotonicity} now hold, so every honest head and raw
vote continues to descend from \(B_\star\); a stable root evicted by
a taller sibling inside \(B_\star\)'s subtree is a strict descendant
of \(B_\star\), and the walk-standing rule keeps its walk inside
that subtree as well.

In the base case, \(B_\star\) is available-confirmed at every honest
validator.  Thereafter every honest vote descends from \(B_\star\), so
Lemma~\ref{lem:tsq-liveness} keeps \(B_\star\) available-confirmed at
every honest evaluation: the deepest available-confirmed block contains
\(B_\star\), and Lemma~\ref{lem:tsq-uniqueness} makes every other
available-confirmed block compatible with it.  If the action-state Simplex root is an ancestor of
\(B_\star\), then \(B_\star\) is an official-confirmation candidate.  If the
Simplex root descends from \(B_\star\), every official-confirmation candidate
and every fallback block between the Simplex root and action root descends
from \(B_\star\).  With no conflicting grade~0, the deepest candidate, or
the fallback when no candidate exists, makes every nonempty honest SG head and
every honest deepest official confirmation
descend from \(B_\star\).  Every official confirmation is an ancestor of the
deepest available-confirmed block by
Definition~\ref{def:official-confirmation}, so it is compatible with
\(B_\star\) as well.
This closes the SG, Goldfish, Simplex-root, and height-filter parts of the
induction.
\end{proof}

\begin{lemma}[An honest proposal extending the current finality roots is unobstructed]
\label{lem:post-healing-proposer}
After the drain action in
Lemma~\ref{lem:first-confirmation-persists}, consider a round with an
honest round proposer --- the common-recognition event of
Definition~\ref{def:recovery-timing} --- whose proposal has parent
\(Z\) and block \(B\).  Suppose every honest selection- and vote-state
Simplex root is an ancestor of \(Z\); at each of the round's
cutoffs, every honest validator's store maximum, active finality
outcomes, and processed finalized evidence equal the proposer's at that
same cutoff, and none of these inputs changes from the selection cutoff
through the action; every
action-state Simplex root is an
ancestor of \(B\); and the proposal path, with each receiver's root,
remains in each receiver's
store-global viable subtree at every per-round state through the
action.  Then the proposal is unobstructed.

In plain words, after healing a commonly recognized honest proposal
succeeds once it extends the finality roots that have already been
accepted.
\end{lemma}

\begin{proof}
The equal-and-constant-state premise gives every receiver
\(S_{u,\mathrm{vote}}^r=S_{p,\mathrm{sel}}^r\): the two states are
read at different cutoffs, but the premise makes their inputs equal
across validators at each cutoff and unchanged between the cutoffs, and
the selection of Definition~\ref{def:finality-root} is a function of
the activated pairs \((J,h_j)\), \((F,h_F)\), and the store maximum
alone.  Likewise
Lemma~\ref{lem:aged-containment} supplies the candidate-tree
containment clause of Definition~\ref{def:usable-honest-round}.
By Lemma~\ref{lem:first-confirmation-persists}, every eligible honest
nonempty SG head in the preceding batch descends from \(B_\star\), so by
Lemma~\ref{lem:root-no-regression} no graded root conflicts with
\(B_\star\).  If the proposer holds grade~2 for some root --- in
particular, when every honest validator emitted a nonempty head, honest
weight \(W-b\geq m\) guarantees grade~2 for \(B_\star\) itself --- that
root is compatible with \(B_\star\) and the graded branch proceeds.
Otherwise --- whether the batch carries few honest heads or none --- no
grade-2 root exists, and the proposer supplies
its selection-state Simplex root \(A_p^r=S_{p,\mathrm{sel}}^r\) as an
ungraded base root, relaying its backing
certificate as ordinary evidence; the selection-agreement premise gives
every receiver \(S_{u,\mathrm{vote}}^r=S_{p,\mathrm{sel}}^r=A_p^r\), and
the equality disjunct of item~2 in Definition~\ref{def:stable-root}
accepts.

Lemma~\ref{lem:g2-g1} supplies every receiver's upper check for the
proposer's deepest grade-2 root.  Lemma~\ref{lem:g3-g2},
Lemma~\ref{lem:direct-roots-chain}, and deepest grade-2 selection put every
receiver's grade-3 lower root below the proposed root.  The Simplex-root
premises and Definition~\ref{def:stable-root} move each receiver only forward
on the chain ending at \(Z\) at vote time and at \(B\) at action time.  The
viability premise supplies the remaining candidate-tree guards.
Recognition, the first conjunct of
Definition~\ref{def:usable-honest-round}, holds by
Lemma~\ref{lem:common-recognition} from the honest-round-proposer
premise.  These facts satisfy every guard in
Definition~\ref{def:usable-honest-round}, so the proposal is
unobstructed.
\end{proof}

\begin{corollary}[After the drain, an honest proposal reaches every official confirmation]
\label{cor:post-drain-confirms}
After the drain action, let \(B\) be the proposal of an honest-proposer
round satisfying Lemma~\ref{lem:post-healing-proposer}.  Then one of the
following holds: the
store-global maximum rises above \(\sigma[B].h+1\) during the round ---
a height advance that supersedes the proposal; or a justification on a
\(B_\star\)-descendant branch conflicting with \(B\), at a height
admitting selection, activates before the closing action and redirects
closing selections --- by Lemma~\ref{lem:target-uniqueness} at most one
such justification exists per height, so this alternative occurs at most
once per context height, and the redirected root still descends from
\(B_\star\); or at that
round's closing action every honest validator has an official
confirmation, and its deepest official confirmation contains \(B\).

In plain words, healed rounds behave like ordinary rounds: an honest
proposal is confirmed everywhere unless the chain outruns it or one
late sibling justification redirects it, and both exceptions are
bounded.
\end{corollary}

\begin{proof}
The proposal is unobstructed by Lemma~\ref{lem:post-healing-proposer}, so
Lemma~\ref{lem:honest-proposer-votes} gives the exact honest vote batch;
Lemma~\ref{lem:recovery-rooted-persistence}, whose root and candidate-tree
premises hold by Lemma~\ref{lem:first-confirmation-persists}, keeps every
honest head and vote in \(B\)'s subtree through the round; and
Lemma~\ref{lem:recovery-rooted-confirmation} available-confirms \(B\) at
every honest validator.  The confirmation mechanics are the shared
Lemma~\ref{lem:gate-closure}; its premises are established below for
the case in which neither post-drain exception occurs.
Lemma~\ref{lem:honest-proposal-confirms}'s pre-confirmation case space
--- a height-filter removal at the recovery heights, a conflicting
Simplex root off \(B_\star\), and the unredeemed
height-\((H{-}1)\) justification --- is excluded after the drain by
Lemma~\ref{lem:first-confirmation-persists}; the post-drain case space
is the following two alternatives.  A sibling
\(B_\star\)-descendant branch may raise the store maximum during the
round and remove \(B\)'s subtree from the viability band before the
closing action: the first alternative.  Or a justification on a sibling
\(B_\star\)-descendant branch, assembled from pairs included after
\(a_{r-1}\) and accepted by \(a_r\), may activate for the closing
round and satisfy the selection guard there, redirecting closing
Simplex roots off \(B\): the second alternative --- persistence
constrains such a justification only to \(B_\star\)-compatibility, not
to \(B\)'s chain, so it cannot be excluded, and
Lemma~\ref{lem:target-uniqueness} bounds it to one per height.  When
neither occurs, every justification selectable within the round and its
closing action is an ancestor of \(B\), the candidate-tree premises
hold through the closing action, and
Lemma~\ref{lem:rederivation-ancestry} covers every mid-round
re-derivation, so the roots used by
the round's later slots stay compatible with \(B\).  These facts, with
the available confirmation above, the compatible grades and processed
finality of Lemma~\ref{lem:first-confirmation-persists}, and admission
accepting as at the drain action, supply
Lemma~\ref{lem:gate-closure}'s premises, and that lemma gives every
honest validator an official confirmation whose deepest contains
\(B\).
\end{proof}

\begin{lemma}[A saved target transfers to another chain that contains it]
\label{lem:saved-target-transfers}
Let an honest validator have saved nonempty target \(T_0\) at height \(k\).
Every chain that is still at height \(k\) and contains \(T_0\) has
current-height target \(T_0\).

In plain words, the first block of one height has the same name on every
extension that is still at that height.
\end{lemma}

\begin{proof}
The signing rule saved \(T_0\) only when it was the first block of height
\(k\) on the signer's chain.  If another height-\(k\) chain contains
\(T_0\), Lemma~\ref{lem:chain-target-uniqueness} makes \(T_0\) its first
block of height \(k\) as well.  Definition~\ref{def:current-height-target}
therefore gives that chain the same target.
\end{proof}

\begin{lemma}[A lock is an ancestor of the proposal or reveals one old justification]
\label{lem:lock-alignment}
Let \(B\) be a valid source proposal with post-state height
\(k\in\{H-1,H,H+1\}\), and suppose an honest validator officially
confirms a descendant of \(B\).  Call a finality lock \(T_L\)
\emph{off-\(B\)} when \(T_L\) is not an ancestor of \(B\).  If that
validator has a finality lock \(T_L\) at
height \(k\), formed before the first recovery confirmation, then either \(T_L\preceq B\), or \(k=H-1\) and the lock reveals
the unique justification \(\mathsf{JC}(H-1,T_L)\) on another branch.  The
second case is impossible at \(k=H\) and at \(k=H+1\).

In plain words, the only lock that can stop progress on a confirmed recovery
proposal is the one old justification immediately below the nonjustifiable
height.
\end{lemma}

\begin{proof}
Definitions~\ref{def:height-vote-rule} and~\ref{def:finality-vote-rule} write
the lock only after the validator has signed target \(T_L\) and learned
\(\mathsf{JC}(k,T_L)\).  If \(T_L\) lies on the chain ending at the official
confirmation, two cases arise.  While that chain is still at height
\(k\), Lemma~\ref{lem:saved-target-transfers} makes \(T_L\) the
chain's height-\(k\) target.  If the chain has advanced past \(k\), its
height-\(k\) target is its unique first height-\(k\) block by
Lemma~\ref{lem:chain-target-uniqueness}; the source proposal \(B\) lies
on that chain with \(\sigma[B].h=k\), so the chain's height-\(k\)
transition occurred at or before \(B\) and that first block is an
ancestor of \(B\), and a target saved at height \(k\) on the chain is
that same first block.  In both cases \(T_L\) equals the source
proposal's target
and is an ancestor of \(B\).

If \(T_L\) does not lie on that chain, the retained signed justification is
the stated off-branch witness.  Lemma~\ref{lem:no-height-h-justification} excludes a
height-\(H\) justification, and a height-\((H{+}1)\) justification would
carry an honest official height-\((H{+}1)\) confirmation by
Lemma~\ref{lem:height-vote-confirmation}, so before the first recovery
confirmation no lock exists at \(H+1\) at all.  Therefore \(k=H-1\).
Lemma~\ref{lem:target-uniqueness} makes this off-branch justification unique.
\end{proof}

\begin{lemma}[A confirmed source proposal gets a progress quorum]
\label{lem:common-progress}
Let \(B\) be the source proposal of an action, let
\(\operatorname{ctx}(B)=(k,T_B,\nu_B)\) with
\(k\in\{H-1,H,H+1\}\), and suppose every honest validator officially confirms
a descendant of \(B\).  Suppose no honest validator has an off-\(B\) finality
lock at \(k\).  Then one of the following holds:
\begin{enumerate}
 \item some valid descendant of \(B\) has already advanced above height \(k\);
 \item every honest validator uses \(\operatorname{ctx}(B)\), and every honest
 current-height pair is nonempty and is either \((k,\bot)\) or
 \((k,T_i)\) with \(T_i\preceq B\).  If one valid strict descendant
\(D\succ B\),
 whose prestate is still at height \(k\), includes those pairs, they set
 progress bits of total weight at least \(W-b\geq q\).
\end{enumerate}

In plain words, validators may name different old targets, but every such
target is already on the proposal chain and therefore helps that branch
advance.
\end{lemma}

\begin{proof}
Fix one honest validator and its deepest official confirmation
\(Q_i^r\), the block whose authenticated state its action signs.
If \(\bar\sigma_i.h>k\), its official confirmation contains \(B\) and lies on
the same walk chain, so \(Q_i^r\) is a valid descendant of \(B\) that has
already advanced above \(k\).  This is item~1.  Otherwise
\(\bar\sigma_i.h=k=\sigma[B].h\).  If this equality holds for every honest
validator, Definition~\ref{def:ordinary-current-target} makes each of them
use the same height, target, and nonjustifiable flag from \(\sigma[B]\).
We prove item~2 by fixing one such validator.  There are five cases, in the
order used by
Definition~\ref{def:height-vote-rule}.
\begin{enumerate}
 \item With timeout history, the validator repeats \((k,\bot)\).  Its official
   confirmation contains \(B\), whose post-state has height \(k\), so the
   confirmation-height guard permits the repeat.
 \item With a finality lock, the premise and
   Lemma~\ref{lem:lock-alignment} give \(T_L\preceq B\), so the validator
   repeats \((k,T_L)\).
 \item With an unlocked saved target, the validator repeats it if it lies on
   its official-confirmation chain.  Since that chain contains \(B\) and is
   still at height \(k\), Lemma~\ref{lem:saved-target-transfers} makes the
   saved target equal to \(T_B\), hence ancestral to \(B\).  If the saved
   target is off that chain, the rule signs timeout
   because the confirmation has reached \(k\).  This switch is not slashable,
   by Lemma~\ref{lem:signer-safety}.
 \item With no history and \(\nu_B=\mathsf{true}\), the rule signs timeout.
 \item With no history and \(\nu_B=\mathsf{false}\), the rule signs
   \((k,T_B)\), and \(T_B\neq\bot\): under this lemma's premises every
   honest validator officially confirms a descendant of \(B\) and its
   closed confirmation state \(\bar\sigma_i\) is at height
   \(k=\sigma[B].h\), so
   Definition~\ref{def:ordinary-current-target} takes its first branch,
   \(Y_i=B\) and \(\sigma_i=\sigma_a[B]\), the closed state the
   definition assigns.  Here \(B.\slot<\slot(a)\): \(B\) is the preceding round's
   first-slot proposal and \(\slot(a)\) is the current round's second
   slot, so the closure makes at least one call.  The closure's first
   \textsc{process\_slot} call
   materializes the height-\(k\) target when it was still pending ---
   \(T_k\gets B\) fires because \(B.\slot\geq\sigma[B].s_k\), the
   height-\(k\) transition having been consumed at a block at slot
   \(\sigma[B].s_k\leq B.\slot\) (Lemma~\ref{lem:empty-slot-noop}) ---
   so \(\sigma_a[B].T_k\) is nonempty and the empty-target fallback of
   the definition cannot arise here.
   Hence \(T_B=\sigma_a[B].T_k\), nonempty,
   and the definition of proposal context with
   Lemma~\ref{lem:same-slot-target} gives \(T_B\preceq B\).
\end{enumerate}
These cases exhaust the signing rule and prove the first claim.  On the
including branch,
\[
 T_i\preceq B\preceq D.\parent .
\]
Definition~\ref{def:vote-contribution} therefore sets the progress bit for
every included honest pair, including every exact target and timeout.
Honest weight is \(W-b\geq q\) by
Lemma~\ref{lem:integer-thresholds}.
\end{proof}

\begin{lemma}[A confirmed source proposal fixes one height context]
\label{lem:usable-common-action}
Let \(B\) be an unobstructed honest proposal.  In the next action, which
grades \(B\)'s SG batch, suppose every honest official confirmation contains
\(B\) and every honest validator's closed confirmation state
\(\bar\sigma_i=\sigma_a[Q_i^r]\) is
still at height \(\sigma[B].h\).  Then \(B\)
is that action's source proposal and every honest
validator uses the same authenticated context
\(\operatorname{ctx}(B)=(k_B,T_B,\nu_B)\), even when their confirmed tips are
different descendants of \(B\).  That later action may have processed
additional blocks from both slots of its round, so the condition that every
closed confirmation state \(\bar\sigma_i\) is still at height \(\sigma[B].h\) is an explicit
premise and must be
checked before this lemma is invoked.

In plain words, the proof needs one common proposal block, not one common
descendant tip.
\end{lemma}

\begin{proof}
The source-proposal clause of
Definition~\ref{def:ordinary-current-target} carries \(B\) through exactly
the action that grades its SG batch.  Valid-block replay gives one immutable
post-state \(\sigma[B]\).  The same definition derives \(k_B\), \(T_B\), and
\(\nu_B\) only from that state and \(B\)'s slot.  Different descendants of
\(B\) therefore cannot change the triple.
\end{proof}

\begin{lemma}[A proposal-descendant inclusion advances one height]
\label{lem:common-action-advances}
Under the delivery and inclusion clauses of
Assumption~\ref{ass:frontier-opportunities}, if item~2 of
Lemma~\ref{lem:common-progress} holds and its pairs are included in a valid descendant of \(B\)
whose prestate is still at height \(k\), the direct height check advances
that branch from \(k\) to \(k+1\).  If the branch has already advanced, the
claim is already true for that branch.  A certificate on a sibling branch is
not used for either conclusion.

In plain words, progress is proved on the branch that includes the votes, not
in a store merely because another branch contains a certificate.
\end{lemma}

\begin{proof}
Item~2 of Lemma~\ref{lem:common-progress} sets progress bits of weight at least \(q\).
If the exact target bits also reach \(q\) at an ordinary height,
Definition~\ref{def:height-outcome} processes that branch first.  Otherwise
the progress branch applies.  Both branches call
\textsc{advance\_height} exactly once and reset both participation arrays.
\end{proof}

\begin{lemma}[A selected branch one height behind catches up]
\label{lem:height-catchup}
Suppose the store-global maximum is \(H\) and an unobstructed honest proposal
\(B\) has post-state height \(H-1\).  Its closing action signs the
height-\(H-1\) pairs.  By the including honest proposal that the
inclusion clause of Assumption~\ref{ass:frontier-opportunities}
supplies on a descendant of \(B\), and, only if needed, one retry after the
unique height-\(H-1\) justification is first admitted and redirects the
selection --- the separately charged redirection event --- either that
descendant advances to
height \(H\), the permitted heights move above \(H\), or the first recovery
confirmation has already occurred.

In plain words, validators need not begin on the maximum-height branch, but
a sibling certificate is not treated as advancing their selected branch.
\end{lemma}

\begin{proof}
Apply Lemma~\ref{lem:honest-proposal-confirms}.  In item~1,
Definition~\ref{def:ordinary-current-target} makes \(B\) the next action's
source proposal.  If no off-\(B\) finality lock exists,
apply Lemma~\ref{lem:common-progress}.  If some honest closed confirmation state \(\bar\sigma_i\) has
height above \(H-1\), its first item says that a valid
descendant of \(B\) is already at height \(H\).  Otherwise every honest action
state is still at \(\sigma[B].h=H-1\), so
Lemma~\ref{lem:usable-common-action} supplies the common context and item~2
of Lemma~\ref{lem:common-progress} applies;
Lemma~\ref{lem:common-action-advances} then advances the descendant that
includes the honest pairs from \(H-1\) to \(H\).

No off-\(B\) lock at \(H-1\) can suppress a counted action's height
pair.  A finality lock at \((H-1,T_L)\) is recorded only when its holder
signs the finality pair, which
Definition~\ref{def:finality-vote-rule} reads from the holder's
fork-choice action state: the lock therefore postdates an action at
which \(\mathsf{JC}(H-1,T_L)\) was the holder's \emph{activated} latest
justification.  If that activation predates the window, the holder
relayed the certificate under Definition~\ref{def:earlier-votes} before
the first counted round, and it is activated at every honest validator
from the first counted round on.  If the certificate is instead first
admitted during the window, its first admission is the unique
redirection charged as case~3 of
Lemma~\ref{lem:honest-proposer-coverage}, no lock exists before the
holder's first action with that activation, and a lock recorded at
action \(a_r\) constrains that action's own height pair --- the
ordering of Definition~\ref{def:fg-rule} --- and every later one; the
certificate, admitted at the holder by \(a_{r-1}\) and relayed, is
activated at every honest validator by the following action.  In either case, the round at which the certificate first activates at
any honest validator is the charged redirection round: its action may
see suppressed height pairs at the activating validators --- a lock
recorded there constrains that same action --- while the one-round
stagger leaves others not yet active.  From the following round on,
every honest selection state has the
certificate active; with store maximum \(H\) --- this lemma's standing
hypothesis, and if the maximum rises instead, the item-2 alternative
below ends this case --- the guard
\(\hmax=h_j+1\) selects \(T_L\) everywhere, every honest proposal
extends \(T_L\), and its context satisfies
\(T_L\preceq B\) --- so any lock an honest validator then holds is
on-\(B\), and no counted action after the charged round is
suppressed within this lemma's maximum-\(H\) span.  The lock case of Lemma~\ref{lem:common-progress} is
therefore vacuous at counted rounds, and the height-\((H-1)\)
redirection remains the
separately charged selection event.  If item~2 of
Lemma~\ref{lem:honest-proposal-confirms} occurs, the global maximum has moved
to at least \(H+1\), so every later viable proposal already has post-state
height at least \(H\).  This is a band change, not a transition applied to
\(B\)'s branch.  Item~4 is the first recovery confirmation.  These cases
prove the claim.
\end{proof}

\begin{lemma}[Recovery reaches the successor of the nonjustifiable height]
\label{lem:leave-nonjustifiable-height}
Let \(B\) be an unobstructed honest proposal with post-state height \(H\).
Its closing action signs the height-\(H\) pairs.  By the including
honest proposal that the inclusion clause of
Assumption~\ref{ass:frontier-opportunities} supplies on a descendant of
\(B\), and, only if needed,
one retry after the unique height-\(H-1\) justification is relayed, either
that descendant advances to \(H+1\) or the first recovery confirmation has
already occurred.

In plain words, pre-existing votes at \(H\) may help or be repeated; they do
not prevent the timeout advance from the nonjustifiable height.
\end{lemma}

\begin{proof}
Apply Lemma~\ref{lem:honest-proposal-confirms}.  In item~1,
Definition~\ref{def:ordinary-current-target} makes \(B\) the source proposal.
Lemma~\ref{lem:no-height-h-justification} excludes both a height-\(H\)
justification and an off-\(B\) height-\(H\) finality lock.
Apply Lemma~\ref{lem:common-progress}.  If some honest closed confirmation state \(\bar\sigma_i\) has
height above \(H\), its first item says that a valid
descendant has already advanced to \(H+1\).  Otherwise every honest action
state remains at \(H\), so Lemma~\ref{lem:usable-common-action} fixes
\(\operatorname{ctx}(B)\) at height \(H\), and under item~2 of
Lemma~\ref{lem:common-progress} the
honest pairs have progress weight at least \(q\), including every
saved-target and timeout case.  Since \(H\) is nonjustifiable,
Definition~\ref{def:height-outcome} disables its target branch and
Lemma~\ref{lem:common-action-advances} advances the including descendant by
progress.

Item~3 is the unique height-\(H-1\) redirection.  After its carrying state is
relayed, the next honest proposal extends that root and the case cannot
recur.  Item~4 is the stated early confirmation.  Item~2 is inapplicable to
a height-\(H\) proposal: removal from the permitted \(h_{\max}-1\) or
\(h_{\max}\) heights would require a state
at \(H+2\), whose honest height-\(H+1\) confirmed witness is already item~4.
\end{proof}

\begin{lemma}[An honest proposal at the recovery height forces the first confirmation]
\label{lem:honest-proposer-forces-first-confirmation}
If the first recovery confirmation has not occurred and an honest round
proposer produces an unobstructed block \(B\) with post-state height \(H+1\),
the next action, which grades \(B\)'s SG batch, produces that first
confirmation.

In plain words, an honest recovery-height proposal forces liveness; the first
confirmation remains the proposer-independent safety event.
\end{lemma}

\begin{proof}
If an accepted state had already advanced beyond \(H+1\), its advancing
quorum would contain an honest confirmed witness at \(H+1\) by
Lemma~\ref{lem:height-increase-confirmation}, contradicting the premise.
Apply Lemma~\ref{lem:honest-proposal-confirms}.  Its items~2 and~3 are
inapplicable at height \(H+1\): item~3 is a selection event, and the
selection guard \(\hmax=h_j+1\) of Definition~\ref{def:finality-root}
fails for the height-\((H{-}1)\) justification once every honest maximum
is at least \(H+1\), while item~2 would require a state at \(H+2\) as in
the preceding lemma; no proposal can reintroduce either, because item~2
of Definition~\ref{def:stable-root} accepts an ungraded proposed root
only when it equals the receiver's own selection.  Item~4 is already the
conclusion.
In the remaining case, Lemma~\ref{lem:closing-action-existence} produces
an official confirmation at every honest validator at the closing action,
whose deepest official confirmation contains \(B\).
The block \(B\) has post-state height \(H+1\), so these confirmations
satisfy Definition~\ref{def:first-recovery-confirmation}.
\end{proof}

\begin{lemma}[The recovery height advances after the first confirmation]
\label{lem:recovery-height-justifies}
After the drain action of
Lemma~\ref{lem:first-confirmation-persists}, a valid descendant of
\(B_\star\) advances from \(H+1\) to \(H+2\) once it includes the honest
height-\(H+1\) pairs, unless it had already advanced.  The advance may be by
an exact target quorum or only by progress.

In plain words, an old off-branch target may turn into a timeout, but every
honest pair still helps the healed branch leave \(H+1\).
\end{lemma}

\begin{proof}
If the including descendant has already advanced from \(H+1\), the lemma is
complete.  It remains to consider a descendant whose prestate is still at
height \(H+1\) after the drain action.

There was no honest height pair above
\(H\) before the relay point by
Assumption~\ref{ass:frontier-opportunities}.  Between that point and
\(r_\star\), an honest height-\(H+1\) pair would require an official
confirmation at height \(H+1\) by
Lemma~\ref{lem:height-vote-confirmation}, contradicting the definition of
\(r_\star\).  Thus no honest validator has height-\(H+1\) history before
action \(r_\star\), which first evaluates the first slot's available
confirmation and then releases its SG/FG attestation.

Lemma~\ref{lem:first-confirmation-persists} gives every honest validator a
nonempty deepest official confirmation on the \(B_\star\) chain, with no
conflicting active grade~0 in the drain action.  If some valid descendant
of \(B_\star\) has already advanced above \(H+1\), the lemma's
conclusion already holds for that descendant.  Otherwise no
\(B_\star\)-descendant has advanced, so every honest deepest official
confirmation --- a \(B_\star\)-descendant --- has closed state at
height exactly \(H+1\), and every honest context height is
\(k_i=H+1\).  For each honest validator, the signing rule then has
five relevant cases; the empty-history case with \(\nu=\mathsf{true}\)
cannot arise at an ordinary restart height.
\begin{enumerate}
 \item A prior height-\(H+1\) timeout is repeated.
 \item A saved height-\(H+1\) target on the \(B_\star\) chain is that chain's
   height-\(H+1\) target, by Lemma~\ref{lem:saved-target-transfers}, and is
   repeated.
 \item A saved target off the \(B_\star\) chain is unlocked, because no
   height-\(H+1\) justification existed before \(r_\star\), by
   Lemma~\ref{lem:no-hidden-h-advance} and the preceding paragraph; it changes
   to timeout.  The switch creates no E1 evidence, by
   Lemma~\ref{lem:signer-safety}, because no finality entry exists at that
   height.
 \item A validator with no history signs the chain's current target.
 \item A finality lock at \(H+1\), created at or after the drain when
   a block includes the released height-\((H{+}1)\) pairs and the
   resulting justification activates, locks the first
   height-\((H{+}1)\) block of \(B_\star\)'s chain, an ancestor of
   the including descendant, so the validator repeats that pair and the
   progress bit is still set.
\end{enumerate}
Thus every honest pair is either a timeout or names a target ancestral to
the including descendant.  Definition~\ref{def:vote-contribution} sets every
honest progress bit.  Their weight is \(W-b\geq q\), so the direct height
check in Definition~\ref{def:height-outcome} advances to \(H+2\).  If the
exact target bits also reach \(q\), the target branch runs first and creates
the justification; otherwise the progress branch makes the same height
advance without one.
\end{proof}

\begin{definition}[Ordinary-height target opportunity]
\label{def:fresh-ordinary-action}
After action \(r_\star\), let \(h_0\) be at least every height at which an
honest validator signed a current-height pair before the recovery relay
point.  A later action is
\emph{an ordinary-height target opportunity} at height
\(h_\star>h_0\) when all of the following hold:
\begin{enumerate}
 \item it is the closing action of an honest-proposer round satisfying
   the premises of Lemma~\ref{lem:post-healing-proposer} --- including
   the per-cutoff cross-validator equality and intra-round constancy of
   store maxima, active finality outcomes, and processed finalized
   evidence, assumed with this opportunity --- and the round is not
   superseded by a store-maximum rise;
 \item the round's proposal \(B\) has
   \(\operatorname{ctx}(B)=(h_\star,T_\star,\mathsf{false})\);
 \item every honest validator's durable history at \(h_\star\) is
   either empty or contains only the same saved target \(T_\star\),
   with no timeout and no finality lock;
 \item every honest validator's closed confirmation state
   \(\bar\sigma_i\) at that closing action is still at height
   \(h_\star\); and
 \item no block on the branch that first includes the resulting target
   pairs advances past \(h_\star\) before including every honest one
   of them --- no progress-only quorum moves the confirmed branch past
   \(h_\star\) from the closing action until the height-\(h_\star\)
   tallies are processed with the full honest target weight.  This
   extends the still-at-height condition that
   Lemma~\ref{lem:usable-common-action} requires be checked before it
   is invoked through the inclusion itself.
\end{enumerate}  Confirmation coverage
is not part of this definition: Corollary~\ref{cor:post-drain-confirms}
derives from unobstructedness that every honest official confirmation at
this action contains \(B\), unless one of that corollary's two bounded
exceptions --- a superseding maximum rise or a late sibling
justification redirect --- occurs; an opportunity in the sense of this
definition is one at which neither does, and the recurrence supplied by
Assumption~\ref{ass:clean-target-opportunity} covers the wait.  The including
descendant is supplied by the inclusion clause of
Assumption~\ref{ass:frontier-opportunities} and is not part of this
definition.  Since \(K\geq2\), consecutive heights include a
non-periodic one, so \(\nu_B=\mathsf{false}\) restricts the choice of
\(h_\star\), not the execution.
From these first target signatures until
\(\mathsf{JC}(h_\star,T_\star)\) has been processed by every honest
validator, no honest validator signs a timeout at \(h_\star\).

In plain words, an earlier action may already have saved the same target, but
no honest validator has chosen a timeout or another target.
\end{definition}

\begin{assumption}[Clean target and finality inclusion]
\label{ass:clean-target-opportunity}
Once honest SG heads and raw Goldfish votes remain in one subtree, then after
any finite target, timeout, and finality histories there is, within
\(\rho\) rounds, an ordinary-height target opportunity in the sense of
Definition~\ref{def:fresh-ordinary-action}.  After the resulting
justification is common and all honest validators sign its finality pair, an
honest slot proposer includes all those finality pairs, within
\(\rho\) rounds of the justification becoming common, in a block
whose prestate still names that justification as the latest unfinalized
one.
Until that inclusion, every honest validator's fork-choice action state
--- \(\Sigma_{i,\mathrm{act}}^r\), the state
Definition~\ref{def:finality-vote-rule} reads --- has
processed \(\mathsf{JC}(h_\star,T_\star)\) and names it as its latest
unfinalized justification.  No
earlier direct height check on that inclusion branch installs a newer
justification, whether from one block's attestations or from participation
accumulated across several blocks.

This is a condition on the proposer sequence after the chain has healed.  It
is not implied merely by an honest proposer in every \(\rho\)-round
interval: earlier actions at a new height can leave different honest
validators with different saved targets or timeouts.  Seven parts of this
assumption are genuinely assumed:
\begin{enumerate}
 \item the absence, before the counted round's closing action, of a
   late sibling-justification redirect --- the second exception of
   Corollary~\ref{cor:post-drain-confirms};
 \item the still-at-height condition of
   Definition~\ref{def:fresh-ordinary-action} --- that no progress-only
   quorum advances the confirmed branch past the restart height from the
   closing action until the including block processes the full honest
   target weight;
 \item the occurrence, within \(\rho\) rounds, of an
   honest-proposer round at the restart height satisfying
   Lemma~\ref{lem:post-healing-proposer}'s premises ---
   unobstructedness with the cross-validator equality conditions ---
   whose closing action is not superseded by a maximum rise;
 \item the history alignment there;
 \item the no-timeout clause of
   Definition~\ref{def:fresh-ordinary-action}, which constrains honest
   actions after the opportunity until the justification is common;
 \item the common latest-justification action states through the
   inclusion; and
 \item the finality-pair inclusion ordering of the previous paragraph.
\end{enumerate}
Given item~1, confirmation coverage and the ordinary
including descendant are derived, not assumed ---
Corollary~\ref{cor:post-drain-confirms} and the inclusion clause of
Assumption~\ref{ass:frontier-opportunities} supply them.

In plain words, ordinary finality restarts at a later height where honest
histories are already aligned, and finality votes are included before a
newer justification replaces their subject in the including state.
\end{assumption}

\begin{lemma}[An ordinary-height target opportunity justifies its target]
\label{lem:fresh-ordinary-justify}
An ordinary-height target opportunity creates
\(\mathsf{JC}(h_\star,T_\star)\).

In plain words, all honest validators choose or repeat the same target and
the including block obtains a target quorum.
\end{lemma}

\begin{proof}
The still-at-height clause of
Definition~\ref{def:fresh-ordinary-action} puts every honest signed
context state at height \(h_\star\), so Definition~\ref{def:ordinary-current-target}
gives every honest validator \(k_i=h_\star\).  For an empty history,
Definition~\ref{def:height-vote-rule} takes its ordinary-height case and
signs \((h_\star,T_\star)\).  For a saved-target history, the opportunity
premise says that the saved value is \(T_\star\); because every official
confirmation contains \(B\) and \(T_\star\preceq B\), the same rule repeats
that pair.  The premise excludes the timeout and finality-lock cases
in the durable history at the action's start, and the action itself
writes no lock at \(h_\star\): the finality pair it evaluates first
concerns the latest unfinalized justification, whose height is below
\(h_\star\) because \(\mathsf{JC}(h_\star,T_\star)\) does not
yet exist, so the current-height rule reads no same-action lock at
\(h_\star\).
Thus every honest validator signs the exact pair
\((h_\star,T_\star)\).  Their weight is \(W-b\geq q\) by
Lemma~\ref{lem:integer-thresholds}.  The extended still-at-height clause
of Definition~\ref{def:fresh-ordinary-action} keeps the inclusion branch
at height \(h_\star\) until a block processes the height-\(h_\star\)
tallies with every honest target pair --- in particular, no block of that
branch advances first on a progress-only subset of the pairs --- and the
inclusion clause of
Assumption~\ref{ass:frontier-opportunities} supplies that block.  Every
included vote sets both
participation bits under Definition~\ref{def:vote-contribution}, the
target weight reaches \(q\), and at the
direct height check the target branch is checked before
progress and creates
\(\mathsf{JC}(h_\star,T_\star)\).
\end{proof}

\begin{lemma}[An honest inclusion finalizes the latest aligned justification]
\label{lem:honest-proposer-finalizes}
Suppose one SG/FG action ends with every honest validator holding the same
latest unfinalized justification \(\mathsf{JC}(h,T)\), with saved target \(T\)
and no timeout at \(h\).  Then every honest validator signs finality pair
\((h,T)\).  If an honest block proposer includes those pairs before any block
changes the latest justification, that block finalizes \(T\) before
processing any simultaneous target or progress quorum.

In plain words, once all honest validators can finalize the same latest
justification, an honest inclusion finalizes it even if the same block
also advances a later height.
\end{lemma}

\begin{proof}
The common action states and histories satisfy every guard in
Definition~\ref{def:finality-vote-rule}, so honest finality-pair weight is
\(W-b\geq q\).  The inclusion premise gives the honest block a prestate that
still names \((h,T)\) as latest.
Definition~\ref{def:block-attestation-processing} therefore puts every
included honest signer in \(P\).  Definition~\ref{def:height-outcome} checks
finality before target and progress, and finalizes \(T\).
\end{proof}

\begin{lemma}[Honest proposers restart finality after healing]
\label{lem:post-healing-finality}
After the drain action in
Lemma~\ref{lem:first-confirmation-persists},
Assumption~\ref{ass:clean-target-opportunity} creates and finalizes a
justification at a later ordinary height.

In plain words, once SG and the available chain stay in
\(B_\star\)'s subtree,
finality restarts at the next assumed height where honest target histories
are already aligned.
\end{lemma}

\begin{proof}
Lemma~\ref{lem:recovery-height-justifies} first gives a valid
\(B_\star\)-descendant at height \(H+2\).
Lemma~\ref{lem:first-confirmation-persists} keeps every honest SG head and
raw Goldfish vote in the persistent subtree of \(B_\star\) from the drain
onward.  Its induction also keeps later Simplex roots and stable roots
compatible with that subtree.

Lemma~\ref{lem:post-healing-proposer} makes an honest round proposal on
that subtree unobstructed whenever its premises hold, and
Corollary~\ref{cor:post-drain-confirms} then confirms it everywhere,
unless one of its two bounded exceptions occurs.  The
activation rule of Definition~\ref{def:activation-filter} confines
activation staggering to one round, and a staggered round has a finality
advance in flight, which is progress.  A round can also be obstructed by a
store-maximum divergence at its cutoff; as in
Lemma~\ref{lem:honest-proposer-coverage}, such a divergence means a
maximum rise is in flight, common one round later, and each such round
advances the height.  A round can finally be redirected by a late
sibling justification at its closing action ---
Corollary~\ref{cor:post-drain-confirms}'s second alternative --- at most
once per height by Lemma~\ref{lem:target-uniqueness}; once that
justification is active everywhere, the next honest proposal extends its
root, which descends from \(B_\star\), so the pipeline continues on
that branch.  Recurrence of the aligned unobstructed action is
exactly what Assumption~\ref{ass:clean-target-opportunity} postulates;
proposer fairness alone does not supply it, and this proof does not claim
otherwise.
Assumption~\ref{ass:frontier-opportunities} supplies an honest proposer
within \(\rho\) rounds, and its inclusion clause supplies the including
descendant.  What Assumption~\ref{ass:clean-target-opportunity} adds
beyond these derived facts is the occurrence of the unobstructed,
non-superseded, non-redirected round itself, the history alignment at
the restart height,
the common latest-justification action states through the inclusion, and
the finality-pair inclusion ordering.

Assumption~\ref{ass:clean-target-opportunity} now supplies an
ordinary-height target opportunity
\((h_\star,T_\star)\), including the honest target pairs on one valid
descendant.  Lemma
\ref{lem:fresh-ordinary-justify} creates
\(\mathsf{JC}(h_\star,T_\star)\).
Definition~\ref{def:fresh-ordinary-action} makes every honest signer save
\(T_\star\) and sign no timeout at \(h_\star\) before the JC is processed.
Before the JC is processed, Definition~\ref{def:finality-vote-rule} cannot
create a finality lock at that height.  After it is processed, the same
definition can record only \(T_\star\).  Thus every honest finality-lock entry
at \(h_\star\) is empty or \(T_\star\).  Once the JC and its post-state are common,
Definition~\ref{def:finality-vote-rule} makes them all sign finality pair
\((h_\star,T_\star)\).  The same assumption supplies an honest slot proposer
that includes those pairs before a newer justification can be included.
Lemma~\ref{lem:honest-proposer-finalizes} finalizes
\(T_\star\).
\end{proof}

\begin{proposition}[Opportunity budget before the first confirmation]
\label{prop:opportunity-budget}
Under Assumptions~\ref{ass:fixed-electorate}, \ref{ass:sg-fault-bound},
\ref{ass:goldfish-committees},
\ref{ass:recovery-network}, \ref{ass:recovery-goldfish}, and
\ref{ass:frontier-opportunities}, suppose no first recovery confirmation
has yet occurred.  Let \(A\) be the number of distinct finalized
heights \emph{above \(h_F^0\)} first processed at some honest
validator before the first recovery confirmation --- by a finalizing
transition completing inside the window, which is finality progress, or
by the released evidence of a pre-window transition, charged once per
distinct height.  A first-processed finalized height at or below
\(h_F^0\) is excluded because it obstructs nothing: under the fault
bound its root lies on the safe finalized chain below an
already-recorded root (Lemma~\ref{lem:finalized-chain}), so its
processing changes no greatest activated pair and deactivates no
additional grade --- a block conflicting with it conflicts with the
deeper recorded root already processed.  Activated
finalized heights are monotone, and when every pre-confirmation finality
advance keeps \(H-h'>D_{\rm debt}\) --- case~1 of
Lemma~\ref{lem:debt-dichotomy} --- every such height lies below
\(H-D_{\rm debt}\), so \(A\leq H-D_{\rm debt}-h_F^0=:A_{\max}\),
the finalized-height budget below \(H\) --- an a-priori bound
evaluable at the relay point.  Then at most \(7+2A\) \emph{charged
opportunities} are needed to reach the first recovery confirmation.
The charged items are: a possible \(H-1\) catch-up proposal and one
possible repeat of it; a height-\(H\) proposal and one possible
repeat; a height-\(H+1\) proposal; a possible retry for the unique
height-\((H-1)\) justification; one round for the single
store-maximum rise; and at most two obstructed rounds per processed
finalized height.  The confirmation occurs within \((7+2A)\rho+3\)
rounds of the relay point: the proof maps each charged item to at most
one \(\rho\)-round interval, disjointly, with at most three further
closing-action rounds.

In plain words, seven charged proposer opportunities plus two per
processed finalized height force the first confirmation.
\end{proposition}

\begin{proof}
The every-\(\rho\)-round proposer event in
Assumption~\ref{ass:frontier-opportunities} supplies an honest round
proposer in every \(\rho\)-round interval, and every honest-proposer
round before the first confirmation falls in one of the four cases of
Lemma~\ref{lem:honest-proposer-coverage}.

The seven charged opportunities are those enumerated in the
statement, the rise being the single one
Lemma~\ref{lem:pre-confirmation-cap} permits.
The two repeats account for pipelining: a proposal's height pairs are
released only at its closing action, the following round's action, so
the round in between may see an unobstructed honest proposal at the same
height before the pairs are includable; the advance completes at the
assumed includer, and the head-persistence clause of
Assumption~\ref{ass:frontier-opportunities} excludes any further
uncharged same-height opportunity.  The rise consumes its own
opportunity because a higher maximum only \emph{permits} a higher
proposal: the viable band of Definition~\ref{def:finality-root} keeps
both the maximum and the one-behind heights, ordinary Goldfish votes may
still favor the lower branch, and one further honest proposal at the
lower height may still be needed to release its pairs.

The case accounting over Lemma~\ref{lem:honest-proposer-coverage} is as
follows.  Uncharged case-1 rounds, such as sibling-branch proposals
while a release awaits its assumed includer, fall inside intervals the
wall-clock mapping below already counts: the includer clause of
Assumption~\ref{ass:frontier-opportunities}, not proposer fairness
alone, supplies each advance.  Charged case-1 rounds progress through at
most five height proposals --- the \(H-1\) catch-up and its possible
pipelining repeat, the height-\(H\) proposal and its possible repeat,
and the height-\(H+1\) proposal: if the selected proposal height is
\(H-1\), Lemma~\ref{lem:height-catchup} gives a selected branch at
height \(H\) or moves the permitted heights higher; at height \(H\),
Lemma~\ref{lem:leave-nonjustifiable-height} advances a selected
descendant to \(H+1\); and
Lemma~\ref{lem:honest-proposer-forces-first-confirmation} makes the
next unobstructed honest height-\(H+1\) proposal produce the first
confirmation.  Case-2 rounds number at most one, because
Lemma~\ref{lem:pre-confirmation-cap} permits one maximum rise, from
\(H\) to \(H+1\).  Case-3 rounds number at most one, by
Lemma~\ref{lem:target-uniqueness} and the cannot-recur argument of
Lemma~\ref{lem:honest-proposal-confirms}.  Case-4 rounds number at most
\(2A\), two per processed finalized height --- one if its evidence
deactivates a graded candidate root mid-round, and one if its activation
is staggered across the following round boundary.  Only
pre-confirmation processing can obstruct the charged opportunities,
which all precede the confirmation.  These \(5+1+1+2A\) opportunities
are covered by the finite-horizon event in
Assumption~\ref{ass:frontier-opportunities}.

\emph{Wall-clock mapping.}  The intervals consumed are disjoint and
number at most \(7+2A\).  The \(H-1\), \(H\), and \(H+1\)
proposals each arrive within one \(\rho\)-round interval of the
preceding milestone, by the recurrence.  Each of the two pair releases
below \(H+1\) is includable within one further \(\rho\)-round
interval, by the inclusion clause of
Assumption~\ref{ass:frontier-opportunities}; a closing action signs at
the following round's action time, and a later slot of its own round
may already include the released pairs, but the earliest block
guaranteed to be able to include them is a following first-slot
proposal, so the count uses only the guaranteed includer.  The possible
same-height repeat proposal falls inside that same inclusion-wait
interval --- charged in the count above, but consuming no separate
interval --- and so does every uncharged sibling round of the wait.
The retry, the rise, and each of the at most \(2A\) obstructed rounds
consume one interval each.  That is \(3+2+1+1+2A=7+2A\) disjoint
intervals.  The three closing actions of the height steps add at most
three further rounds, and the confirmation occurs at the
height-\((H{+}1)\) proposal's closing action.  Hence the bound
\((7+2A)\rho+3\).
\end{proof}

\begin{theorem}[Recovery reaches \(H+2\) and restores a persistent chain]
\label{thm:eventual-healing}
Under Assumptions~\ref{ass:fixed-electorate},
\ref{ass:sg-fault-bound}, \ref{ass:goldfish-committees},
\ref{ass:recovery-network}, \ref{ass:recovery-goldfish},
\ref{ass:frontier-opportunities}, and
\ref{ass:clean-target-opportunity}, with probability at least
\(1-\operatorname{negl}(\lambda)\) --- taken over the executions of
the cryptographic model, Assumption~\ref{ass:crypto} --- exactly one
of the following three outcomes holds inside every recovery window in
the sense of Definition~\ref{def:recovery-window} of length at least
\((11+2A_{\max})\rho+6\) rounds beyond the relay point ---
\(A_{\max}=H-D_{\rm debt}-h_F^0\), the a-priori processed-height
bound of Proposition~\ref{prop:opportunity-budget}, evaluable at the
relay point --- the end-to-end deadline accounted by
Corollary~\ref{cor:end-to-end} below:

\emph{Outcome (i): healing.}  Every finality advance in the window
other than the confirmed chain's own above-\(H\) finalizations keeps
\(H-h'>D_{\rm debt}\) --- case~1 of
Lemma~\ref{lem:debt-dichotomy}; the confirmed chain's own
finalizations above \(H\) are fresh in-window transitions, fall in
case~2's first subcase, and complete item~5 below rather than replacing
it.  Then all five of the following events occur:
\begin{enumerate}
 \item a selected viable branch catches up from \(H-1\), when necessary, and
   advances through \(H\) to \(H+1\);
 \item the first honest official confirmation \(B_\star\) at height at least
   \(H+1\) occurs.  This is the safety event whether its proposer is honest
   or faulty;
 \item after the action that grades the resulting SG batch, every honest raw
   Goldfish vote and nonempty SG head descends from \(B_\star\), and the
   store-global \(h_{\max}-1\) filter never removes its subtree;
 \item a valid descendant of \(B_\star\) advances from \(H+1\) to \(H+2\);
   this advance may be by target or by progress; and
 \item at the clean ordinary-height target opportunity supplied by
   Assumption~\ref{ass:clean-target-opportunity}, honest votes create a
   justification; after its finality action, the honest inclusion supplied
   by that assumption finalizes its target.
\end{enumerate}

\emph{Outcome (ii): direct restart.}  Some advance \emph{other than
the confirmed chain's own above-\(H\) finalizations} reaches
\(H-h'\leq D_{\rm debt}\), and the first such advance is a
transition completed inside the window: by
Lemma~\ref{lem:debt-dichotomy} it is a fresh finalization within
\(D_{\rm debt}\) of the frontier, and the window concludes with
finality restarted by that advance, in place of items~1--5.

\emph{Outcome (iii): disruption.}  The first such advance is the
activation of a pre-window transition from a released historical
branch: the window ends as a finite disruption, and
Corollary~\ref{cor:recurring-healing} below applies to a subsequent
window at the re-targeted guarded height.

Within outcome~(i): if no earlier first recovery confirmation occurs,
the first confirmation occurs within \((7+2A)\rho+3\) rounds of the
relay point, where \(A\) counts the distinct finalized heights
above \(h_F^0\) first processed at some honest validator before it;
Proposition~\ref{prop:opportunity-budget} defines the charged
opportunities behind this count and proves the bound.  After
\(B_\star\), the drain action signs the height-\(H+1\) pairs and
the next honest descendant proposal reaches \(H+2\).  Finality then
restarts within the separately assumed clean target opportunity;
ordinary honest proposer fairness alone does not establish that
opportunity.

In plain words, the first confirmation supplies safety, one later SG batch
closes stabilization, and the separate clean target opportunity then
restarts finality.
\end{theorem}

\begin{proof}
The three outcomes are exclusive and exhaustive: outcome (i) is the
complement of the event that some advance outside the excluded class
--- the confirmed chain's own above-\(H\) finalizations --- reaches
\(H-h'\leq D_{\rm debt}\), and outcomes (ii) and (iii) split that
event by the nature of its first advance, the two cases of
Lemma~\ref{lem:debt-dichotomy}.

The every-\(\rho\)-round proposer event in
Assumption~\ref{ass:frontier-opportunities} supplies an honest round proposer
in every \(\rho\)-round interval.  Lemma
\ref{lem:honest-proposer-coverage} proves that each such proposer is an
unobstructed opportunity, unless the permitted heights have already moved or the
unique height-\(H-1\) justification causes its one possible retry.

If the selected proposal height is \(H-1\),
Lemma~\ref{lem:height-catchup} gives a selected branch at height \(H\) or
moves the permitted heights higher.  A certificate on a sibling branch is used
only to determine that band; it is never treated as a transition on the
selected branch.  At height \(H\),
Lemma~\ref{lem:leave-nonjustifiable-height} advances a selected descendant to
\(H+1\), unless the first recovery confirmation has already occurred.  These
are the first two possible proposer opportunities.  Justification uniqueness
makes the old-root redirection consume at most one extra opportunity.

If Definition~\ref{def:first-recovery-confirmation} has occurred already,
use that event.  Otherwise
Lemma~\ref{lem:honest-proposer-forces-first-confirmation} makes the next
honest height-\(H+1\) proposal produce it.  This is the third ordinary
opportunity and proves item~2 without assuming that
the proposer of \(r_\star\) is honest.
Lemmas~\ref{lem:first-confirmation-heads} and
\ref{lem:first-confirmation-clears} classify the first SG batch into
compatible and finalized-excluded heads, then remove every live conflicting
grade at the drain action.
Lemma~\ref{lem:first-confirmation-persists} proves item~3 and also proves the
height-filter fact used here: a branch can raise the known maximum to
\(x+1\) only after an honest quorum signer has confirmed height \(x\), so
the \(B_\star\)-subtree is already at the permitted threshold \(x\).

Lemma~\ref{lem:recovery-height-justifies} gives the final
\(H+1\to H+2\) transition on a \(B_\star\)-descendant, proving item~4.
The drain action supplies the signatures; the first later honest descendant
proposal supplies their inclusion, as required by
Assumption~\ref{ass:frontier-opportunities}.
Lemma~\ref{lem:post-healing-finality} then applies
Assumption~\ref{ass:clean-target-opportunity} to justify and finalize its
aligned target.  This proves item~5.
Proposition~\ref{prop:opportunity-budget} supplies the \(7+2A\)
opportunity count and the \((7+2A)\rho+3\)-round bound; its charged
case-1 opportunities are exactly the proposer opportunities used above.
The later finality restart uses the separate clean target opportunity
stated in the theorem.
\end{proof}

\begin{corollary}[End-to-end restart deadline]
\label{cor:end-to-end}
In outcome~(i) of Theorem~\ref{thm:eventual-healing}, finality is
restarted within \((11+2A)\rho+6\) recovery rounds of the relay
point: \((7+2A)\rho+3\) rounds to the first confirmation
(Proposition~\ref{prop:opportunity-budget}); one round to the drain
action; at most \(\rho\) rounds to the honest descendant proposal
whose inclusion completes the \(H+1\to H+2\) advance --- the
inclusion clause of Assumption~\ref{ass:frontier-opportunities}; at
most \(\rho\) rounds to the ordinary-height target opportunity ---
Assumption~\ref{ass:clean-target-opportunity}; at most \(\rho\)
rounds to the including block that processes the height-\(h_\star\)
tallies and creates the justification --- the inclusion clause of
Assumption~\ref{ass:frontier-opportunities}; at most one round for
that justification to become common --- the relay rule of
Assumption~\ref{ass:recovery-network}; one action for the finality
pairs; and at most \(\rho\) rounds to the honest slot proposer that
includes them --- the inclusion deadline of
Assumption~\ref{ass:clean-target-opportunity}.  A window of this
length therefore contains every action and inclusion the theorem uses,
and since \(A\leq A_{\max}\) a priori, a window of
\((11+2A_{\max})\rho+6\) rounds suffices in every execution of
outcome~(i).

In plain words, the window-length clause is the explicit sum of the
assumed per-step deadlines.
\end{corollary}

\begin{proof}
Each summand is the deadline stated by the cited proposition or
assumption clause, and every event of outcome~(i) falls at or before
its summand's deadline; the total is
\((7+2A)\rho+3+1+\rho+\rho+\rho+1+1+\rho=(11+2A)\rho+6\).
\end{proof}

\begin{corollary}[Recovery can be repeated after later disruptions]
\label{cor:recurring-healing}
If every later finite disruption is followed by another recovery window
satisfying Theorem~\ref{thm:eventual-healing} at a later universally
nonjustifiable height, another target is finalized after finitely many
further windows: consecutive disruption outcomes are finitely many, and
every finalizing window finalizes above the previous finalized height.

In plain words, the same bounded recovery applies again after each finite
period of asynchrony.
\end{corollary}

\begin{proof}
Apply Theorem~\ref{thm:eventual-healing} in each recovery window.  A
window ending in outcome (iii) activates a released pre-window
finalized transition.  Honest validators sign finality only through the
protocol, and a fresh in-window finalization is outcome (ii), not a
disruption, so every disruption activates a transition already signed
before its window --- drawn from the finitely many signed objects that
exist at any point of the polynomial horizon of
Assumption~\ref{ass:crypto} --- and each activates a distinct one: an
activated transition is processed at every honest validator and cannot
terminate a later window again.  Consecutive outcome-(iii) windows are
therefore finitely many, so some later window ends in outcome (i) or
(ii) and finalizes a target.  That window's restart height exceeds the
previous finalized height, so finalized height strictly increases
across finalizing windows.
\end{proof}

\section{Incentives and empty votes}\label{sec:incentives}

This section states the one incentive requirement a deployment must meet ---
an empty height pair must not exempt its signer from a stall penalty --- and
records what the paper deliberately leaves open about the reward schedule.  It
proves no result, and no earlier or later result depends on it.

The safety and healing theorems do not depend on rewards or an inactivity
leak.  The relevant distinction is that the empty height pair
\((\bot,\bot)\) sets no progress bit and cannot exempt its signer from a stall
penalty.  Otherwise every validator could emit empty, avoid the penalty, and
prevent height advance forever.  A timeout pair \((h,\bot)\) is not empty and
does set the progress bit.

This paper proves no quantitative inactivity-leak bound.  A deployment may
pay target and finality participation once per round or once per height, but
must specify eligibility, rounding, fork-transition accounting, and the
balances and epoch used.  Paying once per height avoids repeatedly rewarding
one stalled checkpoint.  The liveness proof is independent of that choice.

\section{Implementation requirements}\label{sec:refinement}

This section collects what an implementation must provide for the preceding
proofs to apply to it: what checkpoint sync must authenticate, what blocks and
transport certificates may and may not do, how far the strong G0 gate reaches,
and which parameters are held fixed.  It proves no result, and no earlier or
later result depends on it.

Checkpoint sync authenticates the current height, start slot, target,
nonjustifiable flag, finalized state, latest justification, participation
arrays, and the evidence needed to replay them.
Lemma~\ref{lem:target-bit-compression} makes those two Boolean arrays
sufficient: no per-validator target root is part of the state, and the retained
signed attestations remain the only evidence needed for E1 and E2.  A checkpoint
at a transition block also authenticates that head for the same-slot target
rule.  Unknown-reference attestations are queued and replayed when their
blocks and ancestry arrive.

Blocks contain individual attestations.  A transport certificate is decoded
into those attestations and has no independent state-transition effect.
There is no recovery flag, recovery certificate branch, or
recovery-specific case in block validity or
\textsc{state\_transition}.  Height events are checked directly after every
block and in every empty slot; no block-supplied close flag selects them.
The empty-slot check is a provable no-op (Lemma~\ref{lem:empty-slot-noop})
and is kept for uniformity of the state machine.

Every graded root a validator adopts is backed by an honest gated head
(Lemma~\ref{lem:root-backing}, under
Assumption~\ref{ass:sg-fault-bound}), so an implementation may treat any
adopted graded root as accountable to at least one honest signer.
Strong G0 gates only the recovery SG head and current-height pair.  The
latest-\(J\) finality pair remains independent.  Goldfish proposal views
contain only ordinary blocks and raw Goldfish votes, and receiver validity
does not require a Byzantine proposal's parent to equal the receiver's
recomputed GHOST.

The paper fixes \(V,w\).  Validator churn or changing balances require a
separate overlap theorem and weak-subjectivity assumptions.  Weighted SG and
FG thresholds do not replace Goldfish's source-native unit identity counts.

\section{Conclusion}\label{sec:conclusion}

The finality state machine is the ordinary cumulative per-attestation
machine: exact current targets set both height bits, empty or other on-chain
current-height targets set progress, the latest justification alone has a
finality tally, and each direct height check tests finality, target, then progress.
Recovery adds no recovery-specific state-machine mode; the two permanent
state-machine deltas --- the broader progress predicate and the periodic
nonjustifiable-height flag --- are disclosed in
Definition~\ref{def:direct-height-position}.
Conflicting finalization at any height exposes at least \(2q-W\) provably E1- or
E2-slashable weight, without any synchrony assumption; this is
Theorem~\ref{thm:accountable-safety}, and
Lemma~\ref{lem:signer-safety} makes that weight adversarial.

The SG recovery mechanism uses four strict-majority grades over deadline views at
\(-\Delta,0,+\Delta,+2\Delta\), with grade~3 read on the intersection
of the \(-\Delta\) and \(+\Delta\) views.  It never merges SG heads into Goldfish.
The recovery slot is the \(4\Delta\) slot of Goldfish's optimistic
fast confirmations, with its dedicated confirmation phase before the
slot-view handling, as Definition~\ref{def:recovery-timing} states.  The
profile's departures from the Goldfish source are: the time-shifted
available confirmation replaces both source confirmation rules ---
\(\kappa\)-deep and optimistic fast --- splitting the support freeze
and the evaluation across two slots with a live-participation
denominator, and the source's buffer merge in the dedicated
fast-confirmation phase is dropped, the receipt buffer merging only at
the slot-view freeze; the fork-choice walk chooses within the
viability-filtered candidate
tree, counting weights over the full known tree, where the source's
walk uses an unrestricted child set --- starting the walk at a non-genesis
root is not itself a departure: the source's composed protocol likewise
starts its modified walk at the latest checkpoint~\cite{goldfish};
the discard-both wrapper of
Definition~\ref{def:recovery-proposal} governs proposer equivocation;
and committees are abstracted as stated with
Assumption~\ref{ass:goldfish-committees}.  The
first honest official confirmation at height \(H+1\) is the healing event,
regardless of who proposed its block.  Every honest SG head in that action is
either compatible with the confirmation or already excluded by the
confirmer's finalized root.  Relay makes the excluded heads inert, and the
next eligible batch removes every live conflicting grade.  If no earlier action confirms, an honest round
proposal supplies the exact Goldfish vote batch, and the time-shifted-quorum
available confirmation carries that batch to the action-time head that forces
the event; the coverage lemma proves that no separate unobstructed-round
assumption is needed beyond its listed cases.
The confirmation-gated height rule then keeps the healed subtree within the
one-height viability band.  Ordinary individual FG attestations advance the
chain through \(H+2\), by target or by progress.  A later clean target
opportunity supplies one common target quorum.  A later action signs its
finality pairs, and an honest block includes them and finalizes the target.  The
finality tally counts only on selected states and inclusion branches that processed that justification as their
latest unfinalized justification.  The nonjustifiable-height rule excludes a
height-\(H\) justification, while finalized-chain safety places every older
finalized root an ancestor of the first recovery confirmation.  This is
Theorem~\ref{thm:eventual-healing}.

\begin{thebibliography}{9}
\bibitem{simplex}
B.~Y. Chan and R.~Pass.
\newblock Simplex Consensus: A Simple and Fast Consensus Protocol.
\newblock IACR Cryptology ePrint Archive, Report 2023/463, 2023.

\bibitem{goldfish}
F.~D'Amato, J.~Neu, E.~N.~Tas, and D.~Tse.
\newblock Goldfish: No More Attacks on Ethereum?!
\newblock arXiv:2209.03255, 2022.

\bibitem{accountability}
J.~Neu, E.~N.~Tas, and D.~Tse.
\newblock The Availability-Accountability Dilemma and its Resolution via
Accountability Gadgets.
\newblock 2021.
\end{thebibliography}
